% Applications of vectors in mechanics — Pure Mathematics with Mechanics Upper Sixth — Cours signé OnBuch+
% Official MINESEC syllabus. Compiler avec : tectonic PMM08-applications-of-vectors-in-mechanics.tex
\newcommand{\DOCMATIERE}{Pure Mathematics with Mechanics}
\newcommand{\DOCNIVEAU}{Upper Sixth}
\newcommand{\DOCMODULE}{Upper Sixth — Pure mathematics with mechanics}
\newcommand{\DOCLECON}{8}
\newcommand{\DOCTITRE}{Applications of vectors in mechanics}
% OnBuch+ — préambule « cours signés » ENRICHI (compilé avec tectonic / XeLaTeX)
% Système visuel « Pop » d'OnBuch+ : fond crème (jamais blanc), bordures encre
% épaisses, ombres « dures » décalées, titres Archivo Black en capitales.
% Version MATHÉMATIQUES : graphes (pgfplots/tikz), boîtes pédagogiques
% (définition, propriété, méthode, attention, le savais-tu, exercice résolu,
% exercice, TP, bilan), page de garde et signature OnBuch+ ancrée en bas.
%
% Macros attendues dans main.tex AVANT \input{preamble} :
%   \DOCMATIERE  \DOCNIVEAU  \DOCMODULE  \DOCLECON (numéro)  \DOCTITRE
\documentclass[11pt]{article}

\usepackage{fontspec}
\usepackage[english]{babel}
\usepackage[a4paper,margin=2cm,top=3.4cm,bottom=2.8cm,headheight=1.4cm,footskip=1.3cm]{geometry}
\usepackage{amsmath,amssymb,amsfonts}
\usepackage{mathtools} % \xmapsto, \coloneqq... souvent utilisés par les modèles
\usepackage{cancel} % simplifications barrées (bilans de forces)
% \abs{z} (module, valeur absolue) et \norm{u} (norme), utilisés par les modèles
\providecommand{\abs}{}\let\abs\relax\DeclarePairedDelimiter{\abs}{\lvert}{\rvert}
\providecommand{\norm}{}\let\norm\relax\DeclarePairedDelimiter{\norm}{\lVert}{\rVert}
\usepackage[table]{xcolor} % [table] : fournit aussi \rowcolors (alternance de lignes)
\usepackage{pagecolor}
\usepackage{tikz}
\usetikzlibrary{arrows.meta,positioning,calc,shapes.geometric,shapes.misc,shapes.arrows,decorations.pathmorphing,decorations.pathreplacing,decorations.markings,patterns,shadows,mindmap,trees,fit,backgrounds,matrix,intersections,angles,quotes}
\usepackage{pgfplots}
\usepackage[version=4]{mhchem} % équations nucléaires \ce{^{235}_{92}U}
\usepackage[european,siunitx]{circuitikz} % schémas électriques
\pgfplotsset{compat=1.18}
\usepgfplotslibrary{fillbetween,groupplots} % aires entre courbes (calcul intégral)
\usepackage{enumitem}
\usepackage{fancyhdr}
% [most] charge toutes les bibliothèques tcolorbox (listings, minted, raster,
% documentation...) et gonfle inutilement la mémoire TeX sur un document long
% (~300 boîtes) : on ne charge que ce qui est réellement utilisé.
\usepackage[skins,breakable]{tcolorbox}
\usepackage{graphicx}
\usepackage{colortbl}
\usepackage{booktabs}
\usepackage{tabularx}
\usepackage{diagbox} % case d'en-tête diagonale des tableaux à double entrée
% Colonnes extensibles centrée / gauche / droite (C, L, R), souvent utilisées par les modèles
\newcolumntype{C}{>{\centering\arraybackslash}X}
\newcolumntype{L}{>{\raggedright\arraybackslash}X}
\newcolumntype{R}{>{\raggedleft\arraybackslash}X}
\usepackage{array}
\usepackage{multicol}
\usepackage{siunitx}
% parse-numbers=false : accepte aussi \SI{2,5 \times 10^{-3}}{...}
\sisetup{locale=UK,output-decimal-marker={.},group-minimum-digits=4,parse-numbers=false}
\DeclareSIUnit\ohm{\ensuremath{\symup{\Omega}}} % U+2126 absent de la police
\DeclareSIUnit\division{div} % réglages d'oscilloscope (ms/div, V/div)
\DeclareSIUnit\tr{tr} % tours (vitesse de rotation, ex. \SI{1500}{\tr\per\minute})
\usepackage{qrcode}
% \degree : commande fréquemment utilisée par les modèles pour un angle ou une
% température en dehors de \SI{}{} (non fournie par les paquets chargés).
\providecommand{\degree}{^{\circ}}
% \qed (fin de démonstration), utilisé par les modèles sans amsthm
\providecommand{\qed}{\hfill$\square$}
% \barre : double barre des valeurs interdites dans un tableau de variations (tkz-tab)
\providecommand{\barre}{\ensuremath{\|}}
% environnement proof (amsthm non chargé) utilisé par les modèles
\ifdefined\proof\else\newenvironment{proof}[1][Proof]{\par\noindent\textit{#1.}\ }{\qed\par}\fi
\usepackage{lastpage}
\usepackage{hyperref}
\hypersetup{hidelinks}

% ============================================================
% Palette « Pop » OnBuch+ — mêmes hex que la classe P (plus_ui.dart)
% ============================================================
\definecolor{poporange}{HTML}{F59321}
\definecolor{popdark}{HTML}{E07A0C}
\definecolor{popcream}{HTML}{FFF6E8}
\definecolor{popink}{HTML}{1C1714}
\definecolor{poppurple}{HTML}{7A5AE0}
\definecolor{popgreen}{HTML}{2E9E6A}
\definecolor{poppink}{HTML}{E0457B}
\definecolor{popblue}{HTML}{2F7DE1}
\definecolor{popgold}{HTML}{C98A12}
\definecolor{popmuted}{HTML}{8A7F76}
\definecolor{popline}{HTML}{EADFD2}
\definecolor{popgray}{HTML}{8A7F76}
\colorlet{popgrey}{popgray}
% Alias tolérants (le modèle nomme parfois les couleurs différemment)
\colorlet{popwhite}{white}
\colorlet{popblack}{popink}
\colorlet{popred}{poppink}
\colorlet{popyellow}{popgold}
% Teintes claires pour fonds de tableaux / zones
\colorlet{popblueL}{popblue!10!white}
\colorlet{poporangeL}{poporange!14!white}
\colorlet{popgreenL}{popgreen!12!white}
\colorlet{poppurpleL}{poppurple!12!white}
\colorlet{poppinkL}{poppink!12!white}
\colorlet{popgoldL}{popgold!14!white}

\pagecolor{popcream}

% ============================================================
% Typographie
% ============================================================
\defaultfontfeatures{Ligatures=TeX}
\setmainfont{PlusJakartaSans-Medium.ttf}[Path=../../../fonts/,BoldFont=PlusJakartaSans-Bold.ttf]
% Maths en sans-serif (Fira Math) avec chiffres et lettres droites de Plus
% Jakarta, pour que les formules se fondent dans le texte.
\usepackage[math-style=ISO,bold-style=ISO]{unicode-math}
\setmathfont{FiraMath-Regular.otf}[Path=../../../fonts/]
\setmathfont{PlusJakartaSans-Medium.ttf}[Path=../../../fonts/,range={up/num,up/{Latin,latin},\mathup}]
% (icomma retiré : décimales avec point en anglais)
% Caractères mathématiques Unicode absents des polices de texte (Plus Jakarta,
% Archivo Black) : rendus par leur équivalent mathématique, en texte comme en
% formule — sinon ils s'affichent en carré vide (ex. « Arithmétique dans ℤ »).
\usepackage{newunicodechar}
\newunicodechar{ℝ}{\ensuremath{\mathbb{R}}}
\newunicodechar{ℤ}{\ensuremath{\mathbb{Z}}}
\newunicodechar{ℂ}{\ensuremath{\mathbb{C}}}
\newunicodechar{ℕ}{\ensuremath{\mathbb{N}}}
\newunicodechar{ℚ}{\ensuremath{\mathbb{Q}}}
\newunicodechar{𝔻}{\ensuremath{\mathbb{D}}}
\newunicodechar{α}{\ensuremath{\alpha}}
\newunicodechar{β}{\ensuremath{\beta}}
\newunicodechar{γ}{\ensuremath{\gamma}}
\newunicodechar{δ}{\ensuremath{\delta}}
\newunicodechar{ε}{\ensuremath{\varepsilon}}
\newunicodechar{θ}{\ensuremath{\theta}}
\newunicodechar{λ}{\ensuremath{\lambda}}
\newunicodechar{μ}{\ensuremath{\mu}}
\newunicodechar{σ}{\ensuremath{\sigma}}
\newunicodechar{τ}{\ensuremath{\tau}}
\newunicodechar{φ}{\ensuremath{\varphi}}
\newunicodechar{ψ}{\ensuremath{\psi}}
\newunicodechar{ω}{\ensuremath{\omega}}
\newunicodechar{Σ}{\ensuremath{\Sigma}}
% majuscules produites par \MakeUppercase dans les titres (α-aminés -> Α)
\newunicodechar{Α}{\ensuremath{\alpha}}
\newunicodechar{Β}{\ensuremath{\beta}}
\newunicodechar{Γ}{\ensuremath{\Gamma}}
\newunicodechar{Δ}{\ensuremath{\Delta}}
\newunicodechar{Ε}{\ensuremath{\varepsilon}}
\newunicodechar{Θ}{\ensuremath{\Theta}}
\newunicodechar{Λ}{\ensuremath{\Lambda}}
\newunicodechar{Μ}{\ensuremath{\mu}}
\newunicodechar{Τ}{\ensuremath{\tau}}
\newunicodechar{Φ}{\ensuremath{\Phi}}
\newunicodechar{Ψ}{\ensuremath{\Psi}}
\newunicodechar{Ω}{\ensuremath{\Omega}}
\newunicodechar{↦}{\ensuremath{\mapsto}}
\newunicodechar{∈}{\ensuremath{\in}}
\newunicodechar{∉}{\ensuremath{\notin}}
\newunicodechar{∧}{\ensuremath{\wedge}}
\newunicodechar{⇔}{\ensuremath{\Leftrightarrow}}
\newunicodechar{⟺}{\ensuremath{\Longleftrightarrow}}
\newunicodechar{⇒}{\ensuremath{\Rightarrow}}
\newunicodechar{⟹}{\ensuremath{\Longrightarrow}}
\newunicodechar{∘}{\ensuremath{\circ}}
\newunicodechar{∪}{\ensuremath{\cup}}
\newunicodechar{∩}{\ensuremath{\cap}}
\newunicodechar{≡}{\ensuremath{\equiv}}
\newunicodechar{⊂}{\ensuremath{\subset}}
\newunicodechar{⊥}{\ensuremath{\perp}}
\newunicodechar{∃}{\ensuremath{\exists}}
\newunicodechar{∀}{\ensuremath{\forall}}
\newunicodechar{∛}{\ensuremath{\sqrt[3]{\,}}}
\newunicodechar{∜}{\ensuremath{\sqrt[4]{\,}}}
\newunicodechar{⃗}{\ensuremath{{}^{\to}}}
\newunicodechar{ⁿ}{\ifmmode^{n}\else\textsuperscript{\ensuremath{n}}\fi}
\newunicodechar{ˣ}{\ifmmode^{x}\else\textsuperscript{\ensuremath{x}}\fi}
\newunicodechar{ᵇ}{\ifmmode^{b}\else\textsuperscript{\ensuremath{b}}\fi}
\newunicodechar{ʳ}{\ifmmode^{r}\else\textsuperscript{\ensuremath{r}}\fi}
\newunicodechar{ᵘ}{\ifmmode^{u}\else\textsuperscript{\ensuremath{u}}\fi}
\newunicodechar{ʸ}{\ifmmode^{y}\else\textsuperscript{\ensuremath{y}}\fi}
\newunicodechar{ᵃ}{\ifmmode^{a}\else\textsuperscript{\ensuremath{a}}\fi}
\newunicodechar{ᵉ}{\ifmmode^{e}\else\textsuperscript{\ensuremath{e}}\fi}
\newunicodechar{ᵏ}{\ifmmode^{k}\else\textsuperscript{\ensuremath{k}}\fi}
\newunicodechar{ᵖ}{\ifmmode^{p}\else\textsuperscript{\ensuremath{p}}\fi}
\newunicodechar{ᴺ}{\ifmmode^{N}\else\textsuperscript{\ensuremath{N}}\fi}
\newunicodechar{ᶜ}{\ifmmode^{c}\else\textsuperscript{\ensuremath{c}}\fi}
\newunicodechar{ʰ}{\ifmmode^{h}\else\textsuperscript{\ensuremath{h}}\fi}
\newunicodechar{ᵠ}{\ifmmode^{q}\else\textsuperscript{\ensuremath{q}}\fi}
\newunicodechar{ₙ}{\ifmmode_{n}\else\textsubscript{\ensuremath{n}}\fi}
\newunicodechar{ₐ}{\ifmmode_{a}\else\textsubscript{\ensuremath{a}}\fi}
\newunicodechar{ᵢ}{\ifmmode_{i}\else\textsubscript{\ensuremath{i}}\fi}
\newunicodechar{ᵣ}{\ifmmode_{r}\else\textsubscript{\ensuremath{r}}\fi}
\newunicodechar{ₖ}{\ifmmode_{k}\else\textsubscript{\ensuremath{k}}\fi}
\newunicodechar{ᵦ}{\ifmmode_{\beta}\else\textsubscript{\ensuremath{\beta}}\fi}
\newunicodechar{ₓ}{\ifmmode_{x}\else\textsubscript{\ensuremath{x}}\fi}
\newfontfamily\popdisplay{ArchivoBlack-Regular.ttf}[Path=../../../fonts/]
\newfontfamily\poplogo{SpaceGrotesk-Bold.ttf}[Path=../../../fonts/]
\newfontfamily\popsemi{PlusJakartaSans-SemiBold.ttf}[Path=../../../fonts/]
\newfontfamily\popxbold{PlusJakartaSans-ExtraBold.ttf}[Path=../../../fonts/]
\newfontfamily\popxboldls{PlusJakartaSans-ExtraBold.ttf}[Path=../../../fonts/,LetterSpace=3.0]

\setlist[enumerate]{leftmargin=1.7em,itemsep=5pt,topsep=4pt}
\setlist[itemize]{leftmargin=1.7em,itemsep=4pt,topsep=4pt,label={\color{poporange}\textbullet}}
\setlength{\parindent}{0pt}
\setlength{\parskip}{5pt}
\renewcommand{\arraystretch}{1.25}

% pgfplots : style Pop par défaut
\pgfplotsset{
  every axis/.append style={axis line style={popink,line width=1pt},tick style={popink},
    grid style={popline,line width=0.5pt},label style={font=\small},tick label style={font=\footnotesize}},
  popcurve/.style={poporange,line width=1.8pt,no markers,smooth},
}
% Même style disponible dans un \draw TikZ simple (hors pgfplots)
\tikzset{popcurve/.style={poporange,line width=1.8pt}}
\tikzset{popgrid/.style={popline,very thin}}
\colorlet{popcurve}{poporange} % le nom du style est parfois utilisé comme couleur

% ============================================================
% Pill / squiggle
% ============================================================
\newcommand{\poppill}[3]{%
  \tikz[baseline=-0.55ex]\node[draw=popink,line width=1.2pt,rounded corners=2.6pt,%
    fill=#2,inner xsep=6.5pt,inner ysep=3pt]{%
    \popxboldls\fontsize{7}{7}\selectfont\color{#3}\MakeUppercase{#1}};%
}
\newcommand{\popsquiggle}[1]{%
  \tikz[baseline=-0.4ex]{
    \draw[#1,line width=2.6pt,line cap=round]
      plot [smooth] coordinates {(0,0.09) (0.34,0.01) (0.68,0.21) (1.02,0.01) (1.36,0.10)};
  }%
}
% Mot-clé surligné (marqueur orange sous le texte)
\newcommand{\cle}[1]{{\popxbold\color{popdark}#1}}

% ============================================================
% En-tête / pied
% ============================================================
\pagestyle{fancy}
\fancyhf{}
\renewcommand{\headrulewidth}{0pt}
\renewcommand{\footrulewidth}{0pt}
\lhead{\raisebox{-0.15ex}{\poplogo\fontsize{13}{13}\selectfont\color{popink}OnBuch\color{poporange}+}}
\rhead{\poppill{\DOCMATIERE\ \textbullet\ \DOCNIVEAU\ \textbullet\ Chapter \DOCLECON}{white}{popink}}
\lfoot{\popxbold\fontsize{7.6}{7.6}\selectfont\color{popmuted}\MakeUppercase{Signed course OnBuch+}\ \textbullet\ {\popsemi\DOCTITRE}}
\rfoot{\tikz[baseline=-0.6ex]\node[rounded corners=8.5pt,fill=popink,minimum height=17pt,minimum width=17pt,inner xsep=5pt,inner sep=0]{\popxbold\fontsize{8}{8}\selectfont\color{popcream}\thepage\,/\,\pageref*{LastPage}};}

% ============================================================
% Titres
% ============================================================
\newcommand{\chaptitre}[2]{%
  \begin{tcolorbox}[enhanced,colback=poporange,colframe=popink,boxrule=2.6pt,arc=11pt,
      drop shadow=popink,left=15pt,right=15pt,top=12pt,bottom=11pt,before skip=3pt,after skip=18pt]
    \poppill{#2}{popink}{popcream}\par\vspace{9pt}
    {\raggedright\hyphenpenalty=10000\exhyphenpenalty=10000\popdisplay\fontsize{22}{24}\selectfont\color{popcream}\MakeUppercase{#1}\par}
  \end{tcolorbox}
}
% Section numérotée : pastille orange avec le numéro + titre Archivo
\newcounter{popsec}
\newcommand{\coursec}[1]{%
  \stepcounter{popsec}\setcounter{popsub}{0}%
  \par\vspace{16pt}\noindent
  \begin{minipage}{\linewidth}
  \tikz[baseline=-0.7ex]\node[draw=popink,line width=1.6pt,fill=poporange,rounded corners=4pt,minimum size=22pt,inner sep=0]{\popdisplay\fontsize{12}{12}\selectfont\color{popcream}\thepopsec};%
  \hspace{8pt}{\popdisplay\fontsize{15}{17}\selectfont\color{popink}\MakeUppercase{#1}}\par
  \vspace{4pt}\noindent{\color{poporange}\rule{36pt}{3.6pt}}
  \end{minipage}\par\nopagebreak\vspace{8pt}
}
\newcounter{popsub}
\newcommand{\courssub}[1]{%
  \stepcounter{popsub}%
  \par\vspace{9pt}\noindent{\popxbold\fontsize{11.5}{13}\selectfont\color{popink}{\color{poporange}\thepopsec.\thepopsub}\hspace{6pt}#1}\par\nopagebreak\vspace{4pt}
}

% ============================================================
% Boîtes pédagogiques — même recette : fond blanc, bordure épaisse colorée,
% ombre dure encre, badge plein. Titre optionnel [#1].
% ============================================================
\tcbset{popbase/.style={breakable,enhanced,colback=white,boxrule=2.3pt,arc=9pt,
  shadow={3pt}{-3pt}{0pt}{popink},left=14pt,right=14pt,top=11pt,bottom=11pt,before skip=11pt,after skip=15pt,
  fontupper=\popsemi\fontsize{10.6}{14.5}\selectfont\color{popink},
  fontlower=\popsemi\fontsize{10.6}{14.5}\selectfont\color{popink}}}
% Coche dessinée (le glyphe ✓ n'existe pas dans les polices du document)
\newcommand{\popcheck}[1]{\tikz[baseline=-0.5ex]\draw[#1,line width=1.6pt,line cap=round,line join=round](-0.13,0.0)--(-0.04,-0.09)--(0.13,0.1);}
\providecommand{\coche}{\popcheck{popgreen}}
\providecommand{\resultat}[1]{\par\smallskip\noindent\fcolorbox{popgreen}{popgreenL}{\parbox{\dimexpr\linewidth-2\fboxsep-2\fboxrule\relax}{\textbf{Result.} #1}}\par\smallskip}
\newcommand{\popboxhead}[3]{\poppill{#1}{#2}{white}\ifx\relax#3\relax\else\hspace{6pt}{\popxbold\fontsize{10.6}{12}\selectfont\color{popink}#3}\fi\par\vspace{7pt}}

\newtcolorbox{aretenir}[1][]{popbase,colframe=popgreen,before upper={\popboxhead{Key points}{popgreen}{#1}}}
\newtcolorbox{exemplebox}[1][]{popbase,colframe=poporange,before upper={\popboxhead{Exemple}{poporange}{#1}}}
\newtcolorbox{definition}[1][]{popbase,colframe=popblue,before upper={\popboxhead{Definition}{popblue}{#1}}}
\newtcolorbox{propriete}[1][]{popbase,colframe=poppurple,before upper={\popboxhead{Property}{poppurple}{#1}}}
\newtcolorbox{methode}[1][]{popbase,colframe=popgold,before upper={\popboxhead{Method}{popgold}{#1}}}
\newtcolorbox{attention}[1][]{popbase,colframe=poppink,before upper={\popboxhead{Warning}{poppink}{#1}}}
\newtcolorbox{experience}[1][]{popbase,colframe=popblue,colback=popblueL,before upper={\popboxhead{Experiment}{popblue}{#1}}}
% « Le savais-tu ? » : carte violette pleine (contexte, culture, Cameroun)
\newtcolorbox{savaistu}[1][]{popbase,colback=poppurple,colframe=popink,
  fontupper=\popsemi\fontsize{10.4}{14.2}\selectfont\color{white},
  before upper={\poppill{Did you know?}{white}{poppurple}\ifx\relax#1\relax\else\hspace{6pt}{\popxbold\color{white}#1}\fi\par\vspace{7pt}}}
% Exercice résolu : énoncé en haut, \tcblower puis la solution sur fond crème
\newcounter{popexo}\newcounter{popexr}
\newtcolorbox{exoresolu}[1][]{popbase,colframe=poporange,colbacklower=popcream,
  segmentation style={popink,line width=1.2pt,dashed},
  before upper={\stepcounter{popexr}\popboxhead{Worked example \thepopexr}{poporange}{#1}},
  before lower={\poppill{Solution}{popink}{popcream}\par\vspace{6pt}}}
% Exercice (non résolu en place) — la correction est dans la partie Corrigés
\newtcolorbox{exercice}[1][]{popbase,colframe=popink,
  before upper={\stepcounter{popexo}\popboxhead{Exercise \thepopexo}{popink}{#1}}}
% Corrigé
\newtcolorbox{corrige}[1][]{popbase,colframe=popgreen,colback=popgreenL,
  before upper={\popboxhead{Solution}{popgreen}{#1}}}
% Objectifs / prérequis / bilan
\newtcolorbox{objectifs}{popbase,colframe=popink,colback=poporangeL,
  before upper={\popboxhead{Chapter objectives}{popink}{}}}
\newtcolorbox{prerequis}{popbase,colframe=popink,colback=popblueL,
  before upper={\popboxhead{Prerequisites}{popblue}{}}}
\newtcolorbox{bilan}{popbase,colframe=popink,colback=white,boxrule=2.8pt,
  before upper={\popboxhead{Fiche bilan}{poporange}{L'essentiel à connaître pour l'examen}}}
% Figure encadrée avec légende
\newtcolorbox{popfigure}[1][]{enhanced,breakable=false,colback=white,colframe=popline,boxrule=1.4pt,arc=8pt,
  left=10pt,right=10pt,top=10pt,bottom=8pt,before skip=10pt,after skip=12pt,halign=center,
  lower separated=false,halign lower=center,
  fontlower=\popsemi\fontsize{9}{11.5}\selectfont\color{popmuted},
  #1}
\newcommand{\legende}[1]{\par\vspace{6pt}{\popsemi\fontsize{9}{11.5}\selectfont\color{popmuted}#1}}

% ============================================================
% Signature OnBuch+
% ============================================================
\newcommand{\poplogoinline}[1]{{\poplogo\fontsize{#1}{#1}\selectfont\color{popink}OnBuch\color{poporange}+}}
\newcommand{\signatureonbuch}{%
  \begin{tcolorbox}[enhanced,colback=popink,colframe=popink,boxrule=2.2pt,arc=10pt,
      drop shadow=poporange,left=14pt,right=14pt,top=10pt,bottom=10pt,before skip=0pt,after skip=0pt]
    \tikz[baseline=-0.7ex]\node[circle,fill=popgreen,draw=popcream,line width=1.3pt,minimum size=22pt,inner sep=0]{\popcheck{white}};%
    \hspace{9pt}\begin{minipage}[c]{0.62\linewidth}
      {\popxbold\fontsize{12}{13}\selectfont\color{popcream}Signed\ \textbullet\ {\poplogo OnBuch\color{poporange}+}}\par\vspace{2pt}
      {\raggedright\popsemi\fontsize{8}{10.5}\selectfont\color{popline}\MakeUppercase{OnBuch+ teaching team}\\\MakeUppercase{Follows the official MINESEC syllabus}\par}
    \end{minipage}\hfill
    {\poplogo\fontsize{22}{22}\selectfont\color{popcream}OnBuch\color{poporange}+}
  \end{tcolorbox}%
}

% ============================================================
% Page de garde
% #1 titre · #2 matière · #3 niveau · #4 module · #5 n° leçon · #6 durée ·
% #7 description / objectifs (texte ou itemize)
% ============================================================
\newcommand{\pagegarde}[7]{%
  \thispagestyle{empty}
  \newgeometry{a4paper,margin=1.7cm,top=1.6cm,bottom=1.4cm}%
  \begin{tikzpicture}[remember picture,overlay]
    \fill[poporange!18] ([xshift=-3.2cm,yshift=-3.6cm]current page.north east) circle (4.6cm);
    \fill[poppurple!14] ([xshift=2.4cm,yshift=7.6cm]current page.south west) circle (3.8cm);
  \end{tikzpicture}%
  {\poplogo\fontsize{30}{30}\selectfont\color{popink}OnBuch\color{poporange}+}\hfill
  \raisebox{0.6ex}{\poppill{Signed course \textbullet\ Premium}{popink}{popcream}}\par
  \vspace{1pt}\popsquiggle{poppurple}\par
  \vspace{8pt}
  \begin{tcolorbox}[enhanced,colback=poporange,colframe=popink,boxrule=2.8pt,arc=13pt,
      drop shadow=popink,left=18pt,right=18pt,top=12pt,bottom=13pt,after skip=10pt]
    \poppill{#2 \textbullet\ #3}{popink}{popcream}\hspace{5pt}\poppill{#4}{white}{popink}\par\vspace{8pt}
    {\popdisplay\fontsize{14}{14}\selectfont\color{popink}CHAPTER #5}\par\vspace{4pt}
    {\raggedright\hyphenpenalty=10000\exhyphenpenalty=10000\popdisplay\fontsize{25}{27}\selectfont\color{popcream}\MakeUppercase{#1}\par}
    \vspace{8pt}
    {\popxbold\fontsize{9.5}{9.5}\selectfont\color{popink}\MakeUppercase{Suggested time: #6}}
  \end{tcolorbox}
  \begin{tcolorbox}[enhanced,colback=white,colframe=popink,boxrule=2.2pt,arc=10pt,
      drop shadow=popink,left=16pt,right=16pt,top=10pt,bottom=10pt,after skip=8pt]
    \poppill{In this chapter}{poppurple}{white}\par\vspace{5pt}
    \popsemi\fontsize{9.3}{12.6}\selectfont\color{popink}#7
  \end{tcolorbox}
  \noindent\begin{minipage}[t]{0.487\linewidth}
    \begin{tcolorbox}[enhanced,colback=popgreen,colframe=popink,boxrule=2.2pt,arc=10pt,
        drop shadow=popink,left=11pt,right=11pt,top=10pt,bottom=10pt,height=3.1cm,valign=center]
      \tcbox[colback=white,colframe=popink,boxrule=1.3pt,arc=5pt,boxsep=0pt,nobeforeafter,
        left=3pt,right=3pt,top=3pt,bottom=3pt,valign=center]{\qrcode[height=1.9cm]{https://play.google.com/store/apps/details?id=com.luvvix.onbuch&hl=en}}%
      \hspace{8pt}\begin{minipage}[c]{0.5\linewidth}
        {\popdisplay\fontsize{11}{12.5}\selectfont\color{white}\MakeUppercase{Download OnBuch}}\par\vspace{4pt}
        {\raggedright\popsemi\fontsize{8.4}{11}\selectfont\color{white}Free \textbullet\ Google Play\\AI tutor, past papers, results\par}
      \end{minipage}
    \end{tcolorbox}
  \end{minipage}\hfill
  \begin{minipage}[t]{0.487\linewidth}
    \begin{tcolorbox}[enhanced,colback=poppurple,colframe=popink,boxrule=2.2pt,arc=10pt,
        drop shadow=popink,left=11pt,right=11pt,top=10pt,bottom=10pt,height=3.1cm,valign=center]
      \tikz[baseline=-0.7ex]\node[circle,fill=white,draw=popink,line width=1.3pt,minimum size=1.6cm,inner sep=0]{\popdisplay\fontsize{24}{24}\selectfont\color{poppurple}+};%
      \hspace{8pt}\begin{minipage}[c]{0.6\linewidth}
        {\popdisplay\fontsize{11}{12.5}\selectfont\color{white}\MakeUppercase{Go OnBuch+}}\par\vspace{4pt}
        {\raggedright\popsemi\fontsize{8.4}{11}\selectfont\color{white}All signed courses, unlimited AI tutor, PDF sheets — OnBuch+ tab in the app\par}
      \end{minipage}
    \end{tcolorbox}
  \end{minipage}
  \vspace{8pt}
  \popcontenu
  \vfill
  \signatureonbuch
  \restoregeometry
  \newpage
}

% Rangée « Ce cours contient » (page de garde)
\newcommand{\poptile}[3]{% #1 couleur #2 gros texte #3 légende
  \begin{tcolorbox}[enhanced,colback=#1,colframe=popink,boxrule=2pt,arc=9pt,shadow={2.5pt}{-2.5pt}{0pt}{popink},
      left=6pt,right=6pt,top=9pt,bottom=9pt,halign=center,valign=center,height=1.75cm,width=\linewidth,nobeforeafter]
    {\popdisplay\fontsize{12.5}{13}\selectfont\color{white}\MakeUppercase{#2}}\par\vspace{3pt}
    {\centering\popsemi\fontsize{7.8}{9.5}\selectfont\color{white}#3\par}
  \end{tcolorbox}}
\newcommand{\popcontenu}{%
  \par\noindent{\popxboldls\fontsize{7.5}{7.5}\selectfont\color{popmuted}THIS SIGNED COURSE CONTAINS}\par\vspace{7pt}
  \noindent\begin{minipage}[t]{0.235\linewidth}\poptile{popblue}{Course}{complete, illustrated, step by step}\end{minipage}\hfill
  \begin{minipage}[t]{0.235\linewidth}\poptile{poporange}{Methods}{and worked examples}\end{minipage}\hfill
  \begin{minipage}[t]{0.235\linewidth}\poptile{poppink}{Exercises}{exam style + solutions}\end{minipage}\hfill
  \begin{minipage}[t]{0.235\linewidth}\poptile{popgreen}{Summary sheet}{the essentials to remember}\end{minipage}\par}

% Sommaire
\newtcolorbox{sommaire}{enhanced,colback=white,colframe=popink,boxrule=2.2pt,arc=10pt,
  drop shadow=popink,left=18pt,right=18pt,top=13pt,bottom=13pt,before skip=6pt,after skip=20pt,
  before upper={\poppill{Contents}{poppurple}{white}\par\vspace{9pt}\popsemi\fontsize{10.6}{14.5}\selectfont\color{popink}}}

% Signature de fin de cours — poussée tout en bas de la dernière page
\newcommand{\findecours}{%
  \par\vfill\nopagebreak
  \begin{center}\popsquiggle{poporange}\end{center}\vspace{4pt}
  \signatureonbuch
}
\AtBeginDocument{\def\llbracket{\mathopen{[\mkern-3.5mu[}}\def\rrbracket{\mathclose{]\mkern-3.5mu]}}} % intervalles d'entiers (glyphe ⟦ absent de la police)
% Glyphes absents de Fira Math : redéfinis à partir de glyphes disponibles
\AtBeginDocument{%
  \def\setminus{\mathbin{\backslash}}\let\smallsetminus\setminus
  \def\vdots{\mathord{\vbox{\baselineskip3.2pt\lineskiplimit0pt\kern4pt\hbox{.}\hbox{.}\hbox{.}}}}%
  \def\ddots{\mathinner{\mkern1mu\raise6.5pt\hbox{.}\mkern2mu\raise3.6pt\hbox{.}\mkern2mu\raise0.7pt\hbox{.}\mkern1mu}}%
  \def\triangle{\mathord{\scalebox{0.9}{$\symup{\Delta}$}}}%
  \def\nabla{\mathord{\raisebox{\depth}{\scalebox{1}[-1]{$\symup{\Delta}$}}}}%
  \def\checkmark{\popcheck{popdark}}%
  \def\star{\mathbin{\ast}}\def\bigstar{\mathord{\ast}}%
  \def\diamond{\mathbin{\scalebox{0.6}{$\Diamond$}}}%
  \def\bigcup{\mathop{\mathchoice{\raisebox{-0.3em}{\scalebox{1.7}{$\cup$}}}{\scalebox{1.25}{$\cup$}}{\cup}{\cup}}\displaylimits}%
  \def\bigcap{\mathop{\mathchoice{\raisebox{-0.3em}{\scalebox{1.7}{$\cap$}}}{\scalebox{1.25}{$\cap$}}{\cap}{\cap}}\displaylimits}%
  \def\lesseqqgtr{\mathrel{\lessgtr}}\def\bigoplus{\mathop{\oplus}\displaylimits}%
}
\providecommand{\ssi}{if and only if} % abréviation fréquente
\AtBeginDocument{\providecommand{\wideparen}{\overparen}} % arc de cercle

%% FIGFIX-BEGIN v4 — correctifs globaux des figures (docs/onbuch-generation/tools/figfix)
\usepackage{adjustbox}\usepackage{etoolbox}
% toute figure TikZ/pgfplots plus large que la ligne est réduite à la largeur disponible
\BeforeBeginEnvironment{tikzpicture}{\begin{adjustbox}{max width=\linewidth}}
\AfterEndEnvironment{tikzpicture}{\end{adjustbox}}
% légendes pgfplots lisibles par-dessus les courbes
\pgfplotsset{every axis/.append style={legend style={fill=white,fill opacity=0.92,text opacity=1,draw=black!15,inner sep=3pt}}}
% étiquettes d'arêtes (syntaxe edge["..."]) avec fond blanc
\tikzset{every edge quotes/.append style={fill=white,inner sep=1.2pt}}
%% FIGFIX-END
\begin{document}

\pagegarde{Applications of vectors in mechanics}{Pure Mathematics with Mechanics}{Upper Sixth}{Upper Sixth — Pure mathematics with mechanics}{8}{11 periods}{This chapter extends differentiation and integration from scalars to vectors and applies them to the kinematics of a particle: displacement, velocity and acceleration as vector quantities. It provides the essential mathematical machinery for the whole of GCE Advanced Level mechanics.
\vspace{4pt}
\begin{multicols}{2}\small
\begin{itemize}
\item By the end of this chapter I can differentiate a vector given in component form with respect to a scalar parameter (usually time).
\item By the end of this chapter I can interpret the derivative of a position vector physically as velocity, and the derivative of velocity as acceleration.
\item By the end of this chapter I can integrate a vector function component-wise, including definite integrals, and determine constants of integration from initial conditions.
\item By the end of this chapter I can obtain velocity and displacement vectors from a given acceleration vector by integration.
\item By the end of this chapter I can obtain velocity and acceleration vectors from a given displacement vector by differentiation.
\item By the end of this chapter I can use the chain rule on vectors, e.g. differentiate r(t) where components are composite functions of t.
\end{itemize}\end{multicols}}

\begin{sommaire}
\begin{multicols}{2}
\begin{enumerate}
\item Differentiating vectors: the mechanics of vector calculus
\item Physical applications of vector differentiation: velocity and acceleration
\item Integrating vectors: indefinite and definite integrals
\item Displacement, velocity and acceleration: the kinematic triangle of calculus
\item From differentiation to velocity and acceleration vectors (Lesson 88)
\item From integration to velocity and displacement vectors (Lesson 91) and synthesis
\item Integration activity
\item Methods and worked examples
\item Exercises
\item Solutions to the exercises
\item Summary sheet
\end{enumerate}
\end{multicols}
\end{sommaire}

\begin{objectifs}
\begin{itemize}
\item By the end of this chapter I can differentiate a vector given in component form with respect to a scalar parameter (usually time).
\item By the end of this chapter I can interpret the derivative of a position vector physically as velocity, and the derivative of velocity as acceleration.
\item By the end of this chapter I can integrate a vector function component-wise, including definite integrals, and determine constants of integration from initial conditions.
\item By the end of this chapter I can obtain velocity and displacement vectors from a given acceleration vector by integration.
\item By the end of this chapter I can obtain velocity and acceleration vectors from a given displacement vector by differentiation.
\item By the end of this chapter I can use the chain rule on vectors, e.g. differentiate r(t) where components are composite functions of t.
\item By the end of this chapter I can solve kinematics problems in two dimensions using i, j components (projectile-type and general plane motion).
\item By the end of this chapter I can distinguish clearly between displacement and distance, and between velocity and speed, in vector contexts.
\item By the end of this chapter I can present a complete, well-structured GCE-style solution linking r, v and a through calculus.
\end{itemize}
\end{objectifs}

\begin{prerequis}
\begin{itemize}
\item Scalar differentiation and integration from Lower Sixth pure mathematics (polynomials, trigonometric and exponential functions, chain rule).
\item Vector algebra from Lower Sixth: components in i, j (and k), magnitude of a vector, unit vectors, scalar product.
\item Basic kinematics from Lower Sixth mechanics: constant acceleration formulae (suvat) in one dimension.
\item Trigonometry: resolving quantities into perpendicular components.
\item Sketching and interpreting graphs of motion (displacement–time, velocity–time).
\end{itemize}
\end{prerequis}

\newpage

\coursec{Differentiating vectors: the mechanics of vector calculus}

A bush taxi leaves the Bafoussam motor park and the driver accelerates, brakes, and negotiates bends on the road to Bamenda. At every instant, the vehicle has a position, a velocity and an acceleration — and none of these is just a number: each has a \cle{direction} in space. If the driver only knows his speed from the speedometer, can he predict where the taxi will be ten seconds later around a curve? Clearly not: two vehicles with the same speed of \SI{60}{km.h^{-1}} but different directions will be in very different places after ten seconds. To answer such questions, physicists \cle{differentiate} and \cle{integrate vectors}, exactly as we differentiate and integrate scalars in pure mathematics. This chapter gives you the calculus of vectors that powers all of mechanics.

\courssub{Vector-valued functions: position as a function of time}

Imagine the bush taxi moving in a horizontal plane. Choose an origin $O$ (say, the motor park gate) and fixed perpendicular directions $\vec{i}$ (east) and $\vec{j}$ (north). At time $t$, the position of the taxi is described by the \cle{position vector}
\[
\vec{r}(t) = x(t)\,\vec{i} + y(t)\,\vec{j},
\]
where $x(t)$ and $y(t)$ are ordinary scalar functions of $t$. (The same ideas extend directly to three dimensions with an additional $\vec{k}$ component.)

\begin{definition}[Vector-valued function]
A \cle{vector-valued function} of a scalar parameter $t$ is a rule that assigns to each value of $t$ (in some interval) a vector
\[
\vec{r}(t) = x(t)\,\vec{i} + y(t)\,\vec{j}.
\]
The scalar functions $x(t)$ and $y(t)$ are called the \cle{components} (or coordinates) of $\vec{r}$. As $t$ varies, the tip of $\vec{r}(t)$ traces a curve in the plane called the \cle{path} or \cle{trajectory}.
\end{definition}

In mechanics the parameter $t$ is almost always \cle{time}, measured in seconds, and the components are measured in metres. But the mathematics works for any parameter — an angle, a distance, etc.

\begin{exemplebox}[A first vector function]
Let $\vec{r}(t) = 3t\,\vec{i} + (5 - t^2)\,\vec{j}$. At $t = 0$, $\vec{r}(0) = 5\vec{j}$; at $t = 1$, $\vec{r}(1) = 3\vec{i} + 4\vec{j}$; at $t = 2$, $\vec{r}(2) = 6\vec{i} + \vec{j}$. Plotting these tips gives the beginning of a curve — the path of the moving point.
\end{exemplebox}

\courssub{The derivative of a vector function}

\textbf{Intuition.} Suppose the tip of $\vec{r}(t)$ is at point $P$ at time $t$, and at point $Q$ at the later time $t + \Delta t$. The change in position is
\[
\Delta \vec{r} = \vec{r}(t+\Delta t) - \vec{r}(t) = \vec{PQ},
\]
the chord from $P$ to $Q$. The quotient $\dfrac{\Delta \vec{r}}{\Delta t}$ is the \emph{average} rate of change of position over the interval. As $\Delta t \to 0$, the point $Q$ slides along the curve towards $P$, the chord $\vec{PQ}$ swings round towards the \cle{tangent} at $P$, and the quotient settles down to a limiting vector: the derivative.

\begin{definition}[Derivative of a vector function]
The derivative of $\vec{r}(t)$ with respect to $t$ is the limit
\[
\frac{d\vec{r}}{dt} = \lim_{\Delta t \to 0} \frac{\vec{r}(t+\Delta t) - \vec{r}(t)}{\Delta t},
\]
whenever this limit exists. It is itself a vector, and it points \cle{tangent to the path} at the point $\vec{r}(t)$, in the direction of increasing $t$. The notation $\dot{\vec{r}}$ is also used for $\dfrac{d\vec{r}}{dt}$.
\end{definition}

\begin{popfigure}
\begin{center}
\begin{tikzpicture}[scale=1.1, >=stealth]
% path curve
\draw[thick, popblue] plot[smooth] coordinates {(0.2,0.3) (1.2,1.3) (2.6,1.9) (4.0,1.6) (5.2,0.9)};
% origin
\fill (0,0) circle (1.5pt) node[below left] {$O$};
% points P and Q
\fill[popdark] (1.2,1.3) circle (2pt) node[above left, popdark] {$P$};
\fill[popdark] (2.6,1.9) circle (2pt) node[above, popdark] {$Q$};
% position vectors
\draw[->, thick, popgreen] (0,0) -- (1.2,1.3) node[midway, below left, popgreen] {$\vec{r}(t)$};
\draw[->, thick, popgreen] (0,0) -- (2.6,1.9) node[midway, above right, popgreen] {$\vec{r}(t+\Delta t)$};
% chord
\draw[->, very thick, poporange] (1.2,1.3) -- (2.6,1.9) node[midway, above, poporange] {$\Delta \vec{r}$};
% tangent vector at P
\draw[->, very thick, poppurple] (1.2,1.3) -- (2.9,2.35) node[right, poppurple] {$\dfrac{d\vec{r}}{dt}$};
% tangent line (dashed)
\draw[dashed, popmuted] (0.35,0.75) -- (3.6,2.85);
\end{tikzpicture}
\end{center}
\legende{As $\Delta t \to 0$, the chord $\Delta\vec{r} = \vec{PQ}$ shrinks and turns towards the tangent: the derivative $\frac{d\vec{r}}{dt}$ is tangent to the path at $P$.}
\end{popfigure}

\textbf{The practical rule.} Because $\vec{i}$ and $\vec{j}$ are \emph{fixed} vectors, they pass through the limit unchanged. Writing out the components,
\[
\frac{\vec{r}(t+\Delta t) - \vec{r}(t)}{\Delta t}
= \frac{x(t+\Delta t) - x(t)}{\Delta t}\,\vec{i} + \frac{y(t+\Delta t) - y(t)}{\Delta t}\,\vec{j}.
\]
Letting $\Delta t \to 0$, each scalar quotient tends to the ordinary derivative of that component. Hence the fundamental working rule:

\begin{propriete}[Differentiate component by component]
If $\vec{r}(t) = x(t)\,\vec{i} + y(t)\,\vec{j}$, then
\[
\frac{d\vec{r}}{dt} = \frac{dx}{dt}\,\vec{i} + \frac{dy}{dt}\,\vec{j}.
\]
That is: \textbf{differentiate each component separately and keep $\vec{i}$ and $\vec{j}$ fixed.} The derivation from the limit definition (sketched above) is instructive and may be asked as a short proof.
\end{propriete}

\begin{methode}[How to differentiate a vector]
\begin{enumerate}
\item Identify the parameter (usually time $t$) and write $\vec{r}(t)$ in component form.
\item Differentiate the $x$-component with respect to $t$: this gives the $\vec{i}$-component of $\dfrac{d\vec{r}}{dt}$.
\item Differentiate the $y$-component with respect to $t$: this gives the $\vec{j}$-component.
\item Leave the unit vectors $\vec{i}, \vec{j}$ untouched — they are constant.
\end{enumerate}
\end{methode}

\begin{attention}[Do not differentiate the unit vectors]
A very common mistake is to apply the product rule to $x(t)\vec{i}$ and write something like $\frac{dx}{dt}\vec{i} + x\frac{d\vec{i}}{dt}$ with an unknown $\frac{d\vec{i}}{dt}$. In this course $\vec{i}$ and $\vec{j}$ are \cle{fixed} in space, so
\[
\frac{d\vec{i}}{dt} = \vec{0} \quad \text{and} \quad \frac{d\vec{j}}{dt} = \vec{0}.
\]
(Rotating unit vectors have non-zero derivatives, but that belongs to more advanced mechanics and is not in the syllabus.)
\end{attention}

\courssub{Rules of vector differentiation}

All the familiar rules of scalar calculus survive, because differentiation acts on each component separately.

\begin{propriete}[Linearity]
For any vector functions $\vec{a}(t)$, $\vec{b}(t)$ and constants $\alpha, \beta \in \mathbb{R}$,
\[
\frac{d}{dt}\big(\alpha\,\vec{a} + \beta\,\vec{b}\big) = \alpha\,\frac{d\vec{a}}{dt} + \beta\,\frac{d\vec{b}}{dt}.
\]
In words, \cle{differentiation of vectors is linear}: the derivative of a sum is the sum of the derivatives, and constant factors pass through.
\end{propriete}

\begin{propriete}[Product rule for a scalar times a vector]
If $f(t)$ is a scalar function and $\vec{r}(t)$ a vector function, then
\[
\frac{d}{dt}\big[f(t)\,\vec{r}(t)\big] = \frac{df}{dt}\,\vec{r}(t) + f(t)\,\frac{d\vec{r}}{dt}.
\]
This follows by applying the ordinary product rule to each component of $f\,\vec{r} = (f x)\vec{i} + (f y)\vec{j}$.
\end{propriete}

\begin{propriete}[Chain rule]
If $\vec{r}$ depends on a variable $u$, and $u$ itself depends on $t$ (i.e. $u = u(t)$), then
\[
\frac{d\vec{r}}{dt} = \frac{d\vec{r}}{du}\,\frac{du}{dt}.
\]
The scalar factor $\dfrac{du}{dt}$ simply multiplies the vector $\dfrac{d\vec{r}}{du}$.
\end{propriete}

\begin{exemplebox}[Using the product and chain rules]
Let $\vec{r}(t) = t^2\big(3t\,\vec{i} + 2\,\vec{j}\big)$.
\begin{itemize}
\item \emph{Product rule:} with $f(t) = t^2$ and $\vec{a}(t) = 3t\,\vec{i} + 2\vec{j}$,
\[
\frac{d\vec{r}}{dt} = 2t\big(3t\,\vec{i} + 2\vec{j}\big) + t^2\big(3\,\vec{i}\big) = 9t^2\,\vec{i} + 4t\,\vec{j}.
\]
\item \emph{Check by expanding first:} $\vec{r}(t) = 3t^3\vec{i} + 2t^2\vec{j}$, so $\frac{d\vec{r}}{dt} = 9t^2\vec{i} + 4t\vec{j}$. \checkmark
\end{itemize}
\end{exemplebox}

\courssub{Trigonometric and exponential components; higher derivatives}

Circular and oscillatory motions produce trigonometric components; growth and decay produce exponentials. Nothing new is needed — differentiate each component with the rules of pure mathematics.

\begin{exemplebox}[Differentiating a circular motion]
Let $\vec{r}(t) = \cos t\,\vec{i} + \sin t\,\vec{j}$ (the tip moves on the unit circle). Then
\[
\frac{d\vec{r}}{dt} = -\sin t\,\vec{i} + \cos t\,\vec{j}.
\]
Notice that $\vec{r} \cdot \dfrac{d\vec{r}}{dt} = -\cos t \sin t + \sin t \cos t = 0$: the derivative is \emph{perpendicular} to the position vector, exactly as a tangent to a circle is perpendicular to the radius. Differentiating again,
\[
\frac{d^2\vec{r}}{dt^2} = -\cos t\,\vec{i} - \sin t\,\vec{j} = -\vec{r}(t),
\]
so the second derivative points back towards the origin — the first hint of centripetal acceleration, which we shall meet later.
\end{exemplebox}

\begin{definition}[Second derivative]
The \cle{second derivative} of $\vec{r}(t)$ is obtained by differentiating each component twice:
\[
\frac{d^2\vec{r}}{dt^2} = \frac{d}{dt}\!\left(\frac{d\vec{r}}{dt}\right) = \frac{d^2 x}{dt^2}\,\vec{i} + \frac{d^2 y}{dt^2}\,\vec{j}.
\]
The notation $\ddot{\vec{r}}$ is also used. Higher derivatives are defined similarly.
\end{definition}

\begin{popfigure}
\begin{center}
\begin{tikzpicture}
\begin{axis}[
  width=11cm, height=6.5cm,
  xlabel={$t$}, ylabel={},
  xmin=0, xmax=6.4, ymin=-1.6, ymax=1.6,
  axis lines=middle,
  legend style={at={(0.98,0.98)}, anchor=north east, font=\small},
  samples=100, domain=0:6.28, thick
]
\addplot[popblue] {cos(deg(x))};
\addlegendentry{$x(t)=\cos t$}
\addplot[popgreen] {sin(deg(x))};
\addlegendentry{$y(t)=\sin t$}
\addplot[poporange, dashed] {-sin(deg(x))};
\addlegendentry{$x'(t)=-\sin t$}
\addplot[poppurple, dashed] {cos(deg(x))};
\addlegendentry{$y'(t)=\cos t$}
\end{axis}
\end{tikzpicture}
\end{center}
\legende{For $\vec{r}(t)=\cos t\,\vec{i}+\sin t\,\vec{j}$, each component (solid) is differentiated separately to give the components of $\frac{d\vec{r}}{dt}$ (dashed).}
\end{popfigure}

\begin{exoresolu}[Computing $\frac{d\vec{r}}{dt}$ and $\frac{d^2\vec{r}}{dt^2}$ at a given time]
A particle moves so that its position vector at time $t$ seconds is
\[
\vec{r}(t) = \big(2t^3 - 3t\big)\vec{i} + \big(4e^{-t} + t\big)\vec{j} \quad \text{(metres)}.
\]
\textbf{(a)} Find $\dfrac{d\vec{r}}{dt}$ and $\dfrac{d^2\vec{r}}{dt^2}$. \quad \textbf{(b)} Evaluate $\dfrac{d\vec{r}}{dt}$ at $t = 1$. \quad \textbf{(c)} Find the value of $t > 0$ at which the $\vec{i}$-component of $\dfrac{d\vec{r}}{dt}$ vanishes.
\tcblower
\textbf{(a)} Differentiate each component, keeping $\vec{i}$ and $\vec{j}$ fixed:
\[
\frac{d\vec{r}}{dt} = \big(6t^2 - 3\big)\vec{i} + \big(-4e^{-t} + 1\big)\vec{j}.
\]
Differentiating again:
\[
\frac{d^2\vec{r}}{dt^2} = 12t\,\vec{i} + 4e^{-t}\,\vec{j}.
\]
\textbf{(b)} At $t = 1$:
\[
\left.\frac{d\vec{r}}{dt}\right|_{t=1} = (6 - 3)\vec{i} + \left(1 - \frac{4}{e}\right)\vec{j} = 3\vec{i} + \left(1 - \frac{4}{e}\right)\vec{j} \approx 3\vec{i} - 0.47\vec{j}.
\]
\textbf{(c)} The $\vec{i}$-component of $\dfrac{d\vec{r}}{dt}$ is $6t^2 - 3$. It vanishes when
\[
6t^2 - 3 = 0 \;\Longrightarrow\; t^2 = \tfrac{1}{2} \;\Longrightarrow\; t = \frac{1}{\sqrt{2}} = \frac{\sqrt{2}}{2} \approx 0.707\ \text{s},
\]
taking the positive root since $t > 0$.
\end{exoresolu}

\begin{attention}[GCE exam tips]
\begin{itemize}
\item \textbf{State the parameter.} Always write explicitly that you differentiate \emph{with respect to time $t$}. Marks are awarded for correct notation such as $\frac{d\vec{r}}{dt}$, $\dot{\vec{r}}$.
\item \textbf{Keep the vector form.} Do not drop $\vec{i}$ or $\vec{j}$ in intermediate lines; an answer like "$6t^2 - 3 + (-4e^{-t} + 1)$" confuses components and scores no marks.
\item \textbf{Check units.} If $\vec{r}$ is in metres and $t$ in seconds, then $\frac{d\vec{r}}{dt}$ is in $\mathrm{m\,s^{-1}}$ and $\frac{d^2\vec{r}}{dt^2}$ in $\mathrm{m\,s^{-2}}$.
\item \textbf{Always write the unit vectors} in your final answer for the derivative.
\end{itemize}
\end{attention}

\begin{exercice}[Practice]
\begin{enumerate}
\item Given $\vec{r}(t) = (3t^2 - 2t)\vec{i} + (t^3 + 5)\vec{j}$, find $\dfrac{d\vec{r}}{dt}$ and $\dfrac{d^2\vec{r}}{dt^2}$, and evaluate both at $t = 2$.
\item A point moves with $\vec{r}(t) = 4\cos(2t)\,\vec{i} + 4\sin(2t)\,\vec{j}$. Show that $\dfrac{d^2\vec{r}}{dt^2} = -4\,\vec{r}(t)$, and state what this means geometrically.
\item Given $\vec{r}(t) = e^{2t}\vec{i} + \ln(1+t)\,\vec{j}$, find the time $t$ at which the $\vec{j}$-component of $\dfrac{d\vec{r}}{dt}$ equals $\tfrac{1}{2}$.
\end{enumerate}
\end{exercice}

\begin{corrige}[Exercise 1]
\begin{enumerate}
\item $\dfrac{d\vec{r}}{dt} = (6t - 2)\vec{i} + 3t^2\vec{j}$; $\dfrac{d^2\vec{r}}{dt^2} = 6\vec{i} + 6t\vec{j}$. At $t = 2$: $\dfrac{d\vec{r}}{dt} = 10\vec{i} + 12\vec{j}$ and $\dfrac{d^2\vec{r}}{dt^2} = 6\vec{i} + 12\vec{j}$.
\item $\dfrac{d\vec{r}}{dt} = -8\sin(2t)\vec{i} + 8\cos(2t)\vec{j}$; $\dfrac{d^2\vec{r}}{dt^2} = -16\cos(2t)\vec{i} - 16\sin(2t)\vec{j} = -4\big(4\cos(2t)\vec{i} + 4\sin(2t)\vec{j}\big) = -4\vec{r}(t)$. Geometrically, the second derivative always points from the particle towards the origin (opposite to $\vec{r}$).
\item $\dfrac{d\vec{r}}{dt} = 2e^{2t}\vec{i} + \dfrac{1}{1+t}\vec{j}$. Setting $\dfrac{1}{1+t} = \tfrac12$ gives $1 + t = 2$, so $t = 1$.
\end{enumerate}
\end{corrige}

\begin{aretenir}[Key points of this section]
\begin{itemize}
\item A vector-valued function $\vec{r}(t) = x(t)\vec{i} + y(t)\vec{j}$ assigns a vector to each value of the parameter $t$; its tip traces the path of the moving point.
\item The derivative is defined by $\dfrac{d\vec{r}}{dt} = \lim\limits_{\Delta t \to 0}\dfrac{\vec{r}(t+\Delta t)-\vec{r}(t)}{\Delta t}$ and is computed component by component: $\dfrac{d\vec{r}}{dt} = \dfrac{dx}{dt}\vec{i} + \dfrac{dy}{dt}\vec{j}$.
\item Geometrically, $\dfrac{d\vec{r}}{dt}$ is \cle{tangent} to the path, pointing in the direction of motion.
\item Differentiation is linear; the product rule $\frac{d}{dt}[f\vec{r}] = f'\vec{r} + f\vec{r}\,'$ and the chain rule $\frac{d\vec{r}}{dt} = \frac{d\vec{r}}{du}\cdot\frac{du}{dt}$ hold exactly as for scalars.
\item The second derivative $\dfrac{d^2\vec{r}}{dt^2}$ is obtained by differentiating each component twice.
\item The unit vectors $\vec{i}$ and $\vec{j}$ are fixed: their derivatives are zero.
\end{itemize}
\end{aretenir}

\begin{savaistu}[Why this matters for the road to Bamenda]
Every navigation system — from the GPS in a Douala taxi to the tracking of an aircraft over Mount Cameroon — performs exactly the calculations of this section. The position is stored as a vector function of time; differentiating gives velocity, differentiating again gives acceleration. In the next section we shall give these derivatives their physical names and see Newton's laws come alive through vector calculus.
\end{savaistu}

\coursec{Physical applications of vector differentiation: velocity and acceleration}

In the previous section we mastered the \emph{mechanics} of differentiating a vector function: because the basis vectors $\vec{i}$ and $\vec{j}$ are fixed in space, the derivative of $\vec{r}(t) = x(t)\,\vec{i} + y(t)\,\vec{j}$ is obtained by differentiating each component separately. We now give this operation its physical meaning. This is where vector calculus becomes \emph{kinematics} — the mathematical description of motion. Almost every mechanics question in Paper 3 of the GCE Advanced Level that involves a moving particle relies on the ideas developed in this section.

\courssub{The position vector and displacement}

Consider a particle $P$ moving in a plane. At each instant $t$ the particle occupies a point of the plane, and we describe that point by its position relative to a fixed origin $O$.

\begin{definition}[Position vector]
The \cle{position vector} of a particle $P$ at time $t$, relative to a fixed origin $O$, is
\[ \vec{r}(t) = \overrightarrow{OP} = x(t)\,\vec{i} + y(t)\,\vec{j}, \]
where $x(t)$ and $y(t)$ are the Cartesian coordinates of $P$ at time $t$. As $t$ varies, the tip of $\vec{r}(t)$ traces out the \cle{path} (or \emph{trajectory}) of the particle.
\end{definition}

\begin{definition}[Displacement]
The \cle{displacement} of the particle between times $t_1$ and $t_2$ is the vector
\[ \Delta\vec{r} = \vec{r}(t_2) - \vec{r}(t_1). \]
It is a vector joining the initial position to the final position, independent of the actual path taken.
\end{definition}

Think of a pirogue crossing the Wouri estuary: at each moment its position can be given by two numbers (distances east and north of the jetty), and as time passes these two numbers change, drawing the curve of the journey. The displacement from departure to arrival is simply the straight-line vector from the jetty to the landing point.

\courssub{Velocity: the rate of change of displacement}

\textbf{Intuition.} Suppose that between times $t$ and $t+\delta t$ the particle moves from $P$ to $Q$. Its displacement during this short interval is the vector
\[ \delta\vec{r} = \vec{r}(t+\delta t) - \vec{r}(t) = \overrightarrow{PQ}. \]
Dividing by the elapsed time $\delta t$ gives the \emph{average velocity} over the interval, a vector parallel to the chord $PQ$. As $\delta t \to 0$, the point $Q$ slides along the path towards $P$, the chord $PQ$ tends to the \emph{tangent} to the path at $P$, and the average velocity tends to the instantaneous velocity. This limiting picture is the whole idea.

\begin{definition}[Velocity]
The \cle{velocity} of a particle at time $t$ is the derivative of its position vector with respect to time:
\[ \vec{v}(t) = \frac{d\vec{r}}{dt} = \dot{x}(t)\,\vec{i} + \dot{y}(t)\,\vec{j}. \]
The velocity vector is always \textbf{tangent to the path} at the position of the particle, and points in the direction of motion.
\end{definition}

\begin{propriete}[Speed]
The \cle{speed} of the particle is the magnitude of its velocity:
\[ |\vec{v}(t)| = \sqrt{\dot{x}(t)^{2} + \dot{y}(t)^{2}}. \]
Velocity is a \textbf{vector} (it has magnitude \emph{and} direction); speed is a \textbf{scalar} (a non-negative number). Two particles can have the same speed while moving in completely different directions.
\end{propriete}

\begin{attention}[Exam vocabulary — GCE]
\begin{itemize}
\item ``Find the \emph{velocity}'' $\rightarrow$ give the vector $\vec{v} = \dot{x}\,\vec{i} + \dot{y}\,\vec{j}$.
\item ``Find the \emph{speed}'' $\rightarrow$ give the scalar $|\vec{v}| = \sqrt{\dot{x}^{2} + \dot{y}^{2}}$.
\item ``Find the \emph{magnitude of the velocity}'' $\rightarrow$ same as speed.
\end{itemize}
Writing down only $\dot{x}$ or only $\dot{y}$ when speed is asked loses the mark.
\end{attention}

\courssub{Acceleration: the rate of change of velocity}

\textbf{Intuition.} A car travelling along the winding road from Buea to Limbe may keep its speedometer fixed at $50\ \mathrm{km\,h^{-1}}$, yet its \emph{velocity} is changing continuously because its \emph{direction} is changing. Any change of velocity — in magnitude \emph{or} in direction — constitutes an acceleration.

\begin{definition}[Acceleration]
The \cle{acceleration} of a particle at time $t$ is the derivative of its velocity, i.e.\ the second derivative of its position vector:
\[ \vec{a}(t) = \frac{d\vec{v}}{dt} = \frac{d^{2}\vec{r}}{dt^{2}} = \ddot{x}(t)\,\vec{i} + \ddot{y}(t)\,\vec{j}. \]
\end{definition}

A crucial geometric fact follows. Velocity is always tangent to the path, but acceleration is \emph{not}: acceleration measures how the velocity vector itself is turning, and it points towards the \textbf{concave (inner) side} of the trajectory. Only when the path is a straight line do $\vec{v}$ and $\vec{a}$ share the same line of action.

\begin{attention}[Velocity direction $\neq$ acceleration direction]
Do not confuse the \emph{direction of motion} (the direction of $\vec{v}$, tangent to the path) with the \emph{direction of acceleration}. They coincide only for straight-line motion. On a curved path, $\vec{a}$ always points towards the inside (concave side) of the curve, even when the particle is slowing down.
\end{attention}

\begin{popfigure}
\begin{center}
\begin{tikzpicture}[scale=1.1, >=stealth]
\draw[->, gray] (-0.5,0) -- (7.5,0) node[below left] {$x$};
\draw[->, gray] (0,-0.5) -- (0,4) node[below left] {$y$};
\draw[popblue, thick, domain=0.3:6.8, samples=100] plot (\x, {2.2*sin(0.55*\x r) + 1.4});
\coordinate (P) at (3.2, {2.2*sin(0.55*3.2 r) + 1.4});
\fill[popdark] (P) circle (2pt) node[above right, black] {$P$};
\draw[->, poporange, very thick] (0,0) -- (P) node[midway, below left] {$\vec{r}$};
\draw[->, popgreen, very thick] (P) -- ++(1.6,1.29) node[right] {$\vec{v}$ (tangent)};
\draw[->, poppink, very thick] (P) -- ++(0.55,-1.05) node[below right] {$\vec{a}$ (concave side)};
\draw[dashed, popmuted] (P) -- ++(-1.2,-0.97);
\end{tikzpicture}
\end{center}
\legende{A particle path with the position vector $\vec{r}$, the velocity $\vec{v}$ tangent to the path, and the acceleration $\vec{a}$ pointing towards the concave side of the trajectory.}
\end{popfigure}

\courssub{Finding velocity and acceleration from a position vector}

The procedure is purely mechanical: differentiate each component once for $\vec{v}$, twice for $\vec{a}$, then substitute the required value of $t$.

\begin{methode}[From $\vec{r}(t)$ to $\vec{v}(t)$ and $\vec{a}(t)$]
\begin{enumerate}
\item Write $\vec{r}(t) = x(t)\,\vec{i} + y(t)\,\vec{j}$.
\item Differentiate component-wise: $\vec{v}(t) = \dot{x}(t)\,\vec{i} + \dot{y}(t)\,\vec{j}$.
\item Differentiate again: $\vec{a}(t) = \ddot{x}(t)\,\vec{i} + \ddot{y}(t)\,\vec{j}$.
\item If a specific time $t_0$ is given, substitute $t_0$ into $\vec{v}$ and $\vec{a}$.
\item For speed or magnitude of acceleration, compute the square root of the sum of squares of the components.
\end{enumerate}
\end{methode}

\begin{exoresolu}[Velocity, speed and acceleration from $\vec{r}(t)$]
A particle moves so that its position vector at time $t$ seconds is
\[ \vec{r} = (2t^{2} + 1)\,\vec{i} + (t^{3} - 4t)\,\vec{j} \ \text{metres}. \]
\textbf{(a)} Find the velocity and acceleration vectors at time $t$. \textbf{(b)} Find the velocity, the speed and the acceleration when $t = 2$.
\tcblower
\textbf{(a)} Differentiate component by component:
\[ \vec{v} = \frac{d\vec{r}}{dt} = 4t\,\vec{i} + (3t^{2} - 4)\,\vec{j}, \qquad
\vec{a} = \frac{d\vec{v}}{dt} = 4\,\vec{i} + 6t\,\vec{j}. \]
\textbf{(b)} At $t = 2$:
\[ \vec{v}(2) = 8\,\vec{i} + (12 - 4)\,\vec{j} = 8\,\vec{i} + 8\,\vec{j} \ \mathrm{m\,s^{-1}}, \]
\[ |\vec{v}(2)| = \sqrt{8^{2} + 8^{2}} = \sqrt{128} = 8\sqrt{2} \approx 11.3\ \mathrm{m\,s^{-1}}, \]
\[ \vec{a}(2) = 4\,\vec{i} + 12\,\vec{j} \ \mathrm{m\,s^{-2}}, \qquad |\vec{a}(2)| = \sqrt{16 + 144} = \sqrt{160} = 4\sqrt{10} \approx 12.6\ \mathrm{m\,s^{-2}}. \]
Note that the speed is the \emph{magnitude of the whole velocity vector}, not the value of one component.
\end{exoresolu}

\courssub{Moving parallel to a given direction}

A particle moves \textbf{parallel to the $x$-axis} exactly when its velocity has no $\vec{j}$-component, i.e.\ when $\dot{y} = 0$ (with $\dot{x} \neq 0$). Similarly, motion parallel to the $y$-axis requires $\dot{x} = 0$ (with $\dot{y} \neq 0$). More generally, motion parallel to a direction $p\,\vec{i} + q\,\vec{j}$ requires the components of $\vec{v}$ to be in the ratio $\dot{x} : \dot{y} = p : q$, i.e.\ $q\,\dot{x} = p\,\dot{y}$.

\begin{methode}[When is the particle moving parallel to an axis or a given direction?]
\begin{enumerate}
\item Differentiate $\vec{r}(t)$ to obtain $\vec{v} = \dot{x}\,\vec{i} + \dot{y}\,\vec{j}$.
\item For motion parallel to the $x$-axis, solve $\dot{y} = 0$; for the $y$-axis, solve $\dot{x} = 0$.
\item Check that the \emph{other} component is non-zero at the solution (otherwise the particle is instantaneously at rest, not moving parallel to anything).
\item For a general direction $p\,\vec{i} + q\,\vec{j}$, solve $q\,\dot{x} = p\,\dot{y}$.
\end{enumerate}
\end{methode}

\begin{exoresolu}[Moving parallel to the $x$-axis]
A particle has position vector $\vec{r} = (t^{2} - 3t)\,\vec{i} + (2t^{3} - 6t)\,\vec{j}$ metres at time $t$ s. Find the time(s) at which the particle is moving parallel to the $x$-axis, and state its position at that instant.
\tcblower
Differentiate: $\vec{v} = (2t - 3)\,\vec{i} + (6t^{2} - 6)\,\vec{j}$.

Motion parallel to the $x$-axis requires the $\vec{j}$-component to vanish:
\[ 6t^{2} - 6 = 0 \implies t^{2} = 1 \implies t = 1 \quad (t \geq 0). \]
Check the $\vec{i}$-component at $t = 1$: $2(1) - 3 = -1 \neq 0$, so the particle is genuinely moving (in the negative $x$-direction). Its position is
\[ \vec{r}(1) = (1 - 3)\,\vec{i} + (2 - 6)\,\vec{j} = -2\,\vec{i} - 4\,\vec{j} \ \text{metres}. \]
\end{exoresolu}

\courssub{Enrichment: uniform circular motion — a preview}

A beautiful application — which you will meet again in the chapter on circular motion — is the particle moving with constant angular speed $\omega$ on a circle of radius $R$ centred at $O$:
\[ \vec{r} = R\cos(\omega t)\,\vec{i} + R\sin(\omega t)\,\vec{j}. \]
Differentiating twice:
\[ \vec{v} = -R\omega\sin(\omega t)\,\vec{i} + R\omega\cos(\omega t)\,\vec{j}, \qquad
\vec{a} = -R\omega^{2}\cos(\omega t)\,\vec{i} - R\omega^{2}\sin(\omega t)\,\vec{j} = -\omega^{2}\,\vec{r}. \]
Two remarkable facts appear. First,
\[ \vec{r}\cdot\vec{v} = -R^{2}\omega\cos(\omega t)\sin(\omega t) + R^{2}\omega\sin(\omega t)\cos(\omega t) = 0, \]
so the velocity is \textbf{perpendicular to the radius} (tangent to the circle), with constant speed $|\vec{v}| = R\omega$. Second, $\vec{a} = -\omega^{2}\vec{r}$: the acceleration is directed along the radius \textbf{towards the centre} — this is the famous \cle{centripetal acceleration}, of magnitude $\omega^{2}R = \dfrac{|\vec{v}|^{2}}{R}$. Here the speed is constant, yet the acceleration is not zero: a perfect illustration that acceleration detects changes of \emph{direction} as well as changes of speed.

\begin{popfigure}
\begin{center}
\begin{tikzpicture}[scale=1.2, >=stealth]
\draw[gray] (0,0) circle (2);
\coordinate (P) at (45:2);
\fill[popdark] (0,0) circle (1.5pt) node[below left, black] {$O$};
\fill[popdark] (P) circle (2pt) node[above right, black] {$P$};
\draw[->, poporange, very thick] (0,0) -- (P) node[midway, above left] {$\vec{r}$};
\draw[->, popgreen, very thick] (P) -- ++(135:1.6) node[above left] {$\vec{v} \perp \vec{r}$};
\draw[->, poppink, very thick] (P) -- ++(225:1.4) node[below left] {$\vec{a} = -\omega^{2}\vec{r}$};
\draw[->, popblue] (0.8,0) arc (0:45:0.8) node[midway, right] {$\omega t$};
\end{tikzpicture}
\end{center}
\legende{Uniform circular motion: $\vec{v}$ is perpendicular to $\vec{r}$ (tangent to the circle) while $\vec{a}$ points towards the centre $O$.}
\end{popfigure}

\begin{savaistu}[Why this matters beyond the exam]
The same computation governs a satellite in a circular orbit around the Earth: its velocity is tangent to the orbit while its acceleration — provided by gravity — points towards the centre of the Earth. The mathematics you have just used on $\vec{r} = R(\cos\omega t\,\vec{i} + \sin\omega t\,\vec{j})$ is literally the mathematics of space navigation.
\end{savaistu}

\begin{aretenir}[Key points]
\begin{itemize}
\item Position: $\vec{r} = x\,\vec{i} + y\,\vec{j}$; velocity: $\vec{v} = \dfrac{d\vec{r}}{dt} = \dot{x}\,\vec{i} + \dot{y}\,\vec{j}$; acceleration: $\vec{a} = \dfrac{d\vec{v}}{dt} = \ddot{x}\,\vec{i} + \ddot{y}\,\vec{j}$.
\item $\vec{v}$ is tangent to the path; speed $= |\vec{v}| = \sqrt{\dot{x}^{2} + \dot{y}^{2}}$ is a scalar.
\item $\vec{a}$ points towards the concave side of the trajectory; it coincides with the direction of motion only for straight-line paths.
\item Motion parallel to the $x$-axis: solve $\dot{y} = 0$ (and check $\dot{x} \neq 0$); parallel to the $y$-axis: solve $\dot{x} = 0$ (and check $\dot{y} \neq 0$).
\item Circular motion: $\vec{r} = R(\cos\omega t\,\vec{i} + \sin\omega t\,\vec{j})$ gives $\vec{v} \perp \vec{r}$ and $\vec{a} = -\omega^{2}\vec{r}$ (centripetal).
\end{itemize}
\end{aretenir}

\begin{exercice}[Practice]
A particle moves with position vector $\vec{r} = (t^{3} - 3t)\,\vec{i} + (t^{2} + 2)\,\vec{j}$ metres at time $t$ s.
\begin{enumerate}
\item Find $\vec{v}$ and $\vec{a}$ at time $t$.
\item Find the speed of the particle when $t = 2$.
\item Find the time at which the particle is moving parallel to the $y$-axis.
\item Show that the acceleration is \textbf{not} constant in direction, but that its $\vec{j}$-component is constant. State the direction of the acceleration at $t = 0$.
\end{enumerate}
\end{exercice}

\begin{corrige}[Exercise 1]
\textbf{1.} $\vec{v} = (3t^{2} - 3)\,\vec{i} + 2t\,\vec{j}$; $\vec{a} = 6t\,\vec{i} + 2\,\vec{j}$.

\textbf{2.} At $t = 2$: $\vec{v} = 9\,\vec{i} + 4\,\vec{j}$, so $|\vec{v}| = \sqrt{81 + 16} = \sqrt{97} \approx 9.85\ \mathrm{m\,s^{-1}}$.

\textbf{3.} Parallel to the $y$-axis: $\dot{x} = 0 \Rightarrow 3t^{2} - 3 = 0 \Rightarrow t = 1$ (taking $t \geq 0$). Check: $\dot{y} = 2 \neq 0$, so the motion is genuinely vertical.

\textbf{4.} $\vec{a} = 6t\,\vec{i} + 2\,\vec{j}$. The direction of $\vec{a}$ is given by the ratio of its components: $\tan\theta = \frac{2}{6t} = \frac{1}{3t}$, which depends on $t$. Hence the direction is \textbf{not} constant. However, the $\vec{j}$-component is the constant $2\ \mathrm{m\,s^{-2}}$. At $t = 0$, $\vec{a} = 2\,\vec{j}$, so the acceleration is purely in the positive $y$-direction (parallel to the $y$-axis). As $t$ increases, the direction rotates towards the positive $x$-axis.
\end{corrige}

\coursec{Integrating vectors: indefinite and definite integrals}

In the previous section we saw that differentiation takes us \emph{down} the kinematic chain: from the position vector $\vec{r}(t)$ we obtained the velocity $\vec{v}(t) = \dot{\vec{r}}(t)$, and from the velocity we obtained the acceleration $\vec{a}(t) = \dot{\vec{v}}(t)$. In mechanics, however, the information we actually \emph{measure} is very often the acceleration (from the forces acting, through Newton's second law) or the velocity (from a speedometer, a radar gun, or a ticker-timer in the laboratory). The natural question is therefore the \emph{reverse} one: given $\vec{a}(t)$, can we recover $\vec{v}(t)$? Given $\vec{v}(t)$, can we recover $\vec{r}(t)$?

Since differentiation of a vector is done component by component, it is entirely natural that integration — the inverse operation — should also be done component by component. This section makes that idea precise, shows how the constant of integration becomes a \emph{vector}, and explains how initial conditions fix that constant. This is the mathematical engine behind Lessons 82, 83 and 91 of the syllabus.

\courssub{Definition of the integral of a vector function}

\begin{definition}[Indefinite integral of a vector function]
Let $\vec{r}(t) = x(t)\,\vec{i} + y(t)\,\vec{j}$ be a vector function of the scalar variable $t$, whose components $x(t)$ and $y(t)$ are integrable. The \cle{indefinite integral} of $\vec{r}$ with respect to $t$ is obtained by integrating each component separately:
\[
\int \vec{r}(t)\,dt = \left(\int x(t)\,dt\right)\vec{i} + \left(\int y(t)\,dt\right)\vec{j} + \vec{c},
\]
where $\vec{c}$ is a constant \emph{vector} of integration. In three dimensions, with $\vec{r}(t) = x(t)\vec{i} + y(t)\vec{j} + z(t)\vec{k}$, we simply add the term $\left(\int z(t)\,dt\right)\vec{k}$.
\end{definition}

The intuition is exactly the one you already have from scalar calculus: integration undoes differentiation. Because the unit vectors $\vec{i}, \vec{j}, \vec{k}$ are \emph{constant} (they do not depend on $t$), they behave like ordinary constants during integration and can be taken outside the integral sign:
\[
\int \bigl[x(t)\vec{i} + y(t)\vec{j}\bigr]dt = \vec{i}\int x(t)\,dt + \vec{j}\int y(t)\,dt.
\]

\begin{propriete}[The constant of integration is a vector]
When a vector function is integrated indefinitely, the constant of integration is itself a \cle{constant vector}:
\[
\vec{c} = c_1\,\vec{i} + c_2\,\vec{j} \quad (\text{and } +\,c_3\,\vec{k} \text{ in 3 dimensions}),
\]
where $c_1, c_2, c_3$ are arbitrary real constants. Indeed, each scalar component produces its own constant of integration, and grouping them component-wise gives a constant vector. This result is examinable in the sense that you must \emph{use} it correctly: a GCE examiner will penalise a scalar constant added to a vector expression.
\end{propriete}

\begin{attention}[Scalar constant instead of vector constant]
Two errors recur every year in candidates' scripts:
\begin{itemize}
\item Writing $\displaystyle\int \vec{v}\,dt = \vec{r}(t) + c$ with a \emph{scalar} $c$. A scalar cannot be added to a vector: the constant must be a vector $\vec{c}$.
\item Remembering the constant in one component only, e.g. writing $(2t^2 + c_1)\vec{i} + 3t\,\vec{j}$. \emph{Every} component carries its own constant, so the second component must read $(3t + c_2)\vec{j}$.
\end{itemize}
A quick safety check: differentiate your answer. You must recover the integrand exactly, and the constants must disappear.
\end{attention}

\courssub{Determining the constant vector from initial conditions}

An indefinite integral describes a whole \emph{family} of vector functions differing by a constant vector. Physics selects one member of the family through an \cle{initial condition}: typically the velocity or the position at time $t = 0$.

\begin{methode}[Finding the constant vector from an initial condition]
To determine the constant vector of integration:
\begin{enumerate}
\item Integrate each component of the given vector function, adding a constant vector $\vec{c} = c_1\vec{i} + c_2\vec{j}$.
\item Substitute the given value of $t$ (usually $t = 0$) into the integrated expression.
\item Set the result equal to the given initial vector.
\item Solve \emph{component by component}: equate the $\vec{i}$-components to find $c_1$, the $\vec{j}$-components to find $c_2$.
\item Write the final answer with the constants replaced by their values.
\end{enumerate}
\end{methode}

\begin{exoresolu}[Recovering position from velocity]
A particle moves in a plane with velocity
\[
\vec{v}(t) = (3t^2 - 2t)\,\vec{i} + (4t + 1)\,\vec{j} \quad \mathrm{m\,s^{-1}}.
\]
Given that the particle is at the point with position vector $\vec{r}(0) = 2\vec{i} - \vec{j}$ metres when $t = 0$, find its position vector at time $t$.
\tcblower
\textbf{Solution.} We integrate component by component:
\begin{align*}
\vec{r}(t) &= \int \vec{v}(t)\,dt = \left(\int (3t^2 - 2t)\,dt\right)\vec{i} + \left(\int (4t+1)\,dt\right)\vec{j}\\
&= (t^3 - t^2)\,\vec{i} + (2t^2 + t)\,\vec{j} + \vec{c},
\end{align*}
where $\vec{c} = c_1\vec{i} + c_2\vec{j}$. Apply the initial condition $\vec{r}(0) = 2\vec{i} - \vec{j}$:
\[
\vec{r}(0) = \vec{c} = 2\vec{i} - \vec{j} \quad\Longrightarrow\quad c_1 = 2,\quad c_2 = -1.
\]
Hence
\[
\[\boxed{\vec{r}(t) = (t^3 - t^2 + 2)\,\vec{i} + (2t^2 + t - 1)\,\vec{j}\ \mathrm{m}.}\]
\]
\textbf{Check (differentiate back):} $\dot{\vec{r}}(t) = (3t^2 - 2t)\vec{i} + (4t+1)\vec{j} = \vec{v}(t)$. \checkmark\ Also $\vec{r}(0) = 2\vec{i} - \vec{j}$ as required.
\end{exoresolu}

\courssub{Integrating vectors with polynomial, trigonometric and exponential components}

Each component is an ordinary scalar function, so all the integration techniques of Pure Mathematics apply component-wise. The table below recalls the standard results you will use constantly in this chapter.

\begin{center}
\begin{tabularx}{\linewidth}{|l|X|X|}
\hline
\rowcolor{popblueL} \textbf{Component $f(t)$} & \textbf{$\displaystyle\int f(t)\,dt$} & \textbf{Example of use} \\
\hline
$t^n$ \ ($n \neq -1$) & $\dfrac{t^{n+1}}{n+1}$ & polynomial velocity \\
\hline
$\cos(\omega t)$ & $\dfrac{1}{\omega}\sin(\omega t)$ & simple harmonic motion \\
\hline
$\sin(\omega t)$ & $-\dfrac{1}{\omega}\cos(\omega t)$ & simple harmonic motion \\
\hline
$e^{kt}$ & $\dfrac{1}{k}e^{kt}$ & motion with resistance \\
\hline
$\dfrac{1}{t}$ & $\ln|t|$ & reciprocal acceleration \\
\hline
\end{tabularx}
\end{center}

\begin{exemplebox}[Trigonometric and exponential components]
Find $\displaystyle\int \vec{a}(t)\,dt$ where $\vec{a}(t) = 4\cos(2t)\,\vec{i} + 6e^{3t}\,\vec{j}$.
\[
\int \vec{a}(t)\,dt = \left(4\cdot\tfrac{1}{2}\sin(2t)\right)\vec{i} + \left(6\cdot\tfrac{1}{3}e^{3t}\right)\vec{j} + \vec{c}
= 2\sin(2t)\,\vec{i} + 2e^{3t}\,\vec{j} + \vec{c}.
\]
Notice how the chain-rule factors $\tfrac{1}{\omega}$ and $\tfrac{1}{k}$ appear in each component independently.
\end{exemplebox}

\courssub{Definite integrals of vector functions}

\begin{definition}[Definite integral of a vector function]
The \cle{definite integral} of $\vec{r}(t) = x(t)\vec{i} + y(t)\vec{j}$ between $t = a$ and $t = b$ is evaluated component-wise:
\[
\int_a^b \vec{r}(t)\,dt = \left(\int_a^b x(t)\,dt\right)\vec{i} + \left(\int_a^b y(t)\,dt\right)\vec{j}.
\]
The result is a \emph{constant vector} — not a function of $t$ — and no constant of integration is needed.
\end{definition}

The physical meaning is fundamental and you should memorise it: \textbf{a definite integral of a rate of change gives the total change of the quantity between the two times.} In particular:
\[
\int_{t_1}^{t_2} \vec{a}(t)\,dt = \vec{v}(t_2) - \vec{v}(t_1) \qquad\text{and}\qquad
\int_{t_1}^{t_2} \vec{v}(t)\,dt = \vec{r}(t_2) - \vec{r}(t_1).
\]
The second formula says that the definite integral of velocity is the \cle{displacement} between the two times — a vector. This is precisely the content of Lessons 83 and 84.

For a single component, this has a familiar graphical interpretation: $\int_{t_1}^{t_2} v_x(t)\,dt$ is the (signed) area under the curve $v_x(t)$, and it equals the change in the $x$-coordinate of the particle.

\begin{popfigure}
\begin{center}
\begin{tikzpicture}
\begin{axis}[
    width=11cm, height=6.5cm,
    axis lines=middle,
    xlabel={$t$ (s)}, ylabel={$v_x$ (m\,s$^{-1}$)},
    xmin=0, xmax=5.2, ymin=0, ymax=6.5,
    xtick={1,2,3,4}, ytick={2,4,6},
    domain=0:5, samples=100,
    every axis x label/.style={at={(ticklabel* cs:1.02)}, anchor=west},
    every axis y label/.style={at={(ticklabel* cs:1.02)}, anchor=south},
]
\addplot[popblue, very thick] {1 + 0.8*x + 0.15*x^2};
\addplot[fill=popblueL, draw=none] {1 + 0.8*x + 0.15*x^2} \closedcycle;
\addplot[popblue, very thick] {1 + 0.8*x + 0.15*x^2};
\draw[popdark, dashed] (axis cs:1,0) -- (axis cs:1,1.95);
\draw[popdark, dashed] (axis cs:4,0) -- (axis cs:4,6.6);
\node[popdark] at (axis cs:1,-0.45) {$t_1$};
\node[popdark] at (axis cs:4,-0.45) {$t_2$};
\node[popink, align=center] at (axis cs:2.5,1.6) {$x(t_2)-x(t_1)$\\ $=\displaystyle\int_{t_1}^{t_2}\! v_x\,dt$};
\node[popblue] at (axis cs:3.6,5.6) {$v_x(t)$};
\end{axis}
\end{tikzpicture}
\end{center}
\legende{The signed area under one component curve $v_x(t)$ between $t_1$ and $t_2$ equals the change in the $x$-component of the position: $x(t_2) - x(t_1)$. The same holds independently for the $y$-component.}
\end{popfigure}

\begin{exoresolu}[Definite integral of acceleration: change of velocity]
A particle has acceleration
\[
\vec{a}(t) = 6t\,\vec{i} + (2 - 4t)\,\vec{j} \quad \mathrm{m\,s^{-2}}.
\]
\begin{enumerate}
\item Find the change in velocity between $t = 0$ and $t = 3$.
\item Given that $\vec{v}(0) = \vec{i} + 5\vec{j}$ m\,s$^{-1}$, find $\vec{v}(3)$.
\end{enumerate}
\tcblower
\textbf{Solution.} (1) The change in velocity is the definite integral of the acceleration:
\begin{align*}
\vec{v}(3) - \vec{v}(0) &= \int_0^3 \vec{a}(t)\,dt
= \left(\int_0^3 6t\,dt\right)\vec{i} + \left(\int_0^3 (2-4t)\,dt\right)\vec{j}\\
&= \Bigl[3t^2\Bigr]_0^3\,\vec{i} + \Bigl[2t - 2t^2\Bigr]_0^3\,\vec{j}\\
&= (27 - 0)\,\vec{i} + \bigl((6 - 18) - 0\bigr)\,\vec{j}
= 27\,\vec{i} - 12\,\vec{j}\ \ \mathrm{m\,s^{-1}}.
\end{align*}
(2) Therefore
\[
\vec{v}(3) = \vec{v}(0) + \bigl(27\vec{i} - 12\vec{j}\bigr) = 28\,\vec{i} - 7\,\vec{j}\ \ \mathrm{m\,s^{-1}}.
\]
Note that no constant of integration appeared anywhere: the definite integral gave the \emph{change} directly, and the known initial velocity converted that change into an actual value.
\end{exoresolu}

\begin{attention}[Definite integrals and exam strategy (GCE tip)]
When a question asks only for a \emph{change} between two times (``find the displacement between $t = 1$ and $t = 4$'', ``find the increase in velocity''), go straight to the definite integral: you avoid constants of integration entirely and save precious minutes. Reserve the indefinite-integral-plus-initial-condition method for questions asking for the \emph{general expression} $\vec{v}(t)$ or $\vec{r}(t)$. Mixing the two approaches — for instance inserting limits \emph{and} adding a constant vector — is a classic way to lose marks.
\end{attention}

\courssub{Checking by differentiating back: the inverse nature of the operations}

Differentiation and integration of vectors are inverse operations, exactly as in scalar calculus:
\[
\frac{d}{dt}\left(\int \vec{r}(t)\,dt\right) = \vec{r}(t)
\qquad\text{and}\qquad
\int \dot{\vec{r}}(t)\,dt = \vec{r}(t) + \vec{c}.
\]
This gives you a free, reliable checking procedure in the examination hall: \textbf{differentiate your final answer and verify that you recover the function you started with, and that any given initial condition is satisfied.} If either check fails, the error is almost always a missing constant, a wrong sign in a trigonometric integral, or a forgotten chain-rule factor.

\begin{exemplebox}[A complete check]
Suppose we found $\vec{r}(t) = (t^3 - t^2 + 2)\vec{i} + (2t^2 + t - 1)\vec{j}$ from $\vec{v}(t) = (3t^2 - 2t)\vec{i} + (4t+1)\vec{j}$ with $\vec{r}(0) = 2\vec{i} - \vec{j}$.
\begin{itemize}
\item \emph{Differentiation check:} $\dot{\vec{r}} = (3t^2 - 2t)\vec{i} + (4t+1)\vec{j} = \vec{v}(t)$. \checkmark
\item \emph{Initial condition check:} $\vec{r}(0) = 2\vec{i} - \vec{j}$. \checkmark
\end{itemize}
Both checks pass, so the answer is correct.
\end{exemplebox}

\begin{aretenir}[Key points of this section]
\begin{itemize}
\item Integrate a vector \cle{component by component}: $\displaystyle\int \vec{r}(t)\,dt = \left(\int x\,dt\right)\vec{i} + \left(\int y\,dt\right)\vec{j} + \vec{c}$.
\item The constant of integration is a \cle{constant vector} $\vec{c} = c_1\vec{i} + c_2\vec{j}$, fixed by an initial condition applied component by component.
\item A \cle{definite integral} needs no constant and gives a \emph{change}: $\int_{t_1}^{t_2}\vec{v}\,dt = \vec{r}(t_2) - \vec{r}(t_1)$ (displacement) and $\int_{t_1}^{t_2}\vec{a}\,dt = \vec{v}(t_2) - \vec{v}(t_1)$.
\item All scalar integration techniques (polynomial, trigonometric, exponential) apply to each component independently.
\item Always \cle{check by differentiating back} and by verifying the initial condition.
\end{itemize}
\end{aretenir}

\begin{exercice}[Practice]
\begin{enumerate}
\item Find $\displaystyle\int \bigl[(4t^3 - \sin t)\vec{i} + 3e^{2t}\vec{j}\bigr]dt$.
\item A particle has velocity $\vec{v}(t) = (2t + 1)\vec{i} + 3\cos t\,\vec{j}$ m\,s$^{-1}$ and $\vec{r}(0) = -\vec{i} + 2\vec{j}$ m. Find $\vec{r}(t)$.
\item A particle has acceleration $\vec{a}(t) = 2\vec{i} - 6t\,\vec{j}$ m\,s$^{-2}$. Find the change in velocity between $t = 1$ and $t = 4$.
\end{enumerate}
\end{exercice}

\begin{corrige}[Exercise 1]
\textbf{1.} $\displaystyle\int (4t^3 - \sin t)\,dt = t^4 + \cos t$ and $\displaystyle\int 3e^{2t}\,dt = \tfrac{3}{2}e^{2t}$, so the integral is $(t^4 + \cos t)\vec{i} + \tfrac{3}{2}e^{2t}\vec{j} + \vec{c}$.

\textbf{2.} $\vec{r}(t) = (t^2 + t)\vec{i} + 3\sin t\,\vec{j} + \vec{c}$. At $t = 0$: $\vec{r}(0) = \vec{c} = -\vec{i} + 2\vec{j}$. Hence $\vec{r}(t) = (t^2 + t - 1)\vec{i} + (3\sin t + 2)\vec{j}$ m. Check: $\dot{\vec{r}} = (2t+1)\vec{i} + 3\cos t\,\vec{j} = \vec{v}$. \checkmark

\textbf{3.} $\displaystyle\vec{v}(4) - \vec{v}(1) = \int_1^4 (2\vec{i} - 6t\vec{j})\,dt = \bigl[2t\bigr]_1^4\vec{i} - \bigl[3t^2\bigr]_1^4\vec{j} = 6\vec{i} - 45\vec{j}$ m\,s$^{-1}$.
\end{corrige}

\begin{savaistu}[Why this matters beyond the examination room]
Every time a GPS device tracks a taxi moving through Yaoundé, it performs exactly this mathematics: the device measures velocity (from Doppler shifts of satellite signals) and \emph{integrates} it to reconstruct the position vector of the vehicle. Accelerometers in smartphones and aircraft inertial navigation systems do the same with acceleration, integrating twice — acceleration to velocity, velocity to position. The ``constant vector'' of this section is, in that context, simply the known starting point of the journey.
\end{savaistu}

\coursec{Displacement, velocity and acceleration: the kinematic triangle of calculus}

In the previous section we learnt how to integrate a vector function component by component, and how to evaluate definite vector integrals. We now put this machinery to work on the central problem of kinematics: a particle moves, its position is described by a vector $\vec{r}(t)$, and we want to pass freely between \cle{position}, \cle{velocity} and \cle{acceleration}. Differentiation takes us ``down'' the chain $\vec{r} \to \vec{v} \to \vec{a}$; integration takes us back ``up'' the chain $\vec{a} \to \vec{v} \to \vec{r}$. Mastering this two-way traffic is the key to almost every mechanics question at the GCE Advanced Level.

\courssub{Displacement: the vector change of position}

Imagine a student who leaves the school gate at $A$, walks round the football field, and stops at the tap at $B$. The \emph{distance} she has walked is the length of the path she actually followed. But if we only care about the net effect of the journey --- where she ends up relative to where she started --- we draw a straight arrow from $A$ to $B$. That arrow is her \cle{displacement}.

\begin{definition}[Displacement]
If a particle has position vector $\vec{r}(t_1)$ at time $t_1$ and $\vec{r}(t_2)$ at time $t_2$, its \cle{displacement} over the interval $[t_1, t_2]$ is the vector
\[
\Delta\vec{r} = \vec{r}(t_2) - \vec{r}(t_1).
\]
Displacement is a vector: it has magnitude (the straight-line distance from start to finish) and direction. It must not be confused with the \emph{distance travelled}, which is the scalar length of the path actually followed.
\end{definition}

\begin{popfigure}
\begin{center}
\begin{tikzpicture}[scale=1.1]
\coordinate (A) at (0,0);
\coordinate (B) at (5,1.5);
\draw[very thick, popblue] (A) .. controls (1,2.2) and (2.5,-1.2) .. (3.2,1.8) .. controls (3.8,3.2) and (4.4,2.4) .. (B);
\fill[popdark] (A) circle (2pt) node[below left] {$A=\vec{r}(t_1)$};
\fill[popdark] (B) circle (2pt) node[above right] {$B=\vec{r}(t_2)$};
\draw[-{Stealth}, very thick, poporange] (A) -- (B) node[midway, below, sloped] {displacement $\Delta\vec{r}$};
\node[popblue, align=center] at (2.2,2.4) {path actually followed\\ (distance travelled)};
\end{tikzpicture}
\end{center}
\legende{The displacement $\Delta\vec{r}$ is the straight arrow from $A$ to $B$; the distance travelled is the length of the curved path.}
\end{popfigure}

\begin{attention}[Displacement is not distance]
If a particle moves out and comes back to its starting point, its displacement is $\vec{0}$ even though the distance travelled may be large. Always ask yourself: does the question want the \emph{vector} change of position, or the \emph{scalar} length of the journey?
\end{attention}

\courssub{The kinematic chain: differentiate down, integrate up}

Recall from the earlier sections that velocity and acceleration are defined by differentiation:
\[
\vec{v} = \frac{d\vec{r}}{dt} = \dot{\vec{r}}, \qquad \vec{a} = \frac{d\vec{v}}{dt} = \dot{\vec{v}} = \ddot{\vec{r}}.
\]
Integration reverses each of these steps. Applying the Fundamental Theorem of Calculus component by component gives the following result, which you should be able to state and use fluently.

\begin{propriete}[Recovering velocity and position by integration]
Let $\vec{v}_0 = \vec{v}(0)$ and $\vec{r}_0 = \vec{r}(0)$ be the velocity and position at time $t = 0$. Then for all $t \geq 0$,
\[
\vec{v}(t) = \vec{v}_0 + \int_0^t \vec{a}(\tau)\, d\tau, \qquad
\vec{r}(t) = \vec{r}_0 + \int_0^t \vec{v}(\tau)\, d\tau.
\]
\emph{Proof (examinable as a ``show that''):} since $\dfrac{d\vec{v}}{dt} = \vec{a}$, integrating both sides from $0$ to $t$ gives $\vec{v}(t) - \vec{v}(0) = \displaystyle\int_{0}^{t} \vec{a}(\tau)\, d\tau$, hence the first formula. The second follows identically from $\dfrac{d\vec{r}}{dt} = \vec{v}$.
\end{propriete}

The whole structure is summarised in the following diagram --- learn it, because it tells you at a glance which operation a question is asking for.

\begin{popfigure}
\begin{center}
\begin{tikzpicture}[scale=1.0]
\node[draw, very thick, popblue, rounded corners, fill=popblueL, minimum width=3.2cm, minimum height=1cm] (r) at (0,0) {position $\vec{r}(t)$};
\node[draw, very thick, popgreen, rounded corners, fill=popgreenL, minimum width=3.2cm, minimum height=1cm] (v) at (0,-2.4) {velocity $\vec{v}(t)$};
\node[draw, very thick, poporange, rounded corners, minimum width=3.2cm, minimum height=1cm] (a) at (0,-4.8) {acceleration $\vec{a}(t)$};
\draw[-{Stealth}, very thick, popdark] (r.west) -- ++(-1.6,0) |- (v.west) node[pos=0.25, left] {$\dfrac{d}{dt}$};
\draw[-{Stealth}, very thick, popdark] (v.west) -- ++(-0.9,0) |- (a.west) node[pos=0.25, left] {$\dfrac{d}{dt}$};
\draw[-{Stealth}, very thick, poppurple] (v.east) -- ++(0.9,0) |- (r.east) node[pos=0.25, right] {$\displaystyle\int dt$};
\draw[-{Stealth}, very thick, poppurple] (a.east) -- ++(1.6,0) |- (v.east) node[pos=0.25, right] {$\displaystyle\int dt$};
\end{tikzpicture}
\end{center}
\legende{The kinematic triangle of calculus: differentiate to go down, integrate to go up.}
\end{popfigure}

\courssub{Recovering the constant-acceleration formulae (suvat)}

In Lower Sixth you met the constant-acceleration formulae for motion in a straight line. Vector integration shows us \emph{where they come from} --- and proves that they hold component by component for motion in a plane.

Suppose $\vec{a}$ is a \emph{constant} vector. Integrate once:
\[
\vec{v}(t) = \vec{v}_0 + \int_0^t \vec{a}\, d\tau = \vec{v}_0 + \vec{a}\,t,
\]
since a constant vector comes out of the integral. Writing $\vec{u} = \vec{v}_0$, this is the vector form of $v = u + at$. Integrate again:
\[
\vec{r}(t) = \vec{r}_0 + \int_0^t (\vec{u} + \vec{a}\tau)\, d\tau = \vec{r}_0 + \vec{u}\,t + \tfrac{1}{2}\vec{a}\,t^2,
\]
the vector form of $s = ut + \tfrac{1}{2}at^{2}$. Eliminating $t$ component-wise (or dotting suitably) yields $v^{2} = u^{2} + 2as$ along each fixed direction. The ``suvat'' formulae are therefore nothing but the first two rungs of our integration ladder when $\vec{a}$ is constant.

\begin{aretenir}[Constant acceleration in vector form]
\[
\vec{v} = \vec{u} + \vec{a}t, \qquad \vec{r} = \vec{r}_0 + \vec{u}t + \tfrac{1}{2}\vec{a}t^2.
\]
These apply only when $\vec{a}$ is constant. If $\vec{a}$ depends on $t$, you must integrate the given $\vec{a}(t)$ directly.
\end{aretenir}

\courssub{Motion in a straight line: the parallel special case}

If at every instant the vectors $\vec{r} - \vec{r}_0$, $\vec{v}$ and $\vec{a}$ are all parallel to one fixed direction (say the unit vector $\vec{i}$), the motion takes place on a straight line. We may then drop the vector notation and work with signed scalars:
\[
v = \frac{dx}{dt}, \qquad a = \frac{dv}{dt}, \qquad x(t) = x_0 + \int_0^t v\,d\tau,
\]
where the \emph{sign} of $v$ tells us the direction of travel. This is exactly the one-dimensional calculus you used in Lower Sixth. The vector theory is the generalisation; the straight-line theory is the special case where every vector keeps the same line of action.

\courssub{Displacement by definite integration; displacement versus distance}

Integrating velocity over $[t_1, t_2]$ gives the displacement:
\[
\Delta\vec{r} = \int_{t_1}^{t_2} \vec{v}(t)\, dt.
\]
But the \emph{distance travelled} is the integral of the \emph{speed}:
\[
\text{distance} = \int_{t_1}^{t_2} |\vec{v}(t)|\, dt.
\]

\begin{attention}[The classic trap]
For motion that reverses direction, $\left|\displaystyle\int \vec{v}\,dt\right| \neq \displaystyle\int |\vec{v}|\,dt$. Example: a particle on a line with $v = 2t - 4$ for $0 \le t \le 4$. The displacement is $\displaystyle\int_{0}^{4} (2t-4)\,dt = \left[t^{2} - 4t\right]_{0}^{4} = 0$ m (it returns to its start). But $v < 0$ on $[0,2)$ and $v > 0$ on $(2,4]$, so the distance is $\displaystyle\int_{0}^{2} (4-2t)\,dt + \int_{2}^{4} (2t-4)\,dt = 4 + 4 = 8$ m. To find distance, first locate the times when $v$ changes sign, then split the integral.
\end{attention}

\courssub{Finding the position from the acceleration: the standard method}

\begin{methode}[From $\vec{a}(t)$ to $\vec{r}(t)$]
\begin{enumerate}
\item Integrate $\vec{a}(t)$ component-wise to get $\vec{v}(t) = \displaystyle\int \vec{a}(t)\,dt + \vec{c}_1$.
\item Use the initial velocity $\vec{v}(0) = \vec{v}_0$ to fix the constant vector $\vec{c}_1$.
\item Integrate $\vec{v}(t)$ component-wise to get $\vec{r}(t) = \displaystyle\int \vec{v}(t)\,dt + \vec{c}_2$.
\item Use the initial position $\vec{r}(0) = \vec{r}_0$ to fix $\vec{c}_2$.
\item Answer any further question (position, displacement, speed, meeting point) from the complete $\vec{r}(t)$.
\end{enumerate}
\end{methode}

\begin{exoresolu}[Position from a variable acceleration]
A particle $P$ moves in a plane with acceleration $\vec{a}(t) = (6t)\vec{i} - 4\vec{j}$ m\,s$^{-2}$. At $t = 0$, $P$ is at the point with position vector $2\vec{i} + \vec{j}$ metres and has velocity $\vec{i} + 2\vec{j}$ m\,s$^{-1}$. Find (a) the velocity and (b) the position vector of $P$ at time $t$; (c) the displacement of $P$ during the first 3 seconds.
\tcblower
\textbf{(a)} Integrate the acceleration:
\[
\vec{v}(t) = \int \vec{a}(t)\,dt = \left(3t^2\right)\vec{i} - (4t)\vec{j} + \vec{c}_1.
\]
At $t = 0$: $\vec{v}(0) = \vec{c}_1 = \vec{i} + 2\vec{j}$. Hence
\[
\vec{v}(t) = (3t^2 + 1)\vec{i} + (2 - 4t)\vec{j}.
\]
\textbf{(b)} Integrate again:
\[
\vec{r}(t) = \left(t^3 + t\right)\vec{i} + \left(2t - 2t^2\right)\vec{j} + \vec{c}_2.
\]
At $t = 0$: $\vec{r}(0) = \vec{c}_2 = 2\vec{i} + \vec{j}$. Hence
\[
\vec{r}(t) = \left(t^3 + t + 2\right)\vec{i} + \left(2t - 2t^2 + 1\right)\vec{j}.
\]
\textbf{(c)} Displacement during the first 3 s:
\[
\Delta\vec{r} = \vec{r}(3) - \vec{r}(0) = (27 + 3)\vec{i} + (6 - 18)\vec{j} = 30\vec{i} - 12\vec{j} \text{ m}.
\]
Its magnitude is $|\Delta\vec{r}| = \sqrt{30^{2} + 12^{2}} = \sqrt{1044} = 6\sqrt{29} \approx 32.3$ m.
\end{exoresolu}

\courssub{When and where do two particles meet?}

Two particles \cle{meet} (collide or are at the same point) at time $t$ if and only if their position vectors are \emph{equal as vectors}:
\[
\vec{r}_1(t) = \vec{r}_2(t).
\]
In practice, equate the $\vec{i}$-components and the $\vec{j}$-components separately. This gives two equations; a meeting occurs only if there is a value of $t$ satisfying \emph{both} simultaneously. A common error is to equate only one component --- that merely says the particles share the same $x$-coordinate, not that they are at the same point.

\begin{exoresolu}[A meeting of two particles]
Particle $A$ starts from the origin with velocity $(2\vec{i} + \vec{j})$ m\,s$^{-1}$ and moves with constant velocity. Particle $B$ starts from the point $6\vec{i} + 3\vec{j}$ metres at the same instant with velocity $(\vec{i} + \tfrac{1}{2}\vec{j})$ m\,s$^{-1}$, also constant. Show that the particles meet, and find the time and the position vector of the meeting point.
\tcblower
With constant velocity, $\vec{r}(t) = \vec{r}_0 + \vec{v}t$. Thus
\[
\vec{r}_A(t) = (2t)\vec{i} + (t)\vec{j}, \qquad \vec{r}_B(t) = (6 + t)\vec{i} + \left(3 + \tfrac{t}{2}\right)\vec{j}.
\]
Equating the $\vec{i}$-components:
\[
2t = 6 + t \quad\Longrightarrow\quad t = 6.
\]
\emph{Check} with the $\vec{j}$-components:
\[
t = 3 + \tfrac{t}{2} \quad\Longrightarrow\quad \tfrac{t}{2} = 3 \quad\Longrightarrow\quad t = 6.
\]
Both equations give the \emph{same} value $t = 6$, so the particles do meet. The meeting point is
\[
\vec{r}_A(6) = 12\vec{i} + 6\vec{j} \text{ m}.
\]
(Verification: $\vec{r}_B(6) = (6+6)\vec{i} + (3+3)\vec{j} = 12\vec{i} + 6\vec{j}$ --- identical, as required.)
\end{exoresolu}

The worked exercise above shows the essential verification step: equating components produces \emph{two} equations, and only a common solution certifies a meeting. Had the two equations given different values of $t$, the particles would not meet, even if their paths cross. In the GCE, questions are usually designed so that a meeting does occur --- but the \emph{check} is part of full marks.

\begin{attention}[Exam tip --- GCE ``show that'' questions]
When a question asks you to \emph{show that} a velocity or position formula follows from a given acceleration (or vice versa), state explicitly at each step whether you are \textbf{differentiating} or \textbf{integrating}, e.g.\ ``Integrating $\vec{a}$ with respect to $t$\ldots'' or ``Differentiating $\vec{r}$ with respect to $t$\ldots''. Examiners award method marks for this clarity, and it protects you from mixing up the direction of the chain $\vec{r} \rightleftarrows \vec{v} \rightleftarrows \vec{a}$.
\end{attention}

\begin{exercice}[Consolidation]
\begin{enumerate}
\item A particle has velocity $\vec{v}(t) = (3t^{2} - 2t)\vec{i} + 4\vec{j}$ m\,s$^{-1}$ and starts at $\vec{r}(0) = \vec{i} - \vec{j}$ m. Find $\vec{r}(t)$ and the displacement over $[0, 2]$.
\item A particle moves on a line with $v = 6 - 2t$ m\,s$^{-1}$, $0 \le t \le 5$. Find (a) the displacement, (b) the distance travelled.
\item Particle $P$ has constant acceleration $\vec{a} = 2\vec{i} - \vec{j}$ m\,s$^{-2}$, with $\vec{v}(0) = \vec{i} + 3\vec{j}$, $\vec{r}(0) = \vec{0}$. Find $\vec{r}(t)$ and the speed at $t = 2$.
\item $A$: $\vec{r}_A = (t)\vec{i} + (t^{2})\vec{j}$; $B$: $\vec{r}_B = (4 - t)\vec{i} + (4t - 4)\vec{j}$. Determine whether $A$ and $B$ meet; if so, give the time and place.
\end{enumerate}
\end{exercice}

\begin{corrige}[Exercise 1]
\textbf{1.} $\vec{r}(t) = (t^{3} - t^{2} + 1)\vec{i} + (4t - 1)\vec{j}$; $\Delta\vec{r} = \vec{r}(2) - \vec{r}(0) = 4\vec{i} + 8\vec{j}$ m.

\textbf{2.} (a) $\displaystyle\int_{0}^{5} (6-2t)\,dt = \left[6t - t^{2}\right]_{0}^{5} = 5$ m. (b) $v = 0$ at $t = 3$: distance $= \displaystyle\int_{0}^{3} (6-2t)\,dt + \left|\int_{3}^{5} (6-2t)\,dt\right| = 9 + 4 = 13$ m.

\textbf{3.} $\vec{v} = (2t+1)\vec{i} + (3 - t)\vec{j}$; $\vec{r} = (t^{2} + t)\vec{i} + \left(3t - \tfrac{t^{2}}{2}\right)\vec{j}$. At $t = 2$: $\vec{v} = 5\vec{i} + \vec{j}$, speed $= \sqrt{26} \approx 5.10$ m\,s$^{-1}$.

\textbf{4.} Equate: $t = 4 - t \Rightarrow t = 2$; check $\vec{j}$: $t^{2} = 4t - 4 \Rightarrow (t-2)^{2} = 0 \Rightarrow t = 2$. Consistent: they meet at $t = 2$ s at $2\vec{i} + 4\vec{j}$ m.
\end{corrige}

\begin{aretenir}[Key points of the section]
\begin{itemize}
\item Displacement $\Delta\vec{r} = \vec{r}(t_2) - \vec{r}(t_1)$ is a vector; distance travelled is a scalar, the integral of speed.
\item $\vec{v} = \dot{\vec{r}}$, $\vec{a} = \dot{\vec{v}}$; conversely $\vec{v} = \vec{v}_0 + \displaystyle\int_{0}^{t} \vec{a}\,d\tau$, $\vec{r} = \vec{r}_0 + \displaystyle\int_{0}^{t} \vec{v}\,d\tau$.
\item Constant acceleration: $\vec{v} = \vec{u} + \vec{a}t$, $\vec{r} = \vec{r}_0 + \vec{u}t + \tfrac12 \vec{a}t^{2}$ (the suvat formulae in vector dress).
\item Two particles meet iff $\vec{r}_1(t) = \vec{r}_2(t)$ for one common value of $t$ --- check \emph{both} components.
\end{itemize}
\end{aretenir}

\coursec{From differentiation to velocity and acceleration vectors (Lesson 88)}

In the previous section we established the \cle{kinematic triangle of calculus}: displacement, velocity and acceleration are linked by differentiation in one direction and by integration in the other. We now put the ``differentiation side'' of that triangle to work in a systematic, examination-ready way. The situation is always the same: the position vector of a particle $P$ is given as a function of time,
\[
\vec{r}(t) = x(t)\,\vec{i} + y(t)\,\vec{j},
\]
and we are asked to extract everything the motion contains: the velocity, the acceleration, the speed, the direction of motion, and the special instants at which something remarkable happens (velocity parallel to an axis, acceleration perpendicular to velocity, maximum speed, and so on). This lesson gives you a complete toolkit, illustrated by two detailed case studies and a GCE-style structured question.

\courssub{The systematic procedure: differentiate component-wise, twice}

Recall from the definition of the derivative of a vector that, because $\vec{i}$ and $\vec{j}$ are \emph{constant} unit vectors, differentiating a position vector reduces to differentiating each scalar component separately.

\begin{definition}[Velocity and acceleration from the position vector]
If a particle has position vector $\vec{r}(t) = x(t)\,\vec{i} + y(t)\,\vec{j}$ at time $t$, then its \cle{velocity} and \cle{acceleration} are
\[
\vec{v}(t) = \dot{\vec{r}} = \frac{d\vec{r}}{dt} = \dot{x}\,\vec{i} + \dot{y}\,\vec{j},
\qquad
\vec{a}(t) = \dot{\vec{v}} = \ddot{\vec{r}} = \frac{d^2\vec{r}}{dt^2} = \ddot{x}\,\vec{i} + \ddot{y}\,\vec{j}.
\]
The \cle{speed} is the magnitude of the velocity: $|\vec{v}| = \sqrt{\dot{x}^2 + \dot{y}^2}$.
\end{definition}

\begin{attention}[Dot notation is the language of mechanics]
A dot above a symbol always means differentiation \emph{with respect to time}: $\dot{x} = \dfrac{dx}{dt}$, $\ddot{x} = \dfrac{d^2x}{dt^2}$, $\dot{\vec{v}} = \dfrac{d\vec{v}}{dt}$. Never confuse $\dot{x}$ with $\dfrac{dx}{dy}$; in mechanics the independent variable is almost always $t$.
\end{attention}

\begin{methode}[From $\vec{r}(t)$ to a complete description of the motion]
\begin{enumerate}
\item \textbf{State the definitions.} Write $\vec{v} = \dfrac{d\vec{r}}{dt}$ and $\vec{a} = \dfrac{d\vec{v}}{dt}$ \emph{before} substituting (this earns method marks at the GCE).
\item \textbf{Differentiate each component} of $\vec{r}$ to get $\vec{v} = \dot{x}\,\vec{i} + \dot{y}\,\vec{j}$.
\item \textbf{Differentiate again} to get $\vec{a} = \ddot{x}\,\vec{i} + \ddot{y}\,\vec{j}$.
\item \textbf{Compute the speed} $|\vec{v}| = \sqrt{\dot{x}^2 + \dot{y}^2}$ when required, keeping exact values (surds, $\pi$) until the final numerical answer.
\item \textbf{Answer the specific question} (parallel to an axis, perpendicular vectors, angle of motion, \dots) by translating the geometric condition into an algebraic equation on the components.
\end{enumerate}
\end{methode}

\begin{attention}[Exam tip — GCE mark allocation]
In structured GCE questions, the first mark is frequently awarded for \emph{stating} $\vec{v} = \dfrac{d\vec{r}}{dt}$ (or $\vec{a} = \dfrac{d\vec{v}}{dt}$) explicitly. Writing only the final components, even if correct, can cost you the method mark. Always show the definition, then the substitution, then the result.
\end{attention}

\courssub{Case study 1: a polynomial motion}

\begin{exoresolu}[Motion with $\vec{r}(t) = (t^2 - 3t)\,\vec{i} + 2t^3\,\vec{j}$]
A particle $P$ moves in a plane so that at time $t$ seconds ($t \geq 0$) its position vector, in metres, is
\[
\vec{r}(t) = (t^2 - 3t)\,\vec{i} + 2t^3\,\vec{j}.
\]
\textbf{(a)} Find the velocity and acceleration of $P$ at time $t$.\quad
\textbf{(b)} Find the speed of $P$ when $t = 2$.\quad
\textbf{(c)} Find the time at which the velocity of $P$ is parallel to the vector $\vec{i}$.
\tcblower
\textbf{(a) Velocity and acceleration.} By definition, $\vec{v} = \dfrac{d\vec{r}}{dt}$. Differentiating each component:
\[
\dot{x} = \frac{d}{dt}\left(t^2 - 3t\right) = 2t - 3, \qquad
\dot{y} = \frac{d}{dt}\left(2t^3\right) = 6t^2.
\]
Hence
\[
\vec{v}(t) = (2t - 3)\,\vec{i} + 6t^2\,\vec{j} \ \mathrm{m\,s^{-1}}.
\]
Then $\vec{a} = \dfrac{d\vec{v}}{dt}$ gives
\[
\ddot{x} = 2, \qquad \ddot{y} = 12t
\quad\Longrightarrow\quad
\vec{a}(t) = 2\,\vec{i} + 12t\,\vec{j} \ \mathrm{m\,s^{-2}}.
\]
Note that the acceleration is \emph{not} constant here: its $\vec{j}$-component grows with time.

\textbf{(b) Speed at $t = 2$.} Substituting $t = 2$ into the velocity:
\[
\vec{v}(2) = (4 - 3)\,\vec{i} + 6(4)\,\vec{j} = \vec{i} + 24\,\vec{j}.
\]
Therefore
\[
|\vec{v}(2)| = \sqrt{1^2 + 24^2} = \sqrt{577} \approx \SI{24.0}{m.s^{-1}}.
\]
We kept the exact surd $\sqrt{577}$ and rounded only at the last step.

\textbf{(c) Velocity parallel to $\vec{i}$.} A vector is parallel to $\vec{i}$ exactly when its $\vec{j}$-component is zero (and its $\vec{i}$-component is not). We solve
\[
\dot{y} = 6t^2 = 0 \quad\Longrightarrow\quad t = 0.
\]
At $t = 0$: $\vec{v}(0) = -3\,\vec{i}$, which is indeed parallel to $\vec{i}$ (pointing in the negative $x$-direction). Hence the required time is $t = 0\ \mathrm{s}$.
\end{exoresolu}

\begin{popfigure}
\begin{center}
\begin{tikzpicture}
\begin{axis}[
    width=0.85\linewidth, height=7cm,
    xlabel={$x$ (m)}, ylabel={$y$ (m)},
    xmin=-3.2, xmax=2.5, ymin=-0.5, ymax=9.5,
    grid=both, grid style={popline!40},
    axis lines=middle,
    legend style={at={(0.98,0.02)}, anchor=south east, font=\small},
]
\addplot[popblue, very thick, variable=\t, domain=0:1.9, samples=80]
    ({t*t - 3*t}, {2*t*t*t});
\addlegendentry{trajectory of $P$}
% velocity vectors at t=0, t=1, t=1.5 (scaled by 0.18)
\draw[-{Stealth}, poporange, very thick] (axis cs:0,0) -- (axis cs:-0.54,0);
\draw[-{Stealth}, poporange, very thick] (axis cs:-2,2) -- (axis cs:-2.18,3.08);
\draw[-{Stealth}, poporange, very thick] (axis cs:-2.25,6.75) -- (axis cs:-2.25,9.18);
\fill[popdark] (axis cs:0,0) circle (1.5pt);
\fill[popdark] (axis cs:-2,2) circle (1.5pt);
\fill[popdark] (axis cs:-2.25,6.75) circle (1.5pt);
\node[font=\small, popdark, anchor=west] at (axis cs:0.05,0.15) {$t=0$};
\node[font=\small, popdark, anchor=south east] at (axis cs:-2.05,1.95) {$t=1$};
\node[font=\small, popdark, anchor=east] at (axis cs:-2.3,6.6) {$t=1.5$};
\node[font=\small, poporange, anchor=south] at (axis cs:-0.3,0.05) {$\vec{v}(0)$};
\node[font=\small, poporange, anchor=west] at (axis cs:-2.15,2.6) {$\vec{v}(1)$};
\node[font=\small, poporange, anchor=west] at (axis cs:-2.2,8.1) {$\vec{v}(1.5)$};
\end{axis}
\end{tikzpicture}
\end{center}
\legende{Trajectory of $\vec{r}(t) = (t^2 - 3t)\vec{i} + 2t^3\vec{j}$ for $0 \le t \le 1.9$, with velocity vectors at $t = 0$, $t = 1$ and $t = 1.5$ (scaled for display). Each velocity vector is tangent to the path, as it must be. At $t = 1.5$ the $\vec{i}$-component of $\vec{v}$ vanishes ($2t-3=0$), so the velocity is purely vertical there.}
\end{popfigure}

The figure illustrates a fundamental geometric fact: \emph{the velocity vector is always tangent to the trajectory}. At $t = 0$ the particle is at the origin moving purely in the $-\vec{i}$ direction (part (c) of the case study); as time passes the $\vec{j}$-component $6t^2$ dominates and the motion becomes nearly vertical. Note also that at $t = 1.5$ the particle is at $\vec{r}(1.5) = -2.25\,\vec{i} + 6.75\,\vec{j}$ and its velocity $\vec{v}(1.5) = 13.5\,\vec{j}$ is exactly parallel to $\vec{j}$, since $\dot{x} = 2t - 3 = 0$ there.

\courssub{The direction of motion: angle from the velocity components}

Suppose we want the angle $\theta$ that the motion makes with the positive $x$-axis at a given instant. The direction of motion is the direction of $\vec{v}$, so
\[
\tan\theta = \frac{\dot{y}}{\dot{x}},
\]
choosing the quadrant from the \emph{signs} of $\dot{x}$ and $\dot{y}$. Alternatively, the angle between $\vec{v}$ and a fixed vector $\vec{d}$ follows from the scalar product:
\[
\cos\theta = \frac{\vec{v} \cdot \vec{d}}{|\vec{v}|\,|\vec{d}|}.
\]

\begin{attention}[A classic trap]
The angle of motion is found from the \emph{velocity} components, $\tan\theta = \dfrac{\dot{y}}{\dot{x}}$, \textbf{never} from the position components $\dfrac{y}{x}$. The ratio $y/x$ gives the direction of the \emph{position vector} (the line from the origin to the particle), which generally differs from the direction in which the particle is actually moving. In the figure above at $t = 1$, the position vector points up-left from the origin, while the velocity points up-left but with a quite different slope: $\vec{r}(1) = -2\vec{i} + 2\vec{j}$ has slope $-1$, whereas $\vec{v}(1) = -\vec{i} + 6\vec{j}$ has slope $-6$.
\end{attention}

\begin{exemplebox}[Angle of motion at a given instant]
For the particle of Case study 1, find the angle between the velocity and the vector $\vec{i}$ when $t = 1$.

\emph{Solution.} $\vec{v}(1) = -\vec{i} + 6\vec{j}$. With $\vec{d} = \vec{i}$:
\[
\cos\theta = \frac{\vec{v}\cdot\vec{i}}{|\vec{v}|\,|\vec{i}|} = \frac{-1}{\sqrt{1 + 36}} = \frac{-1}{\sqrt{37}}
\quad\Longrightarrow\quad \theta \approx 99.5^\circ.
\]
The obtuse angle correctly reflects that the particle is moving leftwards ($\dot{x} < 0$) and upwards ($\dot{y} > 0$). Using $\tan\theta = \dot{y}/\dot{x} = -6$ with $\dot{x}<0$, $\dot{y}>0$ places $\theta$ in the second quadrant, giving the same answer.
\end{exemplebox}

\courssub{When is the acceleration perpendicular to the velocity?}

A very frequent GCE request is: ``find the time at which the acceleration is perpendicular to the velocity''. Two non-zero vectors are perpendicular if and only if their scalar product vanishes, so the condition is
\[
\vec{a} \cdot \vec{v} = 0.
\]
This condition has a beautiful physical meaning, which is examinable and worth deriving.

\begin{propriete}[$\vec{a}\cdot\vec{v} = 0$ $\Longleftrightarrow$ the speed is stationary]
Let $\vec{v}(t)$ be the velocity of a particle and $\vec{a} = \dot{\vec{v}}$ its acceleration. Then
\[
\frac{d}{dt}\left(|\vec{v}|^2\right) = 2\,\vec{a}\cdot\vec{v}.
\]
Consequently, at an instant where $\vec{v} \neq \vec{0}$,
\[
\vec{a}\cdot\vec{v} = 0 \quad\Longleftrightarrow\quad \frac{d|\vec{v}|}{dt} = 0,
\]
i.e.\ the \cle{speed} is \emph{stationary} (passing through a maximum, a minimum, or a point of inflection) at that instant. (The derivation below is instructive and its result is quotable in the examination.)
\end{propriete}

\textbf{Derivation.} Since $|\vec{v}|^2 = \vec{v}\cdot\vec{v}$, differentiating with the product rule for the scalar product:
\[
\frac{d}{dt}\left(|\vec{v}|^2\right) = \dot{\vec{v}}\cdot\vec{v} + \vec{v}\cdot\dot{\vec{v}} = 2\,\vec{a}\cdot\vec{v}.
\]
Writing $|\vec{v}| = s$, the left side is $\dfrac{d}{dt}(s^2) = 2s\,\dot{s}$, so $\dot{s} = \dfrac{\vec{a}\cdot\vec{v}}{s}$, which vanishes exactly when $\vec{a}\cdot\vec{v} = 0$ (provided $s \neq 0$). $\blacksquare$

\begin{aretenir}[Interpreting $\vec{a}\cdot\vec{v}$]
\begin{itemize}
\item $\vec{a}\cdot\vec{v} > 0$: the particle is \emph{speeding up} (acceleration has a component along the motion).
\item $\vec{a}\cdot\vec{v} < 0$: the particle is \emph{slowing down}.
\item $\vec{a}\cdot\vec{v} = 0$: the speed is momentarily stationary — the acceleration is purely ``turning'' the velocity, as in circular motion at constant speed.
\end{itemize}
\end{aretenir}

\courssub{Case study 2: trigonometric components and periodic motion}

\begin{exoresolu}[Periodic motion and maximum speed]
A particle moves so that at time $t$ seconds its position vector, in metres, is
\[
\vec{r}(t) = 3\cos t\,\vec{i} + 2\sin t\,\vec{j}, \qquad t \geq 0.
\]
\textbf{(a)} Find $\vec{v}(t)$ and $\vec{a}(t)$.\quad
\textbf{(b)} Show that the motion is periodic with period $2\pi$.\quad
\textbf{(c)} Find the maximum speed of the particle and the times at which it occurs.\quad
\textbf{(d)} Verify that at those instants $\vec{a}\cdot\vec{v} = 0$.
\tcblower
\textbf{(a)} Differentiating component-wise, $\vec{v} = \dfrac{d\vec{r}}{dt}$:
\[
\vec{v}(t) = -3\sin t\,\vec{i} + 2\cos t\,\vec{j},
\qquad
\vec{a}(t) = \frac{d\vec{v}}{dt} = -3\cos t\,\vec{i} - 2\sin t\,\vec{j}.
\]
Notice that $\vec{a}(t) = -\vec{r}(t)$: the acceleration is always directed back towards the origin. This is the vector signature of simple harmonic-type motion.

\textbf{(b)} Since $\cos(t + 2\pi) = \cos t$ and $\sin(t + 2\pi) = \sin t$, we have $\vec{r}(t + 2\pi) = \vec{r}(t)$ for all $t$: the particle retraces its path (an ellipse, since $\dfrac{x^2}{9} + \dfrac{y^2}{4} = \cos^2 t + \sin^2 t = 1$) every $2\pi$ seconds. The motion is periodic with period $2\pi\ \mathrm{s}$.

\textbf{(c)} The speed satisfies
\[
|\vec{v}|^2 = 9\sin^2 t + 4\cos^2 t = 4 + 5\sin^2 t,
\]
using $\cos^2 t = 1 - \sin^2 t$. This is maximal when $\sin^2 t = 1$, i.e.\ $t = \dfrac{\pi}{2} + k\pi$ ($k \in \mathbb{N}$). The maximum speed is
\[
|\vec{v}|_{\max} = \sqrt{4 + 5} = \SI{3}{m.s^{-1}}.
\]
(The minimum speed, $\sqrt{4} = \SI{2}{m.s^{-1}}$, occurs when $\sin t = 0$.)

\textbf{(d)} At $t = \dfrac{\pi}{2}$: $\vec{v} = -3\,\vec{i}$ and $\vec{a} = -2\,\vec{j}$, so
\[
\vec{a}\cdot\vec{v} = (-2)(-3)\ \text{? No — compute component-wise:}\ (-3\cos t)(-3\sin t) + (-2\sin t)(2\cos t).
\]
More directly with the values at $t = \dfrac{\pi}{2}$: $\vec{v} = -3\,\vec{i} + 0\,\vec{j}$ and $\vec{a} = 0\,\vec{i} - 2\,\vec{j}$, hence
\[
\vec{a}\cdot\vec{v} = (0)(-3) + (-2)(0) = 0,
\]
confirming the property: where the speed is stationary (here maximal), acceleration is perpendicular to velocity.
\end{exoresolu}

\begin{savaistu}[Why perpendicular acceleration means constant speed]
In uniform circular motion — a satellite in a circular orbit, or a stone whirled on a string at steady rate — the acceleration points towards the centre while the velocity is tangent to the circle: they are permanently perpendicular, and the speed never changes. The acceleration's only job is to \emph{bend} the direction of motion. This is exactly the content of the property $\vec{a}\cdot\vec{v} = 0 \Leftrightarrow \dfrac{d|\vec{v}|}{dt} = 0$.
\end{savaistu}

\courssub{GCE-style structured practice}

\begin{exercice}[Structured question — from $\vec{r}$ to interpretation]
A particle $P$ moves in the plane so that at time $t$ seconds its position vector, in metres, is
\[
\vec{r}(t) = \left(t^3 - 6t\right)\vec{i} + \left(4t^2\right)\vec{j}, \qquad t \geq 0.
\]
\begin{enumerate}
\item Write down the definitions of $\vec{v}$ and $\vec{a}$ in terms of $\vec{r}$, and hence find $\vec{v}(t)$ and $\vec{a}(t)$.
\item Find the speed of $P$ when $t = 1$, giving your answer in exact form.
\item Find the angle between the velocity of $P$ and the vector $\vec{j}$ when $t = 2$.
\item Find all times $t \geq 0$ at which $\vec{a} \cdot \vec{v} = 0$, and interpret the result physically.
\item Find the time at which the velocity is parallel to $\vec{j}$.
\end{enumerate}
\end{exercice}

\begin{corrige}[Exercise 1]
\textbf{1.} $\vec{v} = \dfrac{d\vec{r}}{dt}$, $\vec{a} = \dfrac{d\vec{v}}{dt}$. Differentiating component-wise:
\[
\vec{v}(t) = \left(3t^2 - 6\right)\vec{i} + 8t\,\vec{j} \ \mathrm{m\,s^{-1}},
\qquad
\vec{a}(t) = 6t\,\vec{i} + 8\,\vec{j}\ \mathrm{m\,s^{-2}}.
\]
\textbf{2.} $\vec{v}(1) = -3\,\vec{i} + 8\,\vec{j}$, so $|\vec{v}(1)| = \sqrt{9 + 64} = \sqrt{73}\ \mathrm{m\,s^{-1}}$.

\textbf{3.} $\vec{v}(2) = 6\,\vec{i} + 16\,\vec{j}$. With $\vec{d} = \vec{j}$:
\[
\cos\theta = \frac{\vec{v}\cdot\vec{j}}{|\vec{v}|\,|\vec{j}|} = \frac{16}{\sqrt{36 + 256}} = \frac{16}{\sqrt{292}} = \frac{8}{\sqrt{73}}
\quad\Longrightarrow\quad \theta \approx 20.6^\circ.
\]

\textbf{4.} Compute the scalar product:
\[
\vec{a}\cdot\vec{v} = 6t\left(3t^2 - 6\right) + 8(8t) = 18t^3 - 36t + 64t = 18t^3 + 28t = 2t\left(9t^2 + 14\right).
\]
For $t > 0$, both factors are strictly positive, so $\vec{a}\cdot\vec{v} = 0$ has \emph{no} solution with $t > 0$; the only solution is $t = 0$. At $t = 0$: $\vec{v}(0) = -6\vec{i}$, $|\vec{v}(0)| = 6\ \mathrm{m\,s^{-1}}$, and $\vec{a}(0) = 8\vec{j}$ is indeed perpendicular to $\vec{v}(0)$. Since $\vec{a}\cdot\vec{v} > 0$ for all $t > 0$, the speed is strictly increasing throughout $t > 0$; hence $t = 0$ is the instant of \emph{minimum speed}. (This part rewards careful reading: the perpendicularity occurs at the endpoint $t = 0$, not at an interior instant.)

\textbf{5.} Velocity parallel to $\vec{j}$ requires the $\vec{i}$-component to vanish:
\[
3t^2 - 6 = 0 \quad\Longrightarrow\quad t = \sqrt{2}\ \mathrm{s} \quad (t \geq 0).
\]
Check: $\vec{v}(\sqrt{2}) = 8\sqrt{2}\,\vec{j} \neq \vec{0}$, indeed parallel to $\vec{j}$.
\end{corrige}

\begin{attention}[Common mistakes to avoid]
\begin{itemize}
\item Differentiating only one component, or mixing up $\dot{x}$ and $\dot{y}$ between $\vec{v}$ and $\vec{a}$.
\item Claiming ``velocity parallel to $\vec{i}$'' means $\vec{v} = \vec{0}$ or $\dot{x} = 0$; the correct condition is $\dot{y} = 0$ (with $\dot{x} \neq 0$).
\item Using $\tan\theta = y/x$ for the direction of motion instead of $\dot{y}/\dot{x}$.
\item Rounding early: keep $\sqrt{73}$, $\sqrt{2}$, $\pi$ exact until the final line.
\item Forgetting to verify that the vector is non-zero when concluding it is ``parallel'' to a direction.
\item Forgetting units: velocity in $\mathrm{m\,s^{-1}}$, acceleration in $\mathrm{m\,s^{-2}}$, when position is in metres and time in seconds.
\end{itemize}
\end{attention}

\begin{aretenir}[Key points of Lesson 88]
\begin{itemize}
\item $\vec{v} = \dot{\vec{r}} = \dot{x}\,\vec{i} + \dot{y}\,\vec{j}$ and $\vec{a} = \ddot{\vec{r}} = \ddot{x}\,\vec{i} + \ddot{y}\,\vec{j}$: differentiate \emph{component-wise}, twice.
\item Speed $= |\vec{v}| = \sqrt{\dot{x}^2 + \dot{y}^2}$; the velocity is always tangent to the trajectory.
\item Direction of motion: $\tan\theta = \dot{y}/\dot{x}$ (velocity components only).
\item Parallel to $\vec{i}$ $\Leftrightarrow \dot{y} = 0$; parallel to $\vec{j}$ $\Leftrightarrow \dot{x} = 0$.
\item $\vec{a}\cdot\vec{v} = 0 \Leftrightarrow$ speed stationary; its sign tells you whether the particle is speeding up or slowing down.
\item State the definitions before substituting — method marks depend on it.
\end{itemize}
\end{aretenir}

\coursec{From integration to velocity and displacement vectors (Lesson 91) and synthesis}

In the previous section we travelled \emph{down} the kinematic ladder: starting from a known position vector $\vec{r}(t)$, we differentiated to obtain the velocity $\vec{v}(t)=\dot{\vec{r}}(t)$ and differentiated again to obtain the acceleration $\vec{a}(t)=\ddot{\vec{r}}(t)$. But in real mechanics problems the roles are very often reversed. A driver in Yaoundé knows the \emph{acceleration} he produces by pressing the pedal; a basketball player knows the \emph{initial velocity} she gives the ball and the acceleration of gravity that acts on it afterwards. Nature gives us accelerations and initial conditions, and asks us to reconstruct velocities and positions. That reconstruction is exactly what \cle{integration} does. This final section is therefore devoted to climbing the ladder \emph{upwards}: from $\vec{a}(t)$ to $\vec{v}(t)$, and from $\vec{v}(t)$ to $\vec{r}(t)$.

\courssub{The systematic procedure: two integrations, two constants}

Suppose a particle moves in a plane with acceleration $\vec{a}(t)$, and suppose we know its \cle{initial velocity} $\vec{u}=\vec{v}(0)$ and its \cle{initial position} $\vec{r}_0=\vec{r}(0)$. Since $\dfrac{d\vec{v}}{dt}=\vec{a}(t)$, the velocity is an antiderivative of the acceleration:
\[
\vec{v}(t)=\int \vec{a}(t)\,dt \;=\; \vec{U}(t)+\vec{c}_1,
\]
where $\vec{U}(t)$ is any antiderivative of $\vec{a}(t)$ and $\vec{c}_1$ is a constant vector. Setting $t=0$ gives $\vec{c}_1=\vec{u}-\vec{U}(0)$, so the initial condition fixes $\vec{c}_1$. Then, since $\dfrac{d\vec{r}}{dt}=\vec{v}(t)$,
\[
\vec{r}(t)=\int \vec{v}(t)\,dt \;=\; \vec{V}(t)+\vec{c}_2,
\]
and the condition $\vec{r}(0)=\vec{r}_0$ fixes the second constant vector $\vec{c}_2$.

\begin{propriete}[Constants of integration and initial conditions]
Each integration of a vector function introduces exactly \cle{one constant vector} (equivalently, one scalar constant per component). Each initial condition — a value of $\vec{v}$ or of $\vec{r}$ at a given time — fixes exactly one constant vector. Hence a motion in the plane is completely determined by the acceleration $\vec{a}(t)$ together with \textbf{two} vector data, typically $\vec{v}(0)$ and $\vec{r}(0)$.
\end{propriete}

\begin{methode}[From acceleration to the complete motion]
\begin{enumerate}
\item \textbf{Integrate} $\vec{a}(t)$ component by component to obtain $\vec{v}(t)=\int\vec{a}\,dt+\vec{c}_1$.
\item \textbf{Apply} $\vec{v}(0)=\vec{u}$ to determine $\vec{c}_1$.
\item \textbf{Integrate} the (now fully known) $\vec{v}(t)$ to obtain $\vec{r}(t)=\int\vec{v}\,dt+\vec{c}_2$.
\item \textbf{Apply} $\vec{r}(0)=\vec{r}_0$ to determine $\vec{c}_2$.
\item \textbf{Check}: differentiate $\vec{r}(t)$ twice and confirm you recover $\vec{a}(t)$, $\vec{v}(0)=\vec{u}$ and $\vec{r}(0)=\vec{r}_0$.
\end{enumerate}
Always work in this order: acceleration first, then velocity, then position — never the reverse.
\end{methode}

\begin{attention}[Where do the constants go?]
A frequent error is to ``absorb'' the initial velocity into the integrand, for example writing $\vec{v}(t)=\int\bigl(\vec{a}(t)+\vec{u}\bigr)dt$. This is wrong: the initial condition is \textbf{not} part of the acceleration. Constants are added \emph{after} integration and are then determined by substituting the given initial values. Integrate first, adjust afterwards.
\end{attention}

\begin{attention}[Check the units on every line]
Integrating with respect to time multiplies units by seconds. Acceleration in $\mathrm{m\,s^{-2}}$ integrated once gives $\mathrm{m\,s^{-1}}$ (a velocity); integrated twice gives $\mathrm{m}$ (a displacement). If after one integration you obtain a quantity still in $\mathrm{m\,s^{-2}}$, you have made an algebraic slip — stop and find it before continuing.
\end{attention}

\courssub{Case study 1: full determination of a motion}

\begin{exoresolu}[A particle with acceleration $\vec{a}=6t\,\vec{i}-2\,\vec{j}$]
A particle $P$ moves in the $xOy$ plane. At time $t$ seconds its acceleration is
\[
\vec{a}(t)=\bigl(6t\,\vec{i}-2\,\vec{j}\bigr)\ \mathrm{m\,s^{-2}}.
\]
Initially $P$ is at the origin, $\vec{r}(0)=\vec{0}$, with velocity $\vec{v}(0)=\vec{i}+\vec{j}$ (in $\mathrm{m\,s^{-1}}$). Find $\vec{v}(t)$ and $\vec{r}(t)$, and the position and speed of $P$ at $t=2$.
\tcblower
\textbf{Step 1: integrate the acceleration.}
\[
\vec{v}(t)=\int\bigl(6t\,\vec{i}-2\,\vec{j}\bigr)dt
=\bigl(3t^{2}+c_1\bigr)\vec{i}+\bigl(-2t+c_2\bigr)\vec{j}.
\]
\textbf{Step 2: apply $\vec{v}(0)=\vec{i}+\vec{j}$.} At $t=0$: $c_1=1$ and $c_2=1$. Hence
\[
\[\boxed{\;\vec{v}(t)=\bigl(3t^{2}+1\bigr)\vec{i}+\bigl(1-2t\bigr)\vec{j}\;}\]
\]
\textbf{Step 3: integrate the velocity.}
\[
\vec{r}(t)=\int\vec{v}(t)\,dt
=\bigl(t^{3}+t+d_1\bigr)\vec{i}+\bigl(t-t^{2}+d_2\bigr)\vec{j}.
\]
\textbf{Step 4: apply $\vec{r}(0)=\vec{0}$.} At $t=0$: $d_1=d_2=0$. Hence
\[
\[\boxed{\;\vec{r}(t)=\bigl(t^{3}+t\bigr)\vec{i}+\bigl(t-t^{2}\bigr)\vec{j}\;}\]
\]
\textbf{Step 5: check.} $\dot{\vec{r}}=\bigl(3t^2+1\bigr)\vec{i}+(1-2t)\vec{j}=\vec{v}$ \checkmark; $\dot{\vec{v}}=6t\,\vec{i}-2\,\vec{j}=\vec{a}$ \checkmark; $\vec{v}(0)=\vec{i}+\vec{j}$ \checkmark; $\vec{r}(0)=\vec{0}$ \checkmark.

\textbf{Values at $t=2$:}
\[
\vec{r}(2)=(8+2)\vec{i}+(2-4)\vec{j}=10\vec{i}-2\vec{j}\ \mathrm{m},
\qquad
\vec{v}(2)=13\vec{i}-3\vec{j}\ \mathrm{m\,s^{-1}}.
\]
The speed is $|\vec{v}(2)|=\sqrt{13^{2}+(-3)^{2}}=\sqrt{178}\approx 13.3\ \mathrm{m\,s^{-1}}$.
\end{exoresolu}

\courssub{Case study 2: when does the particle return to its starting point?}

A classic GCE question: the acceleration depends on $t$, and we must find the time at which the particle is back where it started. The strategy is to build $\vec{r}(t)$ completely, then solve the vector equation $\vec{r}(t)=\vec{r}_0$ — that is, solve \emph{each component} equal to its initial value and look for a common positive solution.

\begin{exoresolu}[Return to the starting point]
A particle starts from the origin with velocity $-4\vec{j}\ \mathrm{m\,s^{-1}}$. Its acceleration at time $t$ s is
\[
\vec{a}(t)=\bigl((6t-8)\,\vec{i}+2\,\vec{j}\bigr)\ \mathrm{m\,s^{-2}}.
\]
Show that the particle returns to its starting point, and find the time at which this happens.
\tcblower
\textbf{Velocity.} $\vec{v}(t)=\int\vec{a}\,dt=(3t^{2}-8t+c_1)\vec{i}+(2t+c_2)\vec{j}$. From $\vec{v}(0)=-4\vec{j}$: $c_1=0$, $c_2=-4$, so
\[
\vec{v}(t)=(3t^{2}-8t)\vec{i}+(2t-4)\vec{j}.
\]
\textbf{Position.} $\vec{r}(t)=\int\vec{v}\,dt=(t^{3}-4t^{2}+d_1)\vec{i}+(t^{2}-4t+d_2)\vec{j}$. From $\vec{r}(0)=\vec{0}$: $d_1=d_2=0$, so
\[
\vec{r}(t)=(t^{3}-4t^{2})\vec{i}+(t^{2}-4t)\vec{j}
=t^{2}(t-4)\vec{i}+t(t-4)\vec{j}.
\]
\textbf{Return condition.} We need $\vec{r}(t)=\vec{0}$, i.e. both components zero:
\[
t^{2}(t-4)=0 \quad\text{and}\quad t(t-4)=0.
\]
For $t>0$, the first equation gives $t=4$, and the second also gives $t=4$. The two conditions have the \textbf{common} positive solution $t=4$ s: the particle returns to the origin after $4$ seconds.

\textbf{Check:} $\vec{r}(4)=(64-64)\vec{i}+(16-16)\vec{j}=\vec{0}$ \checkmark.
\end{exoresolu}

\begin{attention}[``Returns to the starting point'' is a vector condition]
To return to the starting point, \textbf{both} components of $\vec{r}(t)-\vec{r}_0$ must vanish \emph{at the same time}. It is not enough for one component to return to its initial value. For example, if the $\vec{j}$-component of $\vec{r}(t)$ is $2$ m again at some instant, the particle is merely \emph{at the same height} as its starting point — it may be many metres away horizontally. Read the question carefully: ``returns to its starting point'' means $\vec{r}(t)=\vec{r}_0$ as a vector equality; ``is again at the same height'' or ``is again on the $x$-axis'' concerns only one component.
\end{attention}

\courssub{Using definite integrals directly}

When only the \emph{change} between two instants is required, it is often faster to use definite integrals, which absorb the constants automatically. Integrating $\dot{\vec{v}}=\vec{a}$ between $0$ and $T$:
\[
\[\boxed{\;\vec{v}(T)=\vec{u}+\int_{0}^{T}\vec{a}(t)\,dt\;}\]
\qquad\text{and similarly}\qquad
\[\boxed{\;\vec{r}(T)=\vec{r}_0+\int_{0}^{T}\vec{v}(t)\,dt.\;}\]
\]
These formulae say something physically transparent: the \textbf{change of velocity} equals the accumulated (integrated) acceleration, and the \textbf{change of position} equals the accumulated velocity.

\begin{exemplebox}[Quick use of the definite form]
A particle has $\vec{u}=2\vec{i}\ \mathrm{m\,s^{-1}}$ and $\vec{a}(t)=3t\,\vec{i}+4\,\vec{j}\ \mathrm{m\,s^{-2}}$. Its velocity at $T=2$ s is
\[
\vec{v}(2)=2\vec{i}+\int_{0}^{2}\bigl(3t\,\vec{i}+4\,\vec{j}\bigr)dt
=2\vec{i}+\Bigl[\tfrac{3}{2}t^{2}\Bigr]_{0}^{2}\vec{i}+\bigl[4t\bigr]_{0}^{2}\vec{j}
=2\vec{i}+6\vec{i}+8\vec{j}
=8\vec{i}+8\vec{j}\ \mathrm{m\,s^{-1}}.
\]
No constants were needed — the initial value is built into the formula.
\end{exemplebox}

\courssub{Mixed problems: reading areas under $a$--$t$ and $v$--$t$ graphs}

The definite-integral formulae have a beautiful graphical meaning. On an \cle{acceleration--time graph}, the area between the curve and the time axis from $0$ to $T$ equals $\int_0^T a\,dt$, i.e. the \textbf{change in velocity}. On a \cle{velocity--time graph}, the area equals the \textbf{change in position} (displacement). Areas below the axis count negatively. This works component by component for vector motion, and it is the key to problems where the acceleration or velocity is given \emph{piecewise} or \emph{by a graph} rather than by a formula.

\begin{popfigure}
\begin{center}
\begin{tikzpicture}[scale=1.0]
\draw[->,thick] (0,0) -- (6.4,0) node[below left] {$t$ (s)};
\draw[->,thick] (0,-1.6) -- (0,3.2) node[below left] {$a$ ($\mathrm{m\,s^{-2}}$)};
\fill[popblueL,opacity=0.7] (0,0) -- (0,1) -- (2,2) -- (2,0) -- cycle;
\fill[popblueL,opacity=0.7] (2,0) -- (2,2) -- (4,2) -- (4,0) -- cycle;
\fill[popblueL,opacity=0.7] (4,0) -- (4,2) -- (5,0) -- cycle;
\fill[popgreenL,opacity=0.8] (5,0) -- (5.5,-1) -- (5.5,0) -- cycle;
\draw[very thick,popblue] (0,1) -- (2,2) -- (4,2) -- (5.5,-1);
\draw[dashed,popmuted] (2,0) node[below] {$2$} -- (2,2);
\draw[dashed,popmuted] (4,0) node[below] {$4$} -- (4,2);
\draw[dashed,popmuted] (5,0) node[below] {$5$};
\draw[dashed,popmuted] (5.5,0) node[below] {$5.5$} -- (5.5,-1);
\node[popdark] at (1,0.45) {$A_1$};
\node[popdark] at (3,0.9) {$A_2$};
\node[popdark] at (4.35,0.55) {$A_3$};
\node[popdark] at (5.25,-0.35) {$A_4$};
\node[align=center,popink] at (3,2.9) {change in velocity $=A_1+A_2+A_3-A_4$};
\end{tikzpicture}
\end{center}
\legende{Acceleration--time graph (one component). The signed area between the curve and the time axis equals the change of velocity; the part below the axis ($A_4$) counts negatively.}
\end{popfigure}

\begin{exoresolu}[Velocity from a graph]
The $\vec{i}$-component of a particle's acceleration is shown in the figure above: $a$ rises linearly from $1$ to $2\ \mathrm{m\,s^{-2}}$ on $[0,2]$, stays at $2$ on $[2,4]$, then falls linearly, crossing the axis at $t=5$ and reaching $-1$ at $t=5.5$. Given $v(0)=3\ \mathrm{m\,s^{-1}}$, find $v(5.5)$.
\tcblower
Compute the signed area piece by piece. $A_1$ is a trapezium, $A_2$ a rectangle, $A_3$ and $A_4$ triangles:
\begin{align*}
A_1 &= \tfrac{1}{2}(1+2)\times 2 = 3, &
A_2 &= 2\times 2 = 4, \\
A_3 &= \tfrac{1}{2}\times 1\times 2 = 1, &
A_4 &= \tfrac{1}{2}\times 0.5\times 1 = 0.25.
\end{align*}
Hence
\[
v(5.5)=v(0)+A_1+A_2+A_3-A_4=3+3+4+1-0.25=10.75\ \mathrm{m\,s^{-1}}.
\]
Note that the negative part $A_4$ (acceleration below the axis, i.e. a deceleration) is \emph{subtracted}: it reduces the velocity. As a check, the net area on $[4,5.5]$ is $1-0.25=0.75$, which equals the direct integral $\tfrac{1}{2}\bigl(2+(-1)\bigr)\times 1.5=0.75$ \checkmark.
\end{exoresolu}

\courssub{Synthesis: differentiation or integration?}

By now you hold the two directions of the kinematic ladder. The whole skill is choosing the right direction at a glance.

\begin{popfigure}
\begin{center}
\begin{tikzpicture}[scale=1.0, every node/.style={align=center}]
\node[draw,rounded corners,fill=popblueL,minimum width=3.2cm,minimum height=0.9cm] (r) at (0,0) {position $\vec{r}(t)$};
\node[draw,rounded corners,fill=popgreenL,minimum width=3.2cm,minimum height=0.9cm] (v) at (5.2,0) {velocity $\vec{v}(t)$};
\node[draw,rounded corners,fill=popblueL,minimum width=3.2cm,minimum height=0.9cm] (a) at (10.4,0) {acceleration $\vec{a}(t)$};
\draw[->,very thick,poporange] (r) -- node[above] {differentiate} (v);
\draw[->,very thick,poporange] (v) -- node[above] {differentiate} (a);
\draw[->,very thick,popgreen] (v) to[bend left=25] node[below] {integrate $+\ \vec{r}(0)$} (r);
\draw[->,very thick,popgreen] (a) to[bend left=25] node[below] {integrate $+\ \vec{v}(0)$} (v);
\node[popink] at (5.2,-2.1) {Given $\vec{r}$ or $\vec{v}$? \textbf{Differentiate.}\quad Given $\vec{a}$? \textbf{Integrate} (twice for $\vec{r}$).};
\end{tikzpicture}
\end{center}
\legende{The kinematic ladder. Going right costs nothing; going left costs one initial condition per integration.}
\end{popfigure}

\begin{aretenir}[Choosing the right tool]
\begin{itemize}
\item Given $\vec{r}(t)$, asked for $\vec{v}$ or $\vec{a}$: \textbf{differentiate} (once or twice).
\item Given $\vec{v}(t)$, asked for $\vec{a}$: \textbf{differentiate}; asked for $\vec{r}$: \textbf{integrate} and use $\vec{r}(0)$.
\item Given $\vec{a}(t)$, asked for $\vec{v}$ or $\vec{r}$: \textbf{integrate} (once or twice), using $\vec{v}(0)$ then $\vec{r}(0)$.
\item Given a \textbf{graph} of $a(t)$ or $v(t)$: use \textbf{areas} (signed) for changes of velocity or displacement.
\item Only a \emph{change} between two times is asked: use the \textbf{definite-integral} forms $\vec{v}(T)=\vec{u}+\int_0^T\vec{a}\,dt$, $\vec{r}(T)=\vec{r}_0+\int_0^T\vec{v}\,dt$.
\end{itemize}
\end{aretenir}

\courssub{Chapter summary and GCE examination technique}

\begin{aretenir}[Summary of the chapter]
\begin{itemize}
\item Vectors are differentiated and integrated \textbf{component by component}; $\vec{i},\vec{j}$ behave as constants.
\item $\vec{v}=\dot{\vec{r}}$, $\vec{a}=\dot{\vec{v}}=\ddot{\vec{r}}$; conversely $\vec{v}(t)=\vec{u}+\int_0^t\vec{a}\,ds$, $\vec{r}(t)=\vec{r}_0+\int_0^t\vec{v}\,ds$.
\item Each integration introduces one constant vector; each initial condition fixes one.
\item Area under $a$--$t$ graph $=$ change in velocity; area under $v$--$t$ graph $=$ displacement.
\item Units: $\vec{r}$ in m, $\vec{v}$ in $\mathrm{m\,s^{-1}}$, $\vec{a}$ in $\mathrm{m\,s^{-2}}$; integration with respect to time multiplies units by s, differentiation divides by s.
\end{itemize}
\end{aretenir}

\begin{methode}[GCE examination technique for vector kinematics]
\begin{enumerate}
\item \textbf{Declare your notation}: write $\vec{v}(t)=\int\vec{a}\,dt$ before computing, and keep $\vec{i},\vec{j}$ visible throughout.
\item \textbf{Show the constants}: write $+\,c_1$, $+\,c_2$ (or $+\vec{c}$), then a separate line ``$\vec{v}(0)=\ldots$ gives $c_1=\ldots$''. Examiners award method marks for exactly these lines.
\item \textbf{Carry the units}: state final answers as e.g. $\vec{r}(3)=(\ldots)\ \mathrm{m}$, and check the unit at each integration step.
\item \textbf{``Show that the particle passes through \ldots''}: do \emph{not} solve blindly — substitute the candidate time into your $\vec{r}(t)$ and verify the equality component by component.
\item \textbf{Verify by inverse operations}: differentiate your final $\vec{r}(t)$ twice; you must recover the given $\vec{a}(t)$, and $\vec{v}(0)$, $\vec{r}(0)$ must match the data. This 30-second check catches most lost marks.
\end{enumerate}
\end{methode}

\begin{exercice}[Practice]
\begin{enumerate}
\item A particle starts from rest at the point $\vec{i}+2\vec{j}$ m and moves with acceleration $\vec{a}(t)=\bigl(4\,\vec{i}+6t\,\vec{j}\bigr)\ \mathrm{m\,s^{-2}}$. Find $\vec{v}(t)$ and $\vec{r}(t)$, and its distance from the origin at $t=3$.
\item A particle has $\vec{v}(0)=5\vec{j}\ \mathrm{m\,s^{-1}}$, $\vec{r}(0)=\vec{0}$ and $\vec{a}(t)=\bigl(2t\,\vec{i}-10\,\vec{j}\bigr)\ \mathrm{m\,s^{-2}}$. Show that the particle is again on the $x$-axis when $t=1$ s, and find its position at that instant.
\item The velocity of a particle is given by a graph of $v(t)$: $v$ increases linearly from $0$ to $6\ \mathrm{m\,s^{-1}}$ over $[0,4]$ s, then decreases linearly to $0$ at $t=10$ s. Find the total displacement.
\end{enumerate}
\end{exercice}

\begin{corrige}[Exercise 1]
$\vec{v}(t)=\int\vec{a}\,dt=(4t+c_1)\vec{i}+(3t^{2}+c_2)\vec{j}$; rest at $t=0$ gives $c_1=c_2=0$, so $\vec{v}(t)=4t\,\vec{i}+3t^{2}\vec{j}$. Then $\vec{r}(t)=\int\vec{v}\,dt=(2t^{2}+d_1)\vec{i}+(t^{3}+d_2)\vec{j}$; $\vec{r}(0)=\vec{i}+2\vec{j}$ gives $d_1=1$, $d_2=2$. So $\vec{r}(t)=(2t^{2}+1)\vec{i}+(t^{3}+2)\vec{j}$. At $t=3$: $\vec{r}(3)=19\vec{i}+29\vec{j}$ m, and the distance from the origin is $\sqrt{19^{2}+29^{2}}=\sqrt{1202}\approx 34.7$ m.
\end{corrige}

\begin{corrige}[Exercise 2]
$\vec{v}(t)=\int\vec{a}\,dt=(t^{2}+c_1)\vec{i}+(-10t+c_2)\vec{j}$; $\vec{v}(0)=5\vec{j}$ gives $c_1=0$, $c_2=5$, so $\vec{v}(t)=t^{2}\vec{i}+(5-10t)\vec{j}$. Then $\vec{r}(t)=\int\vec{v}\,dt=\tfrac{t^{3}}{3}\vec{i}+(5t-5t^{2})\vec{j}$ after applying $\vec{r}(0)=\vec{0}$. On the $x$-axis the $\vec{j}$-component vanishes: $5t-5t^{2}=5t(1-t)=0$, so $t=1$ s (for $t>0$), as required. Position then: $\vec{r}(1)=\tfrac{1}{3}\vec{i}$ m.
\end{corrige}

\begin{corrige}[Exercise 3]
The displacement is the area under the $v$--$t$ graph: a triangle of base $10$ s and height $6\ \mathrm{m\,s^{-1}}$, giving $\tfrac12\times 10\times 6=30$ m.
\end{corrige}

\begin{savaistu}[Why this chapter matters beyond the GCE]
Everything you will meet in projectile motion, circular motion and simple harmonic motion later this year is built on this single ladder: $\vec{r}\to\vec{v}\to\vec{a}$ by differentiation, and back by integration with initial conditions. When engineers model the trajectory of a rocket or the suspension of a vehicle on the Douala--Yaoundé highway, they solve exactly the problems you have just solved — only with more components and messier functions. Master the ladder now, and the rest of mechanics becomes a matter of climbing it calmly, one rung at a time.
\end{savaistu}

\newpage

\coursec{Integration activity}

\coursec{Integration Activity: Vector Calculus in Mechanics}

\courssub{Aims of the activity}
This integration activity is designed to consolidate all the vector calculus skills developed in Chapter PMM08. By working through a realistic two‑dimensional motion problem set on the Wouri estuary near Douala, you will:
\begin{itemize}
    \item Differentiate a vector function to obtain velocity and acceleration, and interpret the physical meaning of a constant acceleration vector.
    \item Compute speed and direction of motion from the velocity vector.
    \item Distinguish between displacement (a vector) and distance travelled (a scalar) by comparing the magnitude of the displacement with the path length.
    \item Integrate a constant acceleration vector to recover velocity and position vectors, applying initial conditions.
    \item Solve a vector equation to determine whether two moving objects ever meet.
    \item Present a complete solution in the rigorous, step‑by‑step format expected in the GCE Advanced Level examination.
\end{itemize}
The activity is to be carried out in small groups; each group will present its solution on the board for whole‑class remediation.

\courssub{Situation}
\textbf{Tracking a pirogue on the Wouri estuary}

Near Douala, the Wouri estuary provides a natural waterway for pirogues (traditional wooden canoes) transporting goods and passengers. A team of hydrographers models the motion of a pirogue crossing the estuary. They choose a coordinate system with origin at the jetty: the \textbf{i}‑axis points East and the \textbf{j}‑axis points North. Distances are measured in kilometres (km) and time in hours (h).

For the first pirogue, the position vector relative to the jetty is given by
\[
\mathbf{r}_1(t) = (2t + 0.5t^2)\,\mathbf{i} + (3t - t^2)\,\mathbf{j}, \qquad 0 \le t \le 2.
\]

A second pirogue leaves the same jetty at the same instant ($t=0$) but with a different motion: its acceleration is constant and given by
\[
\mathbf{a}_2 = \mathbf{i} - 2\,\mathbf{j} \quad (\mathrm{km\,h^{-2}}),
\]
and it starts from rest at the jetty.

The hydrographers need to analyse both motions in order to predict possible meeting points and to estimate travel distances.

\begin{popfigure}
\begin{tikzpicture}
\begin{axis}[
    axis lines=middle,
    xlabel={$x$ (km)},
    ylabel={$y$ (km)},
    xmin=-1, xmax=7,
    ymin=-5, ymax=5,
    xtick={0,1,...,6},
    ytick={-4,-2,0,2,4},
    grid=major,
    width=10cm,
    height=8cm,
    legend pos=north west,
]
\addplot[domain=0:2, samples=50, popblue, thick] ({2*x+0.5*x^2}, {3*x-x^2});
\addlegendentry{$\mathbf{r}_1(t)$: Pirogue 1}
\addplot[domain=0:2, samples=50, poporange, thick] ({0.5*x^2}, {-x^2});
\addlegendentry{$\mathbf{r}_2(t)$: Pirogue 2}
\addplot[only marks, mark=*, popink] coordinates {(0,0)};
\addlegendentry{Jetty (origin)}
\end{axis}
\end{tikzpicture}
\legende{Trajectories of the two pirogues on the Wouri estuary for $0\le t\le 2$ h.}
\end{popfigure}

\courssub{Tasks}
Work through the following questions in order. Show all steps clearly, using proper vector notation and GCE‑style presentation.

\begin{enumerate}
    \item \textbf{Velocity and acceleration of the first pirogue}
    \begin{enumerate}
        \item Find the velocity vector $\mathbf{v}_1(t)$ and the acceleration vector $\mathbf{a}_1(t)$ by differentiating $\mathbf{r}_1(t)$.
        \item Is the acceleration constant? Justify your answer.
    \end{enumerate}

    \item \textbf{Speed and direction at $t=1$ h}
    \begin{enumerate}
        \item Compute the speed of the first pirogue at $t=1$ h.
        \item Determine the direction of motion at that instant, expressed as an angle measured from East (the positive \textbf{i}‑axis) towards North (the positive \textbf{j}‑axis). Give your answer in degrees to one decimal place.
    \end{enumerate}

    \item \textbf{Displacement versus distance travelled (first pirogue)}
    \begin{enumerate}
        \item Find the displacement vector $\Delta\mathbf{r}_1$ of the first pirogue between $t=0$ and $t=2$ h.
        \item Calculate the magnitude of this displacement.
        \item Without performing the exact integration for the path length, discuss qualitatively whether the actual distance travelled by the pirogue is greater than, equal to, or less than the magnitude of the displacement. Explain your reasoning.
    \end{enumerate}

    \item \textbf{Motion of the second pirogue (integration)}
    \begin{enumerate}
        \item Given $\mathbf{a}_2 = \mathbf{i} - 2\,\mathbf{j}$ and the initial conditions $\mathbf{v}_2(0)=\mathbf{0}$, $\mathbf{r}_2(0)=\mathbf{0}$, find the velocity vector $\mathbf{v}_2(t)$ and the position vector $\mathbf{r}_2(t)$ by integration.
        \item Write the position vector in the form $\mathbf{r}_2(t) = x_2(t)\,\mathbf{i} + y_2(t)\,\mathbf{j}$.
    \end{enumerate}

    \item \textbf{Do the pirogues meet?}
    Determine whether there exists a time $t \in [0,2]$ such that $\mathbf{r}_1(t) = \mathbf{r}_2(t)$. If they meet, give the time and the position; if not, explain why.
\end{enumerate}

\courssub{Model answer}
\begin{exemplebox}[Complete worked solution]

\textbf{1. Velocity and acceleration of the first pirogue}

\begin{align*}
\mathbf{r}_1(t) &= (2t + 0.5t^2)\,\mathbf{i} + (3t - t^2)\,\mathbf{j} \\[4pt]
\mathbf{v}_1(t) &= \frac{d\mathbf{r}_1}{dt} = \frac{d}{dt}(2t + 0.5t^2)\,\mathbf{i} + \frac{d}{dt}(3t - t^2)\,\mathbf{j} \\
                &= (2 + t)\,\mathbf{i} + (3 - 2t)\,\mathbf{j} \quad (\mathrm{km\,h^{-1}}) \\[4pt]
\mathbf{a}_1(t) &= \frac{d\mathbf{v}_1}{dt} = \frac{d}{dt}(2 + t)\,\mathbf{i} + \frac{d}{dt}(3 - 2t)\,\mathbf{j} \\
                &= 1\,\mathbf{i} + (-2)\,\mathbf{j} = \mathbf{i} - 2\,\mathbf{j} \quad (\mathrm{km\,h^{-2}})
\end{align*}

The acceleration vector $\mathbf{a}_1(t) = \mathbf{i} - 2\,\mathbf{j}$ does not depend on $t$; therefore \textbf{the acceleration is constant}.

\vspace{4pt}
\textbf{2. Speed and direction at $t=1$ h}

At $t=1$:
\[
\mathbf{v}_1(1) = (2+1)\,\mathbf{i} + (3-2)\,\mathbf{j} = 3\,\mathbf{i} + 1\,\mathbf{j} \quad (\mathrm{km\,h^{-1}})
\]

Speed is the magnitude of the velocity vector:
\[
\text{Speed} = |\mathbf{v}_1(1)| = \sqrt{3^2 + 1^2} = \sqrt{10} \approx 3.16\ \mathrm{km\,h^{-1}}.
\]

Direction: the angle $\theta$ measured from East (positive \textbf{i}) towards North (positive \textbf{j}) satisfies
\[
\tan\theta = \frac{v_y}{v_x} = \frac{1}{3} \quad\Rightarrow\quad \theta = \arctan\left(\frac{1}{3}\right) \approx 18.4^\circ.
\]
Hence the pirogue is moving at about \textbf{3.16 km\,h$^{-1}$} in a direction \textbf{18.4° North of East}.

\vspace{4pt}
\textbf{3. Displacement versus distance travelled (first pirogue)}

Displacement between $t=0$ and $t=2$:
\[
\mathbf{r}_1(0) = 0\,\mathbf{i} + 0\,\mathbf{j} = \mathbf{0}, \qquad
\mathbf{r}_1(2) = (2\cdot2 + 0.5\cdot4)\,\mathbf{i} + (3\cdot2 - 4)\,\mathbf{j} = (4+2)\,\mathbf{i} + (6-4)\,\mathbf{j} = 6\,\mathbf{i} + 2\,\mathbf{j}.
\]
\[
\Delta\mathbf{r}_1 = \mathbf{r}_1(2) - \mathbf{r}_1(0) = 6\,\mathbf{i} + 2\,\mathbf{j} \quad (\mathrm{km}).
\]

Magnitude of displacement:
\[
|\Delta\mathbf{r}_1| = \sqrt{6^2 + 2^2} = \sqrt{40} = 2\sqrt{10} \approx 6.32\ \mathrm{km}.
\]

\textbf{Qualitative comparison:} The velocity vector $\mathbf{v}_1(t) = (2+t)\,\mathbf{i} + (3-2t)\,\mathbf{j}$ is never zero on $[0,2]$ ($v_x>0$ always, and $v_y=0$ only at $t=1.5$ where $v_x\neq0$). The path is a smooth curve (a parabola in the $xy$‑plane). For any curved path, the distance travelled (arc length) is strictly greater than the straight‑line distance between the endpoints (the magnitude of the displacement). Therefore \textbf{the actual distance travelled is greater than 6.32 km}.

\vspace{4pt}
\textbf{4. Motion of the second pirogue (integration)}

Given constant acceleration $\mathbf{a}_2 = \mathbf{i} - 2\,\mathbf{j}$.

Integrate to find velocity:
\[
\mathbf{v}_2(t) = \int \mathbf{a}_2\,dt = \int (\mathbf{i} - 2\,\mathbf{j})\,dt = t\,\mathbf{i} - 2t\,\mathbf{j} + \mathbf{C}_1.
\]
Initial condition $\mathbf{v}_2(0)=\mathbf{0}$ gives $\mathbf{C}_1 = \mathbf{0}$. Hence
\[
\mathbf{v}_2(t) = t\,\mathbf{i} - 2t\,\mathbf{j} \quad (\mathrm{km\,h^{-1}}).
\]

Integrate to find position:
\[
\mathbf{r}_2(t) = \int \mathbf{v}_2(t)\,dt = \int (t\,\mathbf{i} - 2t\,\mathbf{j})\,dt = \frac{t^2}{2}\,\mathbf{i} - t^2\,\mathbf{j} + \mathbf{C}_2.
\]
Initial condition $\mathbf{r}_2(0)=\mathbf{0}$ gives $\mathbf{C}_2 = \mathbf{0}$. Thus
\[
\mathbf{r}_2(t) = \frac{t^2}{2}\,\mathbf{i} - t^2\,\mathbf{j} \quad (\mathrm{km}).
\]

In component form:
\[
x_2(t) = \frac{t^2}{2}, \qquad y_2(t) = -t^2.
\]

\vspace{4pt}
\textbf{5. Do the pirogues meet?}

We need $\mathbf{r}_1(t) = \mathbf{r}_2(t)$ for some $t\in[0,2]$. Equate components:

\begin{align*}
x\text{-component:} &\quad 2t + 0.5t^2 = \frac{t^2}{2} \quad\Rightarrow\quad 2t + 0.5t^2 = 0.5t^2 \quad\Rightarrow\quad 2t = 0 \quad\Rightarrow\quad t = 0. \\[4pt]
y\text{-component:} &\quad 3t - t^2 = -t^2 \quad\Rightarrow\quad 3t = 0 \quad\Rightarrow\quad t = 0.
\end{align*}

The only solution is $t=0$, which corresponds to the initial instant when both pirogues are at the jetty. For $t>0$, the $x$‑coordinates differ ($2t+0.5t^2 > 0.5t^2$) and the $y$‑coordinates differ ($3t-t^2 > -t^2$ because $3t>0$). Hence \textbf{the two pirogues never occupy the same position again during the interval $0\le t\le 2$}.

\end{exemplebox}

\begin{aretenir}[Key points for GCE presentation]
\begin{itemize}
    \item Always write the vector function in component form before differentiating or integrating.
    \item Show the integration constant explicitly and use initial conditions to determine it.
    \item When asked for speed, give the magnitude of the velocity vector; when asked for direction, give the angle from the positive \textbf{i}‑axis (East) towards the positive \textbf{j}‑axis (North).
    \item Displacement is a vector; distance travelled is a scalar (arc length). For any non‑straight path, distance $>$ magnitude of displacement.
    \item To check if two particles meet, equate their position vectors component‑wise and solve for $t$. Verify that the solution lies in the given time interval.
\end{itemize}
\end{aretenir}

\begin{attention}[Common mistakes to avoid]
\begin{itemize}
    \item Forgetting to include the integration constant when integrating acceleration or velocity.
    \item Confusing speed (scalar) with velocity (vector); speed is the magnitude of velocity.
    \item Giving the direction as a bearing (from North) instead of the angle from East, unless the question specifically asks for a bearing.
    \item Stating that distance travelled equals the magnitude of displacement without checking whether the path is straight.
    \item Solving only one component equation when checking for a meeting; both components must be satisfied simultaneously.
\end{itemize}
\end{attention}

\begin{savaistu}[Did you know?]
The Wouri estuary is the gateway to Douala, Cameroon's economic capital. Hydrographers and port authorities use vector calculus to model currents, plan dredging, and optimise navigation channels for ships and pirogues alike. The constant acceleration model used for the second pirogue is a simplification; in reality, water resistance and engine thrust produce variable acceleration, which would require numerical integration or differential equations — topics you will encounter at university.
\end{savaistu}

\newpage

\coursec{Methods and worked examples}

\courssub{Differentiating a position vector to obtain velocity and acceleration}

\begin{definition}[Velocity and acceleration vectors]
If a particle moves in a plane so that its position vector at time $t$ is
\[ \vec{r}(t) = x(t)\,\vec{i} + y(t)\,\vec{j}, \]
then its \cle{velocity} is the rate of change of displacement (position):
\[ \vec{v}(t) = \dot{\vec{r}}(t) = \frac{d\vec{r}}{dt} = \dot{x}(t)\,\vec{i} + \dot{y}(t)\,\vec{j}, \]
and its \cle{acceleration} is the rate of change of velocity:
\[ \vec{a}(t) = \dot{\vec{v}}(t) = \ddot{\vec{r}}(t) = \ddot{x}(t)\,\vec{i} + \ddot{y}(t)\,\vec{j}. \]
The \cle{speed} is the magnitude of the velocity: $|\vec{v}| = \sqrt{\dot{x}^2 + \dot{y}^2}$. The \cle{direction of motion} at any instant is the direction of $\vec{v}$.
\end{definition}

\begin{propriete}[Componentwise differentiation of vectors]
Since the unit vectors $\vec{i}, \vec{j}, \vec{k}$ are \emph{fixed} (constant in magnitude and direction), a vector is differentiated \cle{component by component}:
\[ \frac{d}{dt}\left[x(t)\,\vec{i} + y(t)\,\vec{j}\right] = \frac{dx}{dt}\,\vec{i} + \frac{dy}{dt}\,\vec{j}. \]
This follows directly from the sum and constant-multiple rules of scalar differentiation, because $\vec{i}$ and $\vec{j}$ behave as constants. (The proof is not examinable, but the result must be used fluently.)
\end{propriete}

\begin{methode}[Differentiating vectors component by component]
To differentiate a vector given in component form:
\begin{enumerate}
\item Write the vector in the form $\vec{r}(t) = x(t)\,\vec{i} + y(t)\,\vec{j}$ (possibly with a $z(t)\,\vec{k}$ term).
\item Differentiate \cle{each component separately} with respect to $t$, treating $\vec{i}, \vec{j}, \vec{k}$ as constants: $\vec{v} = \dot{\vec{r}} = \dot{x}\,\vec{i} + \dot{y}\,\vec{j}$.
\item Differentiate again to obtain acceleration: $\vec{a} = \dot{\vec{v}} = \ddot{x}\,\vec{i} + \ddot{y}\,\vec{j}$.
\item To find a \cle{magnitude} (speed, magnitude of acceleration), use Pythagoras: $|\vec{v}| = \sqrt{\dot{x}^2 + \dot{y}^2}$.
\item To find a \cle{direction}, compute the angle $\theta$ with $\tan\theta = \dfrac{\dot{y}}{\dot{x}}$, checking the quadrant from the signs of the components.
\item Substitute the required value of $t$ only \emph{after} differentiating.
\end{enumerate}
\textbf{Reflexes:} velocity is the derivative of displacement; acceleration is the derivative of velocity. Never differentiate the unit vectors $\vec{i}, \vec{j}$ — they are fixed.
\end{methode}

\begin{exoresolu}[Velocity and acceleration of a particle]
A particle moves so that its position vector at time $t$ seconds is
\[ \vec{r} = (3t^2 - 2t)\,\vec{i} + (t^3 + 4)\,\vec{j} \ \text{metres}. \]
Find (a) the velocity and speed when $t = 2$; (b) the acceleration vector and its magnitude when $t = 2$; (c) the time at which the acceleration is parallel to $\vec{i}$.
\tcblower
\textbf{Solution.}\\
(a) Differentiate each component:
\[ \vec{v} = \dot{\vec{r}} = (6t - 2)\,\vec{i} + 3t^2\,\vec{j}. \]
At $t = 2$: $\vec{v} = (12 - 2)\vec{i} + 12\,\vec{j} = 10\,\vec{i} + 12\,\vec{j}\ \mathrm{m\,s^{-1}}$.\\
Speed: $|\vec{v}| = \sqrt{10^2 + 12^2} = \sqrt{244} = 2\sqrt{61} \approx 15.6\ \mathrm{m\,s^{-1}}$.\\[4pt]
(b) Differentiate again:
\[ \vec{a} = \dot{\vec{v}} = 6\,\vec{i} + 6t\,\vec{j}. \]
At $t = 2$: $\vec{a} = 6\,\vec{i} + 12\,\vec{j}\ \mathrm{m\,s^{-2}}$, with magnitude $|\vec{a}| = \sqrt{36 + 144} = \sqrt{180} = 6\sqrt{5} \approx 13.4\ \mathrm{m\,s^{-2}}$.\\[4pt]
(c) $\vec{a}$ is parallel to $\vec{i}$ when its $\vec{j}$-component vanishes: $6t = 0$, so $t = 0$. At that instant $\vec{a} = 6\,\vec{i}$, which is indeed parallel to $\vec{i}$.
\end{exoresolu}

\begin{popfigure}
\begin{center}
\begin{tikzpicture}
\begin{axis}[
  width=0.72\linewidth, height=6cm,
  axis lines=middle,
  xlabel={$x$ (m)}, ylabel={$y$ (m)},
  xmin=-0.5, xmax=4.5, ymin=3, ymax=8.5,
  xtick={0,1,2,3,4}, ytick={4,5,6,7,8},
  samples=60, domain=0:1.5,
  every axis x label/.style={at={(ticklabel* cs:1.02)}, anchor=west},
  every axis y label/.style={at={(ticklabel* cs:1.02)}, anchor=south},
]
\addplot[popblue, thick] ({3*x^2-2*x},{x^3+4});
\addplot[poporange, ->, thick] coordinates {(2.5,5.125) (4.5,6.875)};
\node[popblue, anchor=south west] at (axis cs:1.6,6.4) {path of the particle};
\node[poporange, anchor=west] at (axis cs:3.6,6.3) {$\vec{v}$ at $t=1$};
\end{axis}
\end{tikzpicture}
\end{center}
\legende{Trajectory of $\vec{r} = (3t^2 - 2t)\vec{i} + (t^3 + 4)\vec{j}$ for $0 \leq t \leq 1.5$. The velocity vector is tangent to the path: at $t = 1$, $\vec{v} = 4\vec{i} + 3\vec{j}$.}
\end{popfigure}

\courssub{Finding when a particle moves in a given direction}

\begin{methode}[Direction of motion conditions]
The \cle{direction of motion} of a particle is the direction of its velocity vector $\vec{v}$.
\begin{enumerate}
\item Compute $\vec{v} = \dot{\vec{r}}$.
\item \emph{Moving parallel to $\vec{i}$} (or to the $x$-axis): set the $\vec{j}$-component of $\vec{v}$ to zero, $\vec{i}$-component non-zero.
\item \emph{Moving parallel to a line $y = mx$}: the components satisfy $\dot{y} = m\dot{x}$.
\item \emph{Moving at $45°$ to the axes}: $|\dot{x}| = |\dot{y}|$.
\item \emph{Instantaneously at rest}: both components of $\vec{v}$ are zero simultaneously.
\item Solve the resulting equation(s) for $t$, and check the domain $t \geq 0$.
\end{enumerate}
\textbf{Reflex:} direction of motion concerns $\vec{v}$, never $\vec{r}$ or $\vec{a}$.
\end{methode}

\begin{exoresolu}[Direction of motion]
A particle has position vector $\vec{r} = (t^2 - 4t)\,\vec{i} + (2t^2 - 6t)\,\vec{j}$ metres at time $t$ s.
(a) Determine whether the particle is ever instantaneously at rest.\\
(b) Find the time at which the particle moves parallel to $\vec{i} + \vec{j}$.
\tcblower
\textbf{Solution.}\\
Velocity: $\vec{v} = \dot{\vec{r}} = (2t - 4)\,\vec{i} + (4t - 6)\,\vec{j}$.\\[4pt]
(a) At rest requires $2t - 4 = 0$ \emph{and} $4t - 6 = 0$ simultaneously, i.e. $t = 2$ and $t = 1.5$: impossible. Hence the particle is \textbf{never} instantaneously at rest. (For example, at $t = 0$, $\vec{v} = -4\vec{i} - 6\vec{j} \neq \vec{0}$.)\\[4pt]
(b) Moving parallel to $\vec{i} + \vec{j}$ means the components of $\vec{v}$ are equal:
\[ 2t - 4 = 4t - 6 \quad\Longrightarrow\quad 2 = 2t \quad\Longrightarrow\quad t = 1. \]
Check: at $t = 1$, $\vec{v} = -2\,\vec{i} - 2\,\vec{j} = -2(\vec{i} + \vec{j})$, which is indeed parallel to $\vec{i} + \vec{j}$ (in the opposite sense). Hence $t = 1$ s.
\end{exoresolu}

\courssub{Integrating acceleration or velocity vectors}

\begin{propriete}[Componentwise integration of vectors]
Integration reverses differentiation, component by component:
\[ \vec{v}(t) = \int \vec{a}(t)\, dt, \qquad \vec{r}(t) = \int \vec{v}(t)\, dt, \]
each integration producing a \cle{constant vector of integration} which is fixed by an initial condition. (Not examinable as a proof; usage is required.)
\end{propriete}

\begin{methode}[Integrating vectors with initial conditions]
To recover velocity from acceleration, or displacement from velocity:
\begin{enumerate}
\item Integrate \cle{each component separately} with respect to $t$:
\[ \vec{v} = \int \vec{a}\, dt, \qquad \vec{r} = \int \vec{v}\, dt. \]
\item Add a \cle{constant vector of integration} $\vec{c} = c_1\vec{i} + c_2\vec{j}$ (one constant per component).
\item Use the given initial condition (e.g. $\vec{v} = \vec{u}$ when $t = 0$) to determine $\vec{c}$.
\item If a second integration is needed, repeat: integrate, add a new constant vector, use the initial position.
\item For a \emph{definite} integral, the constant is not needed: $\displaystyle \vec{v}(t) - \vec{v}(t_0) = \int_{t_0}^{t} \vec{a}\, ds$.
\end{enumerate}
\textbf{Reflexes:} the constant of integration is a \emph{vector}; always check your answer by differentiating back.
\end{methode}

\begin{exoresolu}[From acceleration to displacement]
A particle starts from the point with position vector $2\,\vec{i} - \vec{j}$ m with velocity $3\,\vec{i}\ \mathrm{m\,s^{-1}}$. Its acceleration at time $t$ s is
\[ \vec{a} = 2\,\vec{i} + 6t\,\vec{j} \ \mathrm{m\,s^{-2}}. \]
Find its velocity and position vector at time $t$, and hence its distance from the origin when $t = 2$.
\tcblower
\textbf{Solution.}\\
\textbf{Step 1 — Velocity.} Integrate componentwise:
\[ \vec{v} = \int (2\,\vec{i} + 6t\,\vec{j})\, dt = (2t + c_1)\,\vec{i} + (3t^2 + c_2)\,\vec{j}. \]
Initial condition: $\vec{v}(0) = 3\,\vec{i}$, so $c_1 = 3$, $c_2 = 0$:
\[ \vec{v} = (2t + 3)\,\vec{i} + 3t^2\,\vec{j} \ \mathrm{m\,s^{-1}}. \]
\textbf{Step 2 — Position.} Integrate again:
\[ \vec{r} = \int \vec{v}\, dt = (t^2 + 3t + d_1)\,\vec{i} + (t^3 + d_2)\,\vec{j}. \]
Initial condition: $\vec{r}(0) = 2\,\vec{i} - \vec{j}$, so $d_1 = 2$, $d_2 = -1$:
\[ \vec{r} = (t^2 + 3t + 2)\,\vec{i} + (t^3 - 1)\,\vec{j} \ \text{m}. \]
\textbf{Step 3 — At $t = 2$:} $\vec{r} = (4 + 6 + 2)\vec{i} + (8 - 1)\vec{j} = 12\,\vec{i} + 7\,\vec{j}$ m.\\
Distance from origin: $|\vec{r}| = \sqrt{144 + 49} = \sqrt{193} \approx 13.9$ m.\\[4pt]
\textbf{Check:} $\dot{\vec{r}} = (2t+3)\vec{i} + 3t^2\vec{j} = \vec{v}$ \checkmark\ and $\dot{\vec{v}} = 2\vec{i} + 6t\vec{j} = \vec{a}$ \checkmark
\end{exoresolu}

\courssub{Definite integrals of vectors: change of velocity and impulse}

\begin{methode}[Definite vector integrals in mechanics]
\begin{enumerate}
\item The \cle{change in velocity} between $t_1$ and $t_2$ is
\[ \Delta\vec{v} = \int_{t_1}^{t_2} \vec{a}\, dt. \]
\item The \cle{displacement} between $t_1$ and $t_2$ is
\[ \vec{s} = \int_{t_1}^{t_2} \vec{v}\, dt \quad (\text{note: this is } \vec{r}(t_2) - \vec{r}(t_1), \text{ not the distance travelled}). \]
\item Integrate each component between the limits in the usual way.
\item For \cle{impulse}: $\vec{I} = \int_{t_1}^{t_2} \vec{F}\, dt = m\,\vec{v}(t_2) - m\,\vec{v}(t_1)$ (impulse–momentum principle).
\end{enumerate}
\textbf{Reflex:} displacement from integration is a \emph{vector}; the scalar "distance travelled" requires integrating the \emph{speed} $|\vec{v}|$ and is generally different.
\end{methode}

\begin{exoresolu}[Displacement from a definite integral]
A particle moves with velocity $\vec{v} = (3t^2 - 4)\,\vec{i} + (2t)\,\vec{j}\ \mathrm{m\,s^{-1}}$. Find (a) the displacement of the particle between $t = 0$ and $t = 3$; (b) the change in velocity between $t = 1$ and $t = 2$, given $\vec{a} = 6t\,\vec{i} + 2\,\vec{j}$.
\tcblower
\textbf{Solution.}\\
(a) Displacement:
\begin{align*}
\vec{s} &= \int_0^3 \left[(3t^2 - 4)\,\vec{i} + 2t\,\vec{j}\right] dt \\
&= \left[ (t^3 - 4t)\,\vec{i} + t^2\,\vec{j} \right]_0^3 \\
&= \left[(27 - 12) - 0\right]\vec{i} + \left[9 - 0\right]\vec{j} \\
&= 15\,\vec{i} + 9\,\vec{j} \ \text{m}.
\end{align*}
(b) Change in velocity:
\[ \Delta\vec{v} = \int_1^2 (6t\,\vec{i} + 2\,\vec{j})\, dt = \left[3t^2\,\vec{i} + 2t\,\vec{j}\right]_1^2 = (12 - 3)\vec{i} + (4 - 2)\vec{j} = 9\,\vec{i} + 2\,\vec{j}\ \mathrm{m\,s^{-1}}. \]
\textbf{Check (b):} since $\vec{a} = \dot{\vec{v}}$ with $\vec{v} = (3t^2-4)\vec{i} + 2t\vec{j}$, we have $\vec{v}(2) - \vec{v}(1) = (8\vec{i} + 4\vec{j}) - (-\vec{i} + 2\vec{j}) = 9\vec{i} + 2\vec{j}$ \checkmark
\end{exoresolu}

\courssub{Motion with constant acceleration in vector form}

\begin{propriete}[Vector SUVAT formulas for constant acceleration]
When the acceleration $\vec{a}$ is \cle{constant}, integrating $\vec{a}$ twice and applying the initial conditions $\vec{v}(0) = \vec{u}$, $\vec{r}(0) = \vec{r}_0$ gives:
\[ \vec{v} = \vec{u} + \vec{a}\,t, \qquad \vec{r} = \vec{r}_0 + \vec{u}\,t + \frac{1}{2}\vec{a}\,t^2. \]
These are the vector analogues of the one-dimensional SUVAT equations; each component separately satisfies the scalar SUVAT relations. (The derivation by direct integration is examinable.)
\end{propriete}

\begin{methode}[Constant-acceleration vector formulas]
When $\vec{a}$ is \cle{constant}:
\begin{enumerate}
\item Use $\vec{v} = \vec{u} + \vec{a}\,t$ and $\vec{r} = \vec{r}_0 + \vec{u}\,t + \dfrac{1}{2}\vec{a}\,t^2$.
\item Identify $\vec{u}$ (initial velocity), $\vec{r}_0$ (initial position) and $\vec{a}$ from the question.
\item To find when/where a particle crosses an axis, set the relevant component of $\vec{r}$ to zero and solve for $t$; then substitute back.
\item To find when two particles \cle{collide}, set $\vec{r}_1(t) = \vec{r}_2(t)$: \emph{both} components must agree for the \emph{same} $t$.
\end{enumerate}
\end{methode}

\begin{exoresolu}[Constant acceleration: crossing an axis]
A particle starts from the point $(4, 0)$ m with velocity $\vec{u} = 2\,\vec{i} + 3\,\vec{j}\ \mathrm{m\,s^{-1}}$ and moves under constant acceleration $\vec{a} = -\vec{j}\ \mathrm{m\,s^{-2}}$.
(a) Find its position vector at time $t$.\\
(b) Find the time at which it crosses the $x$-axis and its position and speed at that instant.
\tcblower
\textbf{Solution.}\\
(a) With $\vec{r}_0 = 4\,\vec{i}$:
\[ \vec{r} = \vec{r}_0 + \vec{u}t + \tfrac{1}{2}\vec{a}t^2 = 4\vec{i} + (2\vec{i} + 3\vec{j})t - \tfrac{1}{2}t^2\vec{j} = (4 + 2t)\,\vec{i} + \left(3t - \tfrac{1}{2}t^2\right)\vec{j}. \]
(b) Crossing the $x$-axis: the $\vec{j}$-component is zero:
\[ 3t - \tfrac{1}{2}t^2 = 0 \quad\Longrightarrow\quad t\left(3 - \tfrac{1}{2}t\right) = 0 \quad\Longrightarrow\quad t = 0 \ \text{or}\ t = 6. \]
$t = 0$ is the starting point (on the axis), so the particle crosses again at $t = 6$ s.\\
Position: $\vec{r}(6) = (4 + 12)\vec{i} + 0\,\vec{j} = 16\,\vec{i}$ m, i.e. the point $(16, 0)$.\\
Velocity: $\vec{v} = \vec{u} + \vec{a}t = 2\vec{i} + (3 - t)\vec{j}$, so $\vec{v}(6) = 2\vec{i} - 3\vec{j}\ \mathrm{m\,s^{-1}}$.\\
Speed: $|\vec{v}(6)| = \sqrt{4 + 9} = \sqrt{13} \approx 3.61\ \mathrm{m\,s^{-1}}$.
\end{exoresolu}

\courssub{Collision and interception problems}

\begin{methode}[Collision of two particles]
\begin{enumerate}
\item Write the position vector of each particle as a function of the \emph{same} time variable $t$ (integrate each velocity, or use constant-acceleration formulas).
\item A \cle{collision} occurs when $\vec{r}_A(t) = \vec{r}_B(t)$ for the same value of $t$.
\item Equate the $\vec{i}$-components and the $\vec{j}$-components separately, giving two equations.
\item Solve one equation for $t$, then \emph{verify} that the same $t$ satisfies the other equation. If no common $t$ exists, the particles do not collide (their paths may cross at different times).
\item For an \emph{interception} problem (e.g. a patrol boat intercepting a pirogue on the Wouri), choose the unknown velocity so that the positions coincide, and solve for the unknown components.
\end{enumerate}
\end{methode}

\begin{exoresolu}[Do the particles collide?]
Particle $A$ starts from $2\,\vec{i} + \vec{j}$ m with constant velocity $\vec{i} + 2\,\vec{j}\ \mathrm{m\,s^{-1}}$. Particle $B$ starts from $8\,\vec{i} + 3\,\vec{j}$ m with constant velocity $-\vec{i} + \vec{j}\ \mathrm{m\,s^{-1}}$. Determine whether the particles collide; if so, find the time and position of collision.
\tcblower
\textbf{Solution.}\\
Position vectors at time $t$ (since velocities are constant, $\vec{r} = \vec{r}_0 + \vec{v}t$):
\begin{align*}
\vec{r}_A &= (2 + t)\,\vec{i} + (1 + 2t)\,\vec{j},\\
\vec{r}_B &= (8 - t)\,\vec{i} + (3 + t)\,\vec{j}.
\end{align*}
Collision requires equality of \emph{both} components at the \emph{same} time:
\begin{align*}
\vec{i}\text{-components:} \quad & 2 + t = 8 - t \quad\Longrightarrow\quad 2t = 6 \quad\Longrightarrow\quad t = 3,\\
\vec{j}\text{-components:} \quad & 1 + 2t = 3 + t \quad\Longrightarrow\quad t = 2.
\end{align*}
The two equations give \emph{different} times ($t = 3$ and $t = 2$), so there is no single instant at which both coordinates coincide.\\[4pt]
\textbf{Conclusion:} the particles do \textbf{not} collide. Indeed, at $t = 3$ the $\vec{i}$-components agree (both equal $5$) but the $\vec{j}$-components are $7$ and $6$; at $t = 2$ the $\vec{j}$-components agree (both equal $5$) but the $\vec{i}$-components are $4$ and $6$. The two straight-line paths never pass through the same point.
\end{exoresolu}

\courssub{Mixed problems: force, Newton's second law and vector calculus}

\begin{methode}[From force to motion via $\vec{F} = m\vec{a}$]
\begin{enumerate}
\item Use Newton's second law in vector form: $\vec{a} = \dfrac{\vec{F}}{m}$.
\item If $\vec{F}$ depends on $t$, integrate $\vec{a}(t)$ to get $\vec{v}$, using the initial velocity; integrate again for $\vec{r}$.
\item If $\vec{F}$ is constant, use the constant-acceleration formulas.
\item The \cle{momentum} is $\vec{p} = m\vec{v}$; the \cle{impulse} over $[t_1, t_2]$ is $\vec{I} = \int_{t_1}^{t_2} \vec{F}\, dt = \Delta\vec{p}$.
\item Kinetic energy: $E_k = \dfrac{1}{2}m|\vec{v}|^2$ — remember to use the \emph{speed} (magnitude), not the vector.
\end{enumerate}
\textbf{Reflex:} mass is a scalar — it never changes the \emph{direction} of $\vec{F}$ when computing $\vec{a}$.
\end{methode}

\begin{exoresolu}[Variable force on a particle]
A particle of mass $0.5$ kg moves under the force $\vec{F} = 4t\,\vec{i} + 3\,\vec{j}$ N, where $t$ is the time in seconds. At $t = 0$ the particle is at the origin moving with velocity $2\,\vec{i}\ \mathrm{m\,s^{-1}}$.\\
(a) Find the velocity and position vectors at time $t$.\\
(b) Find the momentum and kinetic energy of the particle when $t = 2$.\\
(c) Find the impulse of the force during the first 2 seconds, and verify the impulse–momentum principle.
\tcblower
\textbf{Solution.}\\
(a) Newton's second law: $\vec{a} = \dfrac{\vec{F}}{m} = \dfrac{4t\,\vec{i} + 3\,\vec{j}}{0.5} = 8t\,\vec{i} + 6\,\vec{j}\ \mathrm{m\,s^{-2}}$.\\[3pt]
Integrate: $\vec{v} = (4t^2 + c_1)\vec{i} + (6t + c_2)\vec{j}$. With $\vec{v}(0) = 2\vec{i}$: $c_1 = 2$, $c_2 = 0$:
\[ \vec{v} = (4t^2 + 2)\,\vec{i} + 6t\,\vec{j}\ \mathrm{m\,s^{-1}}. \]
Integrate again: $\vec{r} = \left(\tfrac{4}{3}t^3 + 2t + d_1\right)\vec{i} + (3t^2 + d_2)\vec{j}$. With $\vec{r}(0) = \vec{0}$: $d_1 = d_2 = 0$:
\[ \vec{r} = \left(\tfrac{4}{3}t^3 + 2t\right)\vec{i} + 3t^2\,\vec{j}\ \text{m}. \]
(b) At $t = 2$: $\vec{v} = (16 + 2)\vec{i} + 12\vec{j} = 18\vec{i} + 12\vec{j}\ \mathrm{m\,s^{-1}}$.\\
Momentum: $\vec{p} = m\vec{v} = 0.5(18\vec{i} + 12\vec{j}) = 9\vec{i} + 6\vec{j}\ \mathrm{kg\,m\,s^{-1}}$.\\
Speed: $|\vec{v}| = \sqrt{324 + 144} = \sqrt{468} = 6\sqrt{13}\ \mathrm{m\,s^{-1}}$.\\
Kinetic energy: $E_k = \tfrac{1}{2}(0.5)(468) = 117$ J.\\[4pt]
(c) Impulse:
\[ \vec{I} = \int_0^2 (4t\,\vec{i} + 3\,\vec{j})\, dt = \left[2t^2\,\vec{i} + 3t\,\vec{j}\right]_0^2 = 8\,\vec{i} + 6\,\vec{j}\ \mathrm{N\,s}. \]
Verification: change of momentum $= \vec{p}(2) - \vec{p}(0) = (9\vec{i} + 6\vec{j}) - 0.5(2\vec{i}) = 9\vec{i} + 6\vec{j} - \vec{i} = 8\vec{i} + 6\vec{j}$ N\,s $= \vec{I}$ \checkmark
\end{exoresolu}

\courssub{Exam-style synthesis problem}

\begin{methode}[Strategy for a full GCE-style vector kinematics question]
\begin{enumerate}
\item \textbf{Read and identify:} what is given ($\vec{r}$, $\vec{v}$ or $\vec{a}$? initial conditions?) and what is required.
\item \textbf{Differentiate or integrate} componentwise as needed, adding constant vectors when integrating.
\item \textbf{Apply conditions} (at rest, parallel to a vector, crossing an axis, collision) by equating components.
\item \textbf{Compute magnitudes and directions} only at the end, with correct units.
\item \textbf{Check:} differentiate your final $\vec{r}$ twice to recover $\vec{a}$; verify initial conditions.
\end{enumerate}
\end{methode}

\begin{exoresolu}[GCE synthesis: full kinematics question]
A particle $P$ of mass $2$ kg moves in a horizontal plane. At time $t$ s its acceleration is
\[ \vec{a} = (6t - 4)\,\vec{i} + 2\,\vec{j}\ \mathrm{m\,s^{-2}}. \]
When $t = 0$, $P$ is at the point with position vector $3\,\vec{j}$ m, moving with velocity $4\,\vec{i}\ \mathrm{m\,s^{-1}}$.\\
(a) Find expressions for $\vec{v}$ and $\vec{r}$ at time $t$.\\
(b) Show that the $\vec{i}$-component of the velocity of $P$ is never zero, and find its minimum value and the position of $P$ at that instant.\\
(c) Find the force acting on $P$ when $t = 2$, and its magnitude.\\
(d) Find the speed of $P$ when $t = 2$ and the angle its direction of motion makes with $\vec{i}$.
\tcblower
\textbf{Solution.}\\
(a) \textbf{Velocity:} integrate $\vec{a}$:
\[ \vec{v} = \int \left[(6t-4)\vec{i} + 2\vec{j}\right] dt = (3t^2 - 4t + c_1)\vec{i} + (2t + c_2)\vec{j}. \]
$\vec{v}(0) = 4\vec{i}$ gives $c_1 = 4$, $c_2 = 0$:
\[ \vec{v} = (3t^2 - 4t + 4)\,\vec{i} + 2t\,\vec{j}\ \mathrm{m\,s^{-1}}. \]
\textbf{Position:} integrate $\vec{v}$:
\[ \vec{r} = (t^3 - 2t^2 + 4t + d_1)\vec{i} + (t^2 + d_2)\vec{j}. \]
$\vec{r}(0) = 3\vec{j}$ gives $d_1 = 0$, $d_2 = 3$:
\[ \vec{r} = (t^3 - 2t^2 + 4t)\,\vec{i} + (t^2 + 3)\,\vec{j}\ \text{m}. \]
(b) The $\vec{i}$-component of $\vec{v}$ is $3t^2 - 4t + 4$. Completing the square:
\[ 3t^2 - 4t + 4 = 3\left(t - \tfrac{2}{3}\right)^2 + \tfrac{8}{3} \geq \tfrac{8}{3} > 0 \quad \text{for all } t, \]
so it is never zero (equivalently, the discriminant is $(-4)^2 - 4(3)(4) = -32 < 0$). Its minimum value is $\tfrac{8}{3}\ \mathrm{m\,s^{-1}}$, attained at $t = \tfrac{2}{3}$ s. Position then:
\[ \vec{r}\left(\tfrac{2}{3}\right) = \left(\tfrac{8}{27} - \tfrac{8}{9} + \tfrac{8}{3}\right)\vec{i} + \left(\tfrac{4}{9} + 3\right)\vec{j} = \tfrac{56}{27}\,\vec{i} + \tfrac{31}{9}\,\vec{j}\ \text{m}. \]
(c) Force: $\vec{F} = m\vec{a} = 2\left[(6t-4)\vec{i} + 2\vec{j}\right]$. At $t = 2$: $\vec{F} = 2(8\vec{i} + 2\vec{j}) = 16\vec{i} + 4\vec{j}$ N.\\
Magnitude: $|\vec{F}| = \sqrt{256 + 16} = \sqrt{272} = 4\sqrt{17} \approx 16.5$ N.\\[4pt]
(d) At $t = 2$: $\vec{v} = (12 - 8 + 4)\vec{i} + 4\vec{j} = 8\vec{i} + 4\vec{j}\ \mathrm{m\,s^{-1}}$.\\
Speed: $|\vec{v}| = \sqrt{64 + 16} = \sqrt{80} = 4\sqrt{5} \approx 8.94\ \mathrm{m\,s^{-1}}$.\\
Direction: $\tan\theta = \dfrac{4}{8} = 0.5$, so $\theta = \arctan(0.5) \approx 26.6°$ above the direction of $\vec{i}$.
\end{exoresolu}

\begin{aretenir}[The kinematic triangle of vector calculus]
\begin{itemize}
\item $\vec{v} = \dot{\vec{r}}$ and $\vec{a} = \dot{\vec{v}} = \ddot{\vec{r}}$: differentiate \emph{component by component}; $\vec{i}, \vec{j}$ are constants.
\item $\vec{v} = \int \vec{a}\, dt$ and $\vec{r} = \int \vec{v}\, dt$: add a \emph{vector} constant and fix it with an initial condition.
\item Displacement between two times: $\vec{s} = \int_{t_1}^{t_2} \vec{v}\, dt$; change of velocity: $\Delta\vec{v} = \int_{t_1}^{t_2} \vec{a}\, dt$.
\item Constant acceleration: $\vec{v} = \vec{u} + \vec{a}t$, $\vec{r} = \vec{r}_0 + \vec{u}t + \tfrac{1}{2}\vec{a}t^2$.
\item Direction of motion = direction of $\vec{v}$; speed $= |\vec{v}|$; collision $\Leftrightarrow \vec{r}_A(t) = \vec{r}_B(t)$ for the same $t$.
\item Newton's second law: $\vec{F} = m\vec{a}$; impulse $\vec{I} = \int \vec{F}\, dt = \Delta(m\vec{v})$.
\end{itemize}
\end{aretenir}

\begin{attention}[Frequent errors in the GCE examination]
\begin{itemize}
\item Differentiating or integrating only one component and forgetting the other.
\item Forgetting the \emph{vector} constant of integration, or replacing it by a scalar.
\item Using $\vec{r}$ instead of $\vec{v}$ for "direction of motion" questions.
\item Confusing displacement $\int \vec{v}\,dt$ (a vector) with distance travelled (a scalar).
\item Declaring a collision when only \emph{one} pair of components agrees — both must agree at the \emph{same} time.
\item Giving a speed with a negative sign, or forgetting units ($\mathrm{m\,s^{-1}}$, $\mathrm{m\,s^{-2}}$, N).
\end{itemize}
\end{attention}

\begin{exercice}[Practice: vector kinematics]
\begin{enumerate}
\item A particle has position vector $\vec{r} = (t^3 - 3t)\,\vec{i} + (2t^2)\,\vec{j}$ m at time $t$ s. Find (a) its velocity and acceleration at time $t$; (b) its speed when $t = 1$; (c) the time at which it moves parallel to $\vec{j}$.
\item A particle starts from rest at the point $\vec{i} + 2\vec{j}$ m and moves with acceleration $\vec{a} = 6t\,\vec{i} - 2\,\vec{j}\ \mathrm{m\,s^{-2}}$. Find its velocity and position vector at time $t$, and its displacement between $t = 0$ and $t = 2$.
\item Particle $A$ starts from $3\vec{j}$ m with constant velocity $2\vec{i} + \vec{j}\ \mathrm{m\,s^{-1}}$. Particle $B$ starts from $6\vec{i}$ m with constant velocity $\vec{i} + 2\vec{j}\ \mathrm{m\,s^{-1}}$. Show that $A$ and $B$ collide, and find the time and position of collision.
\item A particle of mass $3$ kg moves under the force $\vec{F} = 6\,\vec{i} + 12t\,\vec{j}$ N. Given that $\vec{v}(0) = -\vec{j}\ \mathrm{m\,s^{-1}}$, find $\vec{v}(t)$ and the kinetic energy of the particle when $t = 1$.
\end{enumerate}
\end{exercice}

\begin{corrige}[Exercise 1]
\textbf{1.} (a) $\vec{v} = (3t^2 - 3)\vec{i} + 4t\,\vec{j}\ \mathrm{m\,s^{-1}}$; $\vec{a} = 6t\,\vec{i} + 4\,\vec{j}\ \mathrm{m\,s^{-2}}$.\\
(b) At $t = 1$: $\vec{v} = 0\,\vec{i} + 4\,\vec{j}$, speed $= 4\ \mathrm{m\,s^{-1}}$.\\
(c) Parallel to $\vec{j}$: $\vec{i}$-component of $\vec{v}$ zero: $3t^2 - 3 = 0 \Rightarrow t = 1$ (since $t \geq 0$). Check: $\vec{v}(1) = 4\vec{j}$, indeed parallel to $\vec{j}$.\\[4pt]
\textbf{2.} $\vec{v} = \int (6t\vec{i} - 2\vec{j})\,dt = (3t^2 + c_1)\vec{i} + (-2t + c_2)\vec{j}$; rest at $t = 0$ gives $c_1 = c_2 = 0$: $\vec{v} = 3t^2\,\vec{i} - 2t\,\vec{j}\ \mathrm{m\,s^{-1}}$.\\
$\vec{r} = (t^3 + d_1)\vec{i} + (-t^2 + d_2)\vec{j}$; $\vec{r}(0) = \vec{i} + 2\vec{j}$ gives $d_1 = 1$, $d_2 = 2$: $\vec{r} = (t^3 + 1)\vec{i} + (2 - t^2)\vec{j}$ m.\\
Displacement $= \vec{r}(2) - \vec{r}(0) = (9\vec{i} - 2\vec{j}) - (\vec{i} + 2\vec{j}) = 8\vec{i} - 4\vec{j}$ m.\\[4pt]
\textbf{3.} $\vec{r}_A = 2t\,\vec{i} + (3 + t)\vec{j}$; $\vec{r}_B = (6 + t)\vec{i} + 2t\,\vec{j}$. Equate components: $2t = 6 + t \Rightarrow t = 6$; $3 + t = 2t \Rightarrow t = 3$. These differ, so... check again: $\vec{i}$: $2t = 6 + t \Rightarrow t = 6$; $\vec{j}$: $3 + t = 2t \Rightarrow t = 3$. No common time: the particles do \textbf{not} collide. (Note: always verify both equations before concluding.)\\[4pt]
\textbf{4.} $\vec{a} = \vec{F}/m = 2\vec{i} + 4t\,\vec{j}\ \mathrm{m\,s^{-2}}$. $\vec{v} = (2t + c_1)\vec{i} + (2t^2 + c_2)\vec{j}$; $\vec{v}(0) = -\vec{j}$ gives $c_1 = 0$, $c_2 = -1$: $\vec{v} = 2t\,\vec{i} + (2t^2 - 1)\vec{j}\ \mathrm{m\,s^{-1}}$.\\
At $t = 1$: $\vec{v} = 2\vec{i} + \vec{j}$, $|\vec{v}|^2 = 5$, so $E_k = \tfrac{1}{2}(3)(5) = 7.5$ J.
\end{corrige}

\newpage

\coursec{Exercises}

\courssub{I check my knowledge}

\begin{exercice}[Multiple choice]
A particle has position vector $\vec{r}(t) = (t^2 + 1)\vec{i} + 3t\vec{j}$ metres. Its velocity vector at time $t$ is:
\begin{enumerate}
\item[(a)] $(2t)\vec{i} + 3\vec{j}$
\item[(b)] $(t^2 + 1)\vec{i} + 3\vec{j}$
\item[(c)] $2\vec{i}$
\item[(d)] $\left(\tfrac{t^3}{3} + t\right)\vec{i} + \tfrac{3t^2}{2}\vec{j}$
\end{enumerate}
Choose the correct answer and justify briefly.
\end{exercice}

\begin{exercice}[True or false]
For each statement, say whether it is \textbf{true} or \textbf{false} and justify your answer.
\begin{enumerate}
\item If $\vec{r}(t) = x(t)\vec{i} + y(t)\vec{j}$, then $\dot{\vec{r}}(t) = \dot{x}(t)\vec{i} + \dot{y}(t)\vec{j}$, since $\vec{i}$ and $\vec{j}$ are constant vectors.
\item The speed of a particle is the magnitude of its acceleration vector.
\item If $\displaystyle \int_0^1 \vec{v}(t)\,dt = \vec{0}$, then the particle did not move between $t = 0$ and $t = 1$.
\item When integrating a vector function, a constant \emph{vector} of integration must be added.
\end{enumerate}
\end{exercice}

\begin{exercice}[Definitions]
\begin{enumerate}
\item Define the \cle{velocity vector} of a particle whose position vector is $\vec{r}(t)$.
\item Define the \cle{acceleration vector} of the particle.
\item State the relationship between the displacement vector of a particle between times $t_1$ and $t_2$ and its velocity vector $\vec{v}(t)$.
\end{enumerate}
\end{exercice}

\begin{exercice}[Direct calculation]
Given $\vec{r}(t) = (3t^2 - 2t)\vec{i} + (t^3 + 1)\vec{j}$ metres, find:
\begin{enumerate}
\item the velocity vector $\vec{v}(t)$ and the acceleration vector $\vec{a}(t)$;
\item the velocity and acceleration at $t = 2\ \mathrm{s}$;
\item the speed and the direction of motion (angle with $\vec{i}$) at $t = 2\ \mathrm{s}$.
\end{enumerate}
\end{exercice}

\courssub{I practise}

\begin{exercice}[Constant speed, central acceleration]
A particle moves with position vector $\vec{r}(t) = 4\cos(2t)\vec{i} + 4\sin(2t)\vec{j}$ metres.
\begin{enumerate}
\item Find $\vec{v}(t)$ and show that the speed is constant.
\item Find $\vec{a}(t)$ and show that the acceleration is always directed towards the origin.
\item Describe the path of the particle.
\end{enumerate}
\end{exercice}

\begin{exercice}[From acceleration to position]
The acceleration of a particle is $\vec{a} = (6t - 4)\vec{i} + 2\vec{j}\ \mathrm{m\,s^{-2}}$. Initially the particle is at the origin, moving with velocity $5\vec{i}\ \mathrm{m\,s^{-1}}$.
\begin{enumerate}
\item Find its velocity vector $\vec{v}(t)$.
\item Find its position vector $\vec{r}(t)$.
\item Find its distance from the origin when $t = 3\ \mathrm{s}$.
\end{enumerate}
\end{exercice}

\begin{exercice}[Displacement versus distance]
A particle has velocity $\vec{v} = (2t - 4)\vec{i} + 3\vec{j}\ \mathrm{m\,s^{-1}}$.
\begin{enumerate}
\item Find the displacement vector between $t = 0$ and $t = 4\ \mathrm{s}$.
\item Explain why the distance travelled is greater than the magnitude of the displacement, and compute the actual distance travelled.
\end{enumerate}
\end{exercice}

\begin{exercice}[Collision of two particles]
Two particles $A$ and $B$ move in the same plane with position vectors
\[ \vec{r}_A = (t^2 + 1)\vec{i} + 2t\vec{j}, \qquad \vec{r}_B = (3t - 1)\vec{i} + (t + 2)\vec{j} \quad \text{(metres)}. \]
\begin{enumerate}
\item Show that the particles collide, and find the time of collision.
\item Find the position vector of the point of collision.
\item Find the velocity of each particle just before the collision.
\end{enumerate}
\end{exercice}

\courssub{I prepare for the GCE}

\begin{exercice}[GCE-style: kinematics from a position vector]
The position vector of a particle moving in a plane is
\[ \vec{r} = (t^3 - 6t)\vec{i} + (t^2 + 2)\vec{j} \quad \text{metres}, \]
where $t$ is the time in seconds.
\begin{enumerate}
\item[(a)] Find the velocity vector $\vec{v}$ and the acceleration vector $\vec{a}$ of the particle. \hfill (4 marks)
\item[(b)] Find the time at which the particle is moving parallel to the vector $\vec{j}$. \hfill (3 marks)
\item[(c)] Find the angle between the acceleration vector and $\vec{i}$ at $t = 1\ \mathrm{s}$. \hfill (4 marks)
\item[(d)] Find the speed of the particle at $t = 2\ \mathrm{s}$, giving your answer to 3 significant figures. \hfill (2 marks)
\end{enumerate}
\end{exercice}

\begin{exercice}[GCE-style: integration with initial conditions]
Given the velocity $\vec{v}(t) = e^{t}\vec{i} + (1 - t^2)\vec{j}\ \mathrm{m\,s^{-1}}$ of a particle, with $\vec{r}(0) = 2\vec{i} - \vec{j}$ metres:
\begin{enumerate}
\item[(a)] Find the position vector $\vec{r}(t)$ of the particle. \hfill (5 marks)
\item[(b)] Determine the time at which the $y$-component of the velocity vanishes, and find the position of the particle at that instant. \hfill (4 marks)
\item[(c)] Find the acceleration vector at the instant found in (b). \hfill (3 marks)
\end{enumerate}
\end{exercice}

\begin{exercice}[GCE-style: a struck ball — link with projectile motion]
A ball is struck from the origin $O$ so that its acceleration is $\vec{a} = -10\vec{j}\ \mathrm{m\,s^{-2}}$ and its initial velocity is $(8\vec{i} + 6\vec{j})\ \mathrm{m\,s^{-1}}$.
\begin{enumerate}
\item[(a)] By integration, find the velocity vector $\vec{v}(t)$ and the position vector $\vec{r}(t)$ of the ball at time $t$. \hfill (5 marks)
\item[(b)] Find the greatest height reached by the ball above $O$. \hfill (4 marks)
\item[(c)] Find the time of flight, i.e.\ the time at which the ball returns to the horizontal level of $O$, and the horizontal distance travelled (the range). \hfill (4 marks)
\item[(d)] Show that the path of the ball is a parabola by eliminating $t$ between the components of $\vec{r}(t)$. \hfill (3 marks)
\end{enumerate}
\emph{(This exercise connects vector integration with projectile motion — a frequent source of GCE Advanced Level questions.)}
\end{exercice}

\newpage

\coursec{Solutions to the exercises}

\begin{corrige}[Exercise 1 — Multiple choice]
The velocity vector is obtained by differentiating each component of the position vector with respect to time:
\[ \vec{v}(t) = \dot{\vec{r}}(t) = \frac{d}{dt}(t^2 + 1)\,\vec{i} + \frac{d}{dt}(3t)\,\vec{j} = 2t\,\vec{i} + 3\,\vec{j}. \]
The correct answer is \textbf{(a)}. Option (b) is the position vector itself, (c) is the acceleration, and (d) is the result of \emph{integrating} instead of differentiating.
\end{corrige}

\begin{corrige}[Exercise 2 — True or false]
\begin{enumerate}
\item \textbf{True.} Since $\vec{i}$ and $\vec{j}$ are constant (fixed direction and magnitude), they pass through the differentiation: $\dot{\vec{r}} = \dot{x}\,\vec{i} + \dot{y}\,\vec{j}$. This is exactly the rule for differentiating a vector component by component.
\item \textbf{False.} The speed is the magnitude of the \emph{velocity} vector: $v = |\vec{v}| = |\dot{\vec{r}}|$. The magnitude of the acceleration measures the rate of change of the velocity vector, which is a different quantity.
\item \textbf{False.} $\displaystyle\int_0^1 \vec{v}\,dt = \vec{r}(1) - \vec{r}(0)$ is the \emph{displacement}. A zero displacement only means the particle finished where it started; it may have travelled a long path in between (e.g.\ one full revolution of a circle).
\item \textbf{True.} Integrating component by component introduces one arbitrary constant per component, i.e.\ a constant \emph{vector} $\vec{C}$. It is then fixed by an initial condition such as $\vec{r}(0)$.
\end{enumerate}
\end{corrige}

\begin{corrige}[Exercise 3 — Definitions]
\begin{enumerate}
\item The \cle{velocity vector} is the derivative of the position vector with respect to time:
\[ \vec{v}(t) = \dot{\vec{r}}(t) = \frac{d\vec{r}}{dt}. \]
Its direction is tangent to the path and its magnitude is the speed.
\item The \cle{acceleration vector} is the derivative of the velocity vector (second derivative of the position vector):
\[ \vec{a}(t) = \dot{\vec{v}}(t) = \ddot{\vec{r}}(t) = \frac{d^2\vec{r}}{dt^2}. \]
\item The displacement between times $t_1$ and $t_2$ is the definite integral of the velocity:
\[ \vec{r}(t_2) - \vec{r}(t_1) = \int_{t_1}^{t_2} \vec{v}(t)\,dt. \]
This is the vector form of the fundamental theorem of calculus.
\end{enumerate}
\end{corrige}

\begin{corrige}[Exercise 4 — Direct calculation]
\begin{enumerate}
\item Differentiating component by component:
\[ \vec{v}(t) = \dot{\vec{r}} = (6t - 2)\vec{i} + 3t^2\vec{j} \quad (\mathrm{m\,s^{-1}}), \]
\[ \vec{a}(t) = \dot{\vec{v}} = 6\vec{i} + 6t\vec{j} \quad (\mathrm{m\,s^{-2}}). \]
\item At $t = 2$:
\[ \vec{v}(2) = (12 - 2)\vec{i} + 12\vec{j} = 10\vec{i} + 12\vec{j}\ \mathrm{m\,s^{-1}}, \qquad \vec{a}(2) = 6\vec{i} + 12\vec{j}\ \mathrm{m\,s^{-2}}. \]
\item Speed at $t = 2$:
\[ v = |\vec{v}(2)| = \sqrt{10^2 + 12^2} = \sqrt{244} = 2\sqrt{61} \approx 15.6\ \mathrm{m\,s^{-1}}. \]
If $\theta$ is the angle between $\vec{v}(2)$ and $\vec{i}$, then
\[ \tan\theta = \frac{v_y}{v_x} = \frac{12}{10} = 1.2 \quad\Longrightarrow\quad \theta \approx 50.2^{\circ}. \]
The particle moves at about $15.6\ \mathrm{m\,s^{-1}}$ in a direction making an angle of about $50.2^{\circ}$ with $\vec{i}$.
\end{enumerate}
\end{corrige}

\begin{corrige}[Exercise 5 — Constant speed, central acceleration]
\begin{enumerate}
\item Differentiating:
\[ \vec{v}(t) = -8\sin(2t)\vec{i} + 8\cos(2t)\vec{j}\ \mathrm{m\,s^{-1}}. \]
The speed is
\[ v = |\vec{v}| = \sqrt{64\sin^2(2t) + 64\cos^2(2t)} = \sqrt{64} = 8\ \mathrm{m\,s^{-1}}, \]
which does not depend on $t$: the speed is \textbf{constant}.
\item Differentiating again:
\[ \vec{a}(t) = -16\cos(2t)\vec{i} - 16\sin(2t)\vec{j} = -4\bigl(4\cos(2t)\vec{i} + 4\sin(2t)\vec{j}\bigr) = -4\vec{r}(t). \]
Since $\vec{a}$ is a negative scalar multiple of $\vec{r}$, it is always directed from the particle \textbf{towards the origin} (centripetal acceleration).
\item From $x = 4\cos(2t)$ and $y = 4\sin(2t)$:
\[ x^2 + y^2 = 16\cos^2(2t) + 16\sin^2(2t) = 16. \]
The path is the \textbf{circle of centre $O$ and radius $4\ \mathrm{m}$}, described with constant angular speed $\omega = 2\ \mathrm{rad\,s^{-1}}$ (uniform circular motion).
\end{enumerate}
\end{corrige}

\begin{corrige}[Exercise 6 — From acceleration to position]
\begin{enumerate}
\item Integrate the acceleration:
\[ \vec{v}(t) = \int \vec{a}\,dt = \left(3t^2 - 4t\right)\vec{i} + 2t\,\vec{j} + \vec{C}. \]
Initial condition: $\vec{v}(0) = 5\vec{i}$, so $\vec{C} = 5\vec{i}$. Hence
\[ \vec{v}(t) = (3t^2 - 4t + 5)\vec{i} + 2t\,\vec{j}\ \mathrm{m\,s^{-1}}. \]
\item Integrate the velocity:
\[ \vec{r}(t) = \left(t^3 - 2t^2 + 5t\right)\vec{i} + t^2\vec{j} + \vec{D}. \]
Initial condition: $\vec{r}(0) = \vec{0}$, so $\vec{D} = \vec{0}$. Hence
\[ \vec{r}(t) = (t^3 - 2t^2 + 5t)\vec{i} + t^2\vec{j}\ \mathrm{m}. \]
\item At $t = 3$: $\vec{r}(3) = (27 - 18 + 15)\vec{i} + 9\vec{j} = 24\vec{i} + 9\vec{j}$. The distance from the origin is
\[ |\vec{r}(3)| = \sqrt{24^2 + 9^2} = \sqrt{576 + 81} = \sqrt{657} = 3\sqrt{73} \approx 25.6\ \mathrm{m}. \]
\end{enumerate}
\end{corrige}

\begin{corrige}[Exercise 7 — Displacement versus distance]
\begin{enumerate}
\item The displacement is the integral of the velocity:
\[ \vec{s} = \int_0^4 \vec{v}\,dt = \left[t^2 - 4t\right]_0^4 \vec{i} + \left[3t\right]_0^4 \vec{j} = (16 - 16)\vec{i} + 12\vec{j} = 12\vec{j}\ \mathrm{m}. \]
Its magnitude is $|\vec{s}| = 12\ \mathrm{m}$.
\item The $x$-component of the velocity, $2t - 4$, is negative for $t < 2$ and positive for $t > 2$: the particle first moves backwards along the $x$-axis, stops (in $x$) at $t = 2$, then moves forwards, returning to its initial $x$-position. The displacement only records the net change of position, whereas the distance travelled accumulates every metre covered, in both directions. Hence the distance exceeds the magnitude of the displacement.

The distance travelled is the integral of the \emph{speed}:
\[ v(t) = |\vec{v}(t)| = \sqrt{(2t-4)^2 + 9}. \]
Splitting at $t = 2$ and substituting $u = 2t - 4$ ($du = 2\,dt$):
\[ d = \int_0^4 \sqrt{(2t-4)^2 + 9}\,dt = \frac{1}{2}\int_{-4}^{4} \sqrt{u^2 + 9}\,du = \int_0^4 \sqrt{u^2 + 9}\,du. \]
Using $\displaystyle\int \sqrt{u^2 + a^2}\,du = \frac{u}{2}\sqrt{u^2 + a^2} + \frac{a^2}{2}\ln\left(u + \sqrt{u^2 + a^2}\right)$ with $a = 3$:
\[ d = \left[\frac{u}{2}\sqrt{u^2 + 9} + \frac{9}{2}\ln\left(u + \sqrt{u^2 + 9}\right)\right]_0^4 = 10 + \frac{9}{2}\ln\left(\frac{9}{3}\right) = 10 + \frac{9}{2}\ln 3 \approx 14.9\ \mathrm{m}. \]
Indeed $14.9\ \mathrm{m} > 12\ \mathrm{m} = |\vec{s}|$, as expected.
\end{enumerate}
\end{corrige}

\begin{corrige}[Exercise 8 — Collision of two particles]
\begin{enumerate}
\item The particles collide when $\vec{r}_A = \vec{r}_B$, i.e.\ when \emph{both} components are equal at the \emph{same} time:
\[ \begin{cases} t^2 + 1 = 3t - 1 & (\vec{i}\text{-components}) \\ 2t = t + 2 & (\vec{j}\text{-components}) \end{cases} \]
The second equation gives $t = 2$. Check in the first: $t^2 + 1 - 3t + 1 = t^2 - 3t + 2 = (t-1)(t-2)$, which vanishes at $t = 2$ (and at $t = 1$, but $t = 1$ does not satisfy the second equation). Both conditions hold simultaneously at $t = 2\ \mathrm{s}$: \textbf{the particles collide at $t = 2\ \mathrm{s}$}.
\item Position of collision:
\[ \vec{r}_A(2) = 5\vec{i} + 4\vec{j}\ \mathrm{m} \qquad \bigl(= \vec{r}_B(2), \text{ as required}\bigr). \]
\item Velocities: $\vec{v}_A = 2t\vec{i} + 2\vec{j}$ and $\vec{v}_B = 3\vec{i} + \vec{j}$. Just before the collision ($t = 2$):
\[ \vec{v}_A(2) = 4\vec{i} + 2\vec{j}\ \mathrm{m\,s^{-1}}, \qquad \vec{v}_B(2) = 3\vec{i} + \vec{j}\ \mathrm{m\,s^{-1}}. \]
Note that the velocities are different at the collision point — the particles meet but were not moving together.
\end{enumerate}
\end{corrige}

\begin{corrige}[Exercise 9 — GCE-style: kinematics from a position vector]
\textbf{(a)} Differentiate component by component:
\[ \vec{v} = \dot{\vec{r}} = (3t^2 - 6)\vec{i} + 2t\vec{j}\ \mathrm{m\,s^{-1}}, \qquad \vec{a} = \dot{\vec{v}} = 6t\vec{i} + 2\vec{j}\ \mathrm{m\,s^{-2}}. \]
\emph{Mark scheme: M1 for differentiating; A1 for $\vec{v}$; M1 for differentiating again; A1 for $\vec{a}$.}

\textbf{(b)} Motion parallel to $\vec{j}$ means the $\vec{i}$-component of the velocity is zero (and the $\vec{j}$-component non-zero):
\[ 3t^2 - 6 = 0 \quad\Longrightarrow\quad t^2 = 2 \quad\Longrightarrow\quad t = \sqrt{2} \approx 1.41\ \mathrm{s} \quad (t \geq 0). \]
At this time $v_y = 2\sqrt{2} \neq 0$, so the motion is indeed parallel to $\vec{j}$.
\emph{Mark scheme: M1 for setting $v_x = 0$; A1 for $t = \sqrt{2}$; B1 for rejecting the negative root or checking $v_y \neq 0$.}

\textbf{(c)} At $t = 1$: $\vec{a}(1) = 6\vec{i} + 2\vec{j}$. If $\theta$ is the angle between $\vec{a}$ and $\vec{i}$:
\[ \tan\theta = \frac{a_y}{a_x} = \frac{2}{6} = \frac{1}{3} \quad\Longrightarrow\quad \theta \approx 18.4^{\circ}. \]
\emph{Mark scheme: B1 for $\vec{a}(1)$; M1 for using $\tan\theta = a_y/a_x$ (or the scalar product); A1 for $\tan\theta = 1/3$; A1 for $\theta \approx 18.4^{\circ}$.}

\textbf{(d)} At $t = 2$: $\vec{v}(2) = (12 - 6)\vec{i} + 4\vec{j} = 6\vec{i} + 4\vec{j}$, so
\[ v = \sqrt{6^2 + 4^2} = \sqrt{52} \approx 7.21\ \mathrm{m\,s^{-1}} \quad (3\ \text{s.f.}). \]
\emph{Mark scheme: M1 for $|\vec{v}(2)|$; A1 for $7.21\ \mathrm{m\,s^{-1}}$ to 3 s.f. Examiner's expectation: candidates must give the magnitude, not the vector, and respect the required accuracy.}
\end{corrige}

\begin{corrige}[Exercise 10 — GCE-style: integration with initial conditions]
\textbf{(a)} Integrate component by component:
\[ \vec{r}(t) = \int \vec{v}\,dt = e^{t}\vec{i} + \left(t - \frac{t^3}{3}\right)\vec{j} + \vec{C}. \]
Initial condition: $\vec{r}(0) = \vec{i} + \vec{C} = 2\vec{i} - \vec{j}$, hence $\vec{C} = \vec{i} - \vec{j}$. Therefore
\[ \vec{r}(t) = (e^{t} + 1)\vec{i} + \left(t - \frac{t^3}{3} - 1\right)\vec{j}\ \mathrm{m}. \]
\emph{Mark scheme: M1 for integrating both components; A1 for $e^t$; A1 for $t - t^3/3$; M1 for using $\vec{r}(0)$; A1 for the final vector. Warning: $e^0 = 1$, not $0$ — a very common source of error.}

\textbf{(b)} The $y$-component of the velocity vanishes when
\[ 1 - t^2 = 0 \quad\Longrightarrow\quad t = 1\ \mathrm{s} \quad (t \geq 0). \]
Position at $t = 1$:
\[ \vec{r}(1) = (e + 1)\vec{i} + \left(1 - \tfrac{1}{3} - 1\right)\vec{j} = (e + 1)\vec{i} - \tfrac{1}{3}\vec{j} \approx 3.72\vec{i} - 0.333\vec{j}\ \mathrm{m}. \]
\emph{Mark scheme: M1 for solving $1 - t^2 = 0$; A1 for $t = 1$; M1 for substituting into $\vec{r}$; A1 for the position.}

\textbf{(c)} Differentiate the velocity: $\vec{a}(t) = e^{t}\vec{i} - 2t\vec{j}$. At $t = 1$:
\[ \vec{a}(1) = e\,\vec{i} - 2\vec{j} \approx 2.72\vec{i} - 2\vec{j}\ \mathrm{m\,s^{-2}}. \]
\emph{Mark scheme: M1 for differentiating $\vec{v}$; A1 for $\vec{a}(t)$; A1 for the value at $t = 1$.}
\end{corrige}

\begin{corrige}[Exercise 11 — GCE-style: a struck ball — link with projectile motion]
\textbf{(a)} Integrate the acceleration:
\[ \vec{v}(t) = \int (-10\vec{j})\,dt = -10t\,\vec{j} + \vec{C}. \]
With $\vec{v}(0) = 8\vec{i} + 6\vec{j}$: $\vec{C} = 8\vec{i} + 6\vec{j}$, so
\[ \vec{v}(t) = 8\vec{i} + (6 - 10t)\vec{j}\ \mathrm{m\,s^{-1}}. \]
Integrate again:
\[ \vec{r}(t) = 8t\,\vec{i} + \left(6t - 5t^2\right)\vec{j} + \vec{D}. \]
With $\vec{r}(0) = \vec{0}$: $\vec{D} = \vec{0}$, so
\[ \vec{r}(t) = 8t\,\vec{i} + (6t - 5t^2)\vec{j}\ \mathrm{m}. \]
\emph{Mark scheme: M1 for integrating $\vec{a}$; A1 for $\vec{v}$ with constant found from the initial condition; M1 for integrating $\vec{v}$; A1 for $\vec{r}$; B1 for using $\vec{r}(0) = \vec{0}$.}

\textbf{(b)} The greatest height is reached when the vertical component of the velocity vanishes:
\[ 6 - 10t = 0 \quad\Longrightarrow\quad t = 0.6\ \mathrm{s}. \]
Then
\[ y_{\max} = 6(0.6) - 5(0.6)^2 = 3.6 - 1.8 = 1.8\ \mathrm{m}. \]
\emph{Mark scheme: M1 for setting $v_y = 0$; A1 for $t = 0.6$; M1 for substituting into $y(t)$; A1 for $1.8\ \mathrm{m}$.}

\textbf{(c)} The ball returns to the level of $O$ when $y = 0$:
\[ 6t - 5t^2 = t(6 - 5t) = 0 \quad\Longrightarrow\quad t = 0 \ \text{(launch)} \ \text{or}\ t = 1.2\ \mathrm{s}. \]
The time of flight is $T = 1.2\ \mathrm{s}$. The range is the horizontal distance:
\[ R = x(1.2) = 8 \times 1.2 = 9.6\ \mathrm{m}. \]
\emph{Mark scheme: M1 for solving $y = 0$; A1 for $T = 1.2\ \mathrm{s}$ (rejecting $t = 0$); M1 for evaluating $x(T)$; A1 for $R = 9.6\ \mathrm{m}$.}

\textbf{(d)} From $x = 8t$ we get $t = \dfrac{x}{8}$. Substituting into $y = 6t - 5t^2$:
\[ y = 6\left(\frac{x}{8}\right) - 5\left(\frac{x}{8}\right)^2 = \frac{3x}{4} - \frac{5x^2}{64}. \]
This is of the form $y = \alpha x - \beta x^2$ with $\alpha, \beta$ constants: the equation of a \textbf{parabola} (passing through the origin, opening downwards). $\blacksquare$
\emph{Mark scheme: M1 for eliminating $t$; A1 for the Cartesian equation; B1 for the conclusion. Examiner's expectation: the elimination must be shown explicitly, and the conclusion must mention the parabolic nature of $y = \tfrac{3}{4}x - \tfrac{5}{64}x^2$. This is the standard vector route to projectile motion, which recovers the classical formulae $H = \tfrac{u_y^2}{2g}$ and $R = \tfrac{2u_x u_y}{g}$ with $g = 10\ \mathrm{m\,s^{-2}}$.}
\end{corrige}

\newpage

\coursec{Summary sheet}

\begin{popfigure}
\begin{center}
\begin{tikzpicture}[mindmap, grow cyclic, every node/.style={concept, text=white, font=\small},
  root concept/.append style={concept color=popink, font=\bfseries\small, minimum size=3.2cm},
  level 1/.append style={level distance=3.6cm, sibling angle=60},
  level 1 concept/.append style={minimum size=2.3cm, font=\scriptsize}]
\node[root concept]{Vector calculus in mechanics}
  child[concept color=popblue]{ node[align=center]{Differentiating vectors\\ $\dot{\vec r}=\dfrac{d\vec r}{dt}$ componentwise} }
  child[concept color=poporange]{ node[align=center]{Velocity \& acceleration\\ $\vec v=\dot{\vec r}$, $\vec a=\dot{\vec v}$} }
  child[concept color=popgreen]{ node[align=center]{Integrating vectors\\ $\int\vec v\,dt$ + constant $\vec C$} }
  child[concept color=poppurple]{ node[align=center]{Definite integrals\\ $\vec r(t_2)-\vec r(t_1)=\int_{t_1}^{t_2}\vec v\,dt$} }
  child[concept color=popgold]{ node[align=center]{Kinematic triangle\\ $\vec r \xrightarrow{d/dt} \vec v \xrightarrow{d/dt} \vec a$} }
  child[concept color=poppink]{ node[align=center]{Applications\\ projectiles, circular motion, GCE problems} };
\end{tikzpicture}
\end{center}
\legende{Mind map of the chapter: applications of vectors in mechanics.}
\end{popfigure}

\begin{bilan}
\textbf{Key definitions}
\begin{itemize}
\item \cle{Position vector}: $\vec r(t)=x(t)\vec i+y(t)\vec j$ gives the position of a particle at time $t$.
\item \cle{Displacement}: change of position, $\Delta\vec r=\vec r(t_2)-\vec r(t_1)$; a vector, unlike distance travelled (a scalar).
\item \cle{Velocity}: $\vec v=\dfrac{d\vec r}{dt}=\dot{x}\vec i+\dot{y}\vec j$; speed is the magnitude $|\vec v|$.
\item \cle{Acceleration}: $\vec a=\dfrac{d\vec v}{dt}=\dfrac{d^2\vec r}{dt^2}=\ddot{x}\vec i+\ddot{y}\vec j$.
\end{itemize}

\textbf{Formulas and results to know}
\begin{itemize}
\item Differentiate and integrate \textbf{component by component}: if $\vec r=x\vec i+y\vec j$ then $\dot{\vec r}=\dot{x}\vec i+\dot{y}\vec j$ and $\displaystyle\int\vec r\,dt=\Big(\int x\,dt\Big)\vec i+\Big(\int y\,dt\Big)\vec j+\vec C$.
\item Derivative of a dot product: $\dfrac{d}{dt}(\vec u\cdot\vec v)=\dot{\vec u}\cdot\vec v+\vec u\cdot\dot{\vec v}$; in particular $\dfrac{d}{dt}|\vec v|^2=2\vec v\cdot\vec a$.
\item If $|\vec r|$ is constant then $\vec r\cdot\dot{\vec r}=0$: velocity is perpendicular to position (circular motion).
\item Definite integral: $\displaystyle\int_{t_1}^{t_2}\vec v\,dt=\vec r(t_2)-\vec r(t_1)$ (displacement) and $\displaystyle\int_{t_1}^{t_2}\vec a\,dt=\vec v(t_2)-\vec v(t_1)$.
\item The kinematic triangle: $\vec r \xrightarrow{\text{differentiate}} \vec v \xrightarrow{\text{differentiate}} \vec a$; reverse the arrows with integration plus a \textbf{constant vector} fixed by initial conditions.
\end{itemize}

\textbf{Methods}
\begin{enumerate}
\item \emph{Given $\vec r(t)$}: differentiate twice for $\vec v$ and $\vec a$; compute magnitudes for speed and $|\vec a|$; set $\vec v\cdot\vec a=0$ to test perpendicularity; solve $\dot x=0$ or $\dot y=0$ for stationary directions.
\item \emph{Given $\vec a(t)$ (or $\vec v(t)$)}: integrate componentwise, then use the given initial velocity/position to find each constant vector $\vec C$.
\item \emph{Displacement over $[t_1,t_2]$}: integrate $\vec v$ between limits, or subtract position vectors.
\end{enumerate}

\textbf{Common mistakes}
\begin{itemize}
\item Forgetting the constant \textbf{vector} of integration, or treating it as a scalar.
\item Confusing displacement $\int\vec v\,dt$ with distance $\int|\vec v|\,dt$.
\item Differentiating $|\vec r|$ instead of $\vec r$: in general $\dfrac{d|\vec r|}{dt}\neq\Big|\dfrac{d\vec r}{dt}\Big|$.
\item Mixing units: keep $\mathrm{m}$, $\mathrm{s}$, $\mathrm{m\,s^{-1}}$, $\mathrm{m\,s^{-2}}$ consistent.
\end{itemize}

\textbf{GCE tips}
\begin{itemize}
\item Show each component separately; examiners award method marks per component.
\item When asked for \emph{speed}, give $|\vec v|=\sqrt{\dot x^2+\dot y^2}$, not $\vec v$.
\item Quote initial conditions explicitly when determining constants — this earns a mark.
\item Check your answer by differentiating back: $\dfrac{d}{dt}\!\left(\int\vec v\,dt\right)=\vec v$.
\end{itemize}
\end{bilan}

\findecours

\end{document}
